\documentclass[12pt]{book}
\usepackage[utf8]{inputenc}
\usepackage[T1]{fontenc}
\usepackage[danish]{babel}
\usepackage[protrusion=true,expansion=true]{microtype}
\usepackage[margin=1.4in]{geometry}
\usepackage{setspace}
\usepackage{amsmath}
\usepackage{libertinus}
\usepackage{libertinust1math}
\usepackage{textalpha}   % occasional polytonic Greek (e.g. the Homer quotation in \S\,10)
\usepackage{fancyhdr}
\usepackage[colorlinks=true,linkcolor=black,urlcolor=black,
  pdftitle={Almindelig Videnskabslære i Grundtræk},
  pdfauthor={Rasmus Nielsen}]{hyperref}
\usepackage{chngcntr}
\counterwithout{section}{chapter}
\renewcommand{\thesection}{§\,\arabic{section}}
\setlength{\headheight}{15pt}
\pagestyle{fancy}
\fancyhf{}
\fancyhead[LE,RO]{\thepage}
\fancyhead[RE]{\textit{Almindelig Videnskabslære i Grundtræk}}
\fancyhead[LO]{\textit{\leftmark}}
\setlength{\parindent}{1.5em}
\setlength{\parskip}{0pt}
\onehalfspacing

\title{Almindelig Videnskabslære i Grundtræk}
\author{R.\ Nielsen}
\date{Kjøbenhavn: Gyldendalske Boghandel (F.\ Hegel \& Søn), 1880}

\begin{document}
\maketitle
\tableofcontents

\chapter*{Forord}
\addcontentsline{toc}{chapter}{Forord}

% Skjøndt Videnskaberne ere indbyrdes forskjellige, kan den
% principielle, paa Forhold mellem Sandsning og Tænkning grundede,
% Forskjel imellem dem dog ikke være større, end at de alle,
% aprioriske saa vel som empiriske, maae lade sig henføre under
% et fælles, for almindelig Viden gjældende Grundprincip.
%
% Key claim: the question is whether the Videnskabsprincip ultimately
% rests on Sandsning (sensationism) or on the uendelige Videns Idee
% (idealism admitting teleological categories like Anlæg og Udvikling,
% Midler og Formaal). Nielsen's answer is the latter.
%
% Also signals the Darwinist context: the core controversy is whether
% teleology survives in biology (physisk-teleologisk vs. blot physisk).
% And the final paragraph stakes the *uensartede størrelser* position:
% the religious life forms itself by its own spirit, independent of
% Videnskab.
%
% Kjøbenhavn, den 23de August. R. Nielsen.

\noindent Skjøndt Videnskaberne ere indbyrdes forskjellige, kan den
principielle, paa Forhold mellem Sandsning og Tænkning grundede, Forskjel
imellem dem dog ikke være større, end at de alle, aprioriske saa vel som
empiriske, maae lade sig henføre under et fælles, for almindelig Viden
gjældende Grundprincip. Spørgsmaalet bliver da, hvorfra Vidensprincipet
væsenlig stammer; om det oprindelig beroer paa Sandsning eller tillige
paa en for Tingene selv afgjørende Tænkning; om det begrundes paa Stof og
Kraft alene, saa at dets Gyldighed i det Høieste maa indskrænke sig til
Kategorier som Aarsag og Virkning; eller det har sin Rod i den uendelige
Videns Idee, saa at Grundbestemmelser som Anlæg og Udvikling, Midler og
Formaal deri sees at have objectiv Gyldighed. Hvis det Sidste antages,
maa Antagelsen kunne retfærdiggjøres af Videnskabernes indbyrdes
Forbindelse; thi en saadan Forbindelse maa jo væsenlig være bestemt ved
den Maade, hvorpaa Principet udvikler sig. At den almindelige Ideelære
kan dyrkes for sig, Geologien for sig, Historien for sig o.s.v., lader
sig ikke nægte; men Videnskaberne maae dog i Grunden alle hænge sammen:
kan den ene undvære den anden, saa er her jo ingen Sammenhæng.

Den Undersøgelse, hvorpaa der i Videnskabslæren for Tiden særlig maa
lægges Vægt, angaaer Forholdet mellem det Physiske og det Biologiske.
Dette Forhold kan med særligt Hensyn paa Darwinshypothesen opfattes paa
to Maader: deels saaledes, at alt det Teleologiske forsvinder, og kun de
mechanisk-physiske Aarsager anerkjendes; deels saaledes, at det Physiske
underligger en naturlig teleologisk Form, og Selvudviklingen viser sig at
tilfredsstille Anlæget. Alt som Anerkjendelsen af det Darwinske Grundlag
mere og mere udbreder sig, vil Striden mellem den blot physiske og den
physisk-teleologiske Opfattelse, ikke af Darwins personlige Anskuelse,
men af Grundlagets Anvendelighed komme til Gjennembrud; thi de Rester af
Teleologi, som søge Skjul og Tilhold i de fra Darwinismen afvigende
Synsmaader, ville neppe have Haab om at kunne gjælde for sig overfor en
Theori, der ellers i det Hele som i det Enkelte prøver, optager og
tilegner sig alle Forskningens Resultater. Darwinismens Mangler, de være
nu i »Ontogenesen« eller i »Phylogenesen« eller hvor som helst, maae
kunne udfindes og afhjælpes i Conseqvens af Grundanskuelsen selv, ikke
ved udenfra kommende Tilsætninger.

Hvad Erkjendelsen af det Religiøse angaaer, vil det aandelige Liv, for
saa vidt det udvikler sig til sand Forstaaelse af sig selv, ikke undlade
at underkaste sig Videnskabens Control i Alt, hvad der særlig beroer paa
det udvortes Virkelige; men i Forhold til sit indre Væsen vil det
religiøse Liv endnu, som før, vedblive at danne sig af sin egen Aand,
efter sine egne Forbilleder, og gaae sin Gang uafhængigt af Videnskab og
videnskabelige Indsigelser.

Dette er i Korthed Charakteren af de i nærværende Skrift førte
Undersøgelser. Hvor vidt der i disse Undersøgelser, naar afsees fra
Udtryk og Vendinger, som let indfinde sig hos den, der, ifølge Sagens
Natur, umulig kan være saaledes hjemme i de forskjellige Fagvidenskabers
Termini som disses egenlige Dyrkere, maaskee skulde indeholdes Noget, som
ved nærmere Prøvelse turde vise sig at være af almindeligere Betydning,
det overlades herved sagkyndige og velvillige Dommere at afgjøre.

\bigskip
\noindent Kjøbenhavn, den 23de August.
\nopagebreak
\par\hfill R.\ Nielsen.


%%% FØRSTE DEEL: VIDEN OG VIDENSKAB %%%

\part{Viden og Videnskab}

\chapter{Videns Væsen i Almindelighed}

\section{Oprindelige Forudsætninger}
\label{sec:1}
% pp. 1--9
% Opening: ``Vor første Viden kommer, saa at sige, af sig selv.''
% The ascent from immediate knowing through doubt to the demand for
% Videnskabslære. Descartes-style radical doubt leads to bottomless
% skepticism; neither the external world nor the cogito nor the
% principle A=A provides a firm foundation.
% The resolution: the seeker must stop (``Den søgende Individualitet
% maa standse'') and choose how to break off.
% Three paths: (1) return to sund Sands (practical but not scientific);
% (2) some other resolution (left for §2 onward).
% Key structural point: Viden og Videnskab presuppose each other
% in the infinite.

\noindent Vor første Viden kommer, saa at sige, af sig selv. Vi ere
vidende om mange Ting, før vi have nogen Anelse om, hvad Viden egenlig
er. Først naar det gaaer op for os, at der er Noget, vi ikke vide, og som
vi dog gjerne gad vide, Noget, hvorom vi ere i Uvished, skjøndt vi ønske
Vished, begynde vi at lægge Eftertryk paa Ordet Viden. Og dog er der
endnu langt til videnskabelig Viden; vi maae gjøre forberedende Skridt.
Kulturhistorien lærer, at ethvert Folk har maattet gjennemgaae en Række
forskjellige Udviklingstrin og indsamle et rigt Materiale af
forskjelligartet praktisk Viden, inden det naaede saa vidt, at en egenlig
Videnskab derom kunde begynde. En Viden om Naturens Gang, om Himmel og
Jord, om Guder og Mennesker, om Familieliv, Stat og sædelig Orden kan i
visse Maader være betydelig udviklet, uden at dertil svarer nogen
videnskabelig Forstaaelse. Hvad der gjælder om Folkeslag, gjælder ogsaa
om Individer. Den organiske Morphologi lærer os, at Individer af høiere
Orden i Fostertilstanden gjentage, om end med Ændringer og Afkortninger,
de Dannelser, som ere charakteristiske for de lavere og laveste Former,
fra hvilke Arten stammer. Noget Lignende gjælder om Bevidsthedens
Udvikling hos det enkelte Menneske. Som Barn vækkes det til at sondre de
udvortes Ting og fornemme sin egen Tilværelse, saa lærer det at bruge
Sproget; det beriger sig med en Skat af foreløbige Erfaringer, og naaer
omsider saa vidt, at en skoleret Undervisning kan begynde. Underviisningen,
der forbereder en videnskabelig Erkjendelse, er dog endnu ikke
videnskabelig. Kundskaber indsamles og ordnes; men kun lidt efter lidt
kommer den umiddelbare Bevidsthed til at besinde sig paa sin Videns
oprindelige Forudsætninger.

Her kræves en klar Adskillelse mellem den Kreds af Billeder og
Forestillinger, der tilhøre Bevidstheden selv, og Tingene i den udvortes
Verden. Individualiteten med sit tænkende Jeg maa, med sikker Takt efter
given Sprogbrug, kunne tale om Former som Tid og Rum; den maa kunne
skjelne mellem Tingen og dens Egenskaber, Materien og dens Kræfter, maa
kunne forstaae, hvad Sammenhæng er, hvad der menes med Udtryk som Aarsag
og Virkning, Tilfældigt og Nødvendigt o.s.v. Paa disse Betingelser vil
den Paagjældende kunne begynde en virkelig Efterforskning af existerende
Ting og deres Love; han vil, hvis han har Kræfter og naturlig Begavelse,
kunne bearbeide et givet Stof med videnskabelig Dygtighed og komme til
værdifulde Resultater. Men først naar han henvender sin Opmærksomhed paa
Betingelserne selv, paa Bevidsthedens oprindelige Forudsætninger, faaer
han at vide, hvad Viden og Videnskab væsenlig er.

En mærkelig Forandring foregaaer med ham. Hvad før var umiddelbart
indlysende, bliver nu pludselig gaadefuldt. Beskæftiget med Tingenes
indbyrdes Forhold og Bevægelser i Rum og Tid, var han ganske optagen af
særlige, exacte Bestemmelser; hvad Rum og Tid er, syntes ham at forstaae
sig af sig selv. Men nu opstaae Spørgsmaal som disse: Hvad er egenlig
Rummet, hvad er Tiden som Tid? vedkomme disse Formbestemmelser Tingene
selv, eller ere de blot nødvendige for vor Opfattelse? Hvad er Materie,
hvad er Kraft, hvad er en Ting, naar vi skille den fra dens Egenskaber,
et Hele, naar det skilles fra sine Dele? Vi synes at see det Materielle,
men Kræfterne kunne vi ikke see; vi synes at sandse Tingen, men det er
kun dens Egenskaber; vi iagttage Tingenes Forandringer, og deres Følgen
med hinanden, efter hinanden; men hvad berettiger os til at antage
Forandringerne som Virkninger af Aarsager, der ligge i Tingene? Hvad er
egenlig det, vi kalde Aarsag, er det noget Sandseligt eller blot noget
Tænkeligt, Noget, der er i Tilværelsen udenfor os, eller Noget, der blot
er i vor egen Hjerne? Naar slige Spørgsmaal paanøde sig Bevidstheden,
naar Grundlaget for den første naturlige og umiddelbare Viden vakler,
vækkes der i Sindet en Trang til at undersøge Videns og Videnskabens
Væsen.

Idet den paagjældende Tænker begynder at prøve sin Videns Grundsætninger,
tvivle om Alt, for igjennem Tvivlen og Uvisheden at arbeide sig frem til
begrundet Vished, er han udsat for at fare vild. Hvad han end vælger til
Udgangspunkt, erfarer han let, at ethvert Fodfæste glipper. Det
forekommer ham, som kunde der ikke være Tvivl om de materielle
Gjenstandes Realitet; men hvad Vished have vi for, at Gjenstandene
existere udenfor vor Bevidsthed? vi vises hen til vor egen Følelse, vort
Syn, vor Hørelse, kort: til vore egne Sandsninger, vore egne
Bevidsthedstilstande, vor egen subjective, nødvendige Opfattelse. Saa
søger han da i Bevidstheden selv efter noget Fast og mener maaskee i den
Sætning: »Jeg tænker, og tænkende er jeg«, at have et urokkeligt
Standpunkt; men han vil snart finde sig skuffet. Hvad vi kalde vort Jeg
er ganske vist en Eenhed, et indre Noget, der ledsager alle vore
Forestillinger og Tanker, men hvad er det for et Noget? For saa vidt det
er en Bevidsthedens Eenhed, maa det jo være ligesaa foranderligt som
Bevidstheden; men at Bevidstheden, der er organiske Betingelser
underlagt, kan fordunkles, formørkes, ja endog ganske forsvinde, er en
uimodsigelig Kjendsgjerning. Han forsøger at gribe det absolut Visse i
Tænkningen og dens Love; men det gaaer ham ikke bedre. Er der nogen
Tankesætning, som Modsigelsen ikke kan angribe? Ja, her er Lighedens,
Identitetens Grundsætning: A er, hvad A er, A = A. Dog nei, ogsaa det er
en Skuffelse. Sætningen A = A grunder sig paa Sætningen: Jeg = Jeg, og
denne igjen paa Forudsætningen: Bevidsthed = Bevidsthed. Men Bevidstheden
er jo foranderlig; dersom den da forandrer sig saa hastig, at den
befinder sig i een Tilstand, naar den sætter A første Gang, og i en anden
derfra forskjellig, naar den sætter A anden Gang, saa bliver det jo
ganske uvist, om vi i det to Gange satte A virkelig have det Selvsamme;
altsaa: Sætningen A = A bliver uvis.

Hvad følger nu heraf? At den efter Viden og Vished søgende Individualitet
finder Uvished overalt og ender i bundløs Skepsis. At fortsætte paa denne
Vei og tage Sagen alvorligt, fører til Fortvivlelse. Den søgende
Individualitet maa standse; og det kommer nu an paa, hvorledes han bryder
af.

Vender han tilbage til den første Umiddelbarhed, bruger de
almeengjældende Forestillinger paa Takt og slaaer sig til Ro ved den
saakaldte sunde Sands, da kan han udvikle sin Forstand i praktisk
Retning, berige sig med Kundskaber, ja han kan endog med Held studere et
eller andet Fag, gjøre Opdagelser og bringe det meget vidt; men hvad
videnskabelig Viden egenlig er, faaer han aldrig at vide. Der har været
ansete Videnskabsmænd, udmærkede ved sjælden Sands og Skarpsindighed,
som, saa længe Talen var om en Kreds af bestemte Phænomener og Love,
kunde udtrykke sig med al en Videndes Sikkerhed, men som, naar der
handledes om Videns almene Grundlag, om den Gyldighed, der i det Hele kan
tillægges den menneskelige Viden, viste sig hildede i skeptisk
Forvirring.

Viden og Videnskab staae i saa væsenligt Forhold til hinanden, at de
forudsætte hinanden i det Uendelige. Man kan ikke blive tilfulde hjemme i
nogensomhelst Videnskab, være sig empirisk eller apriorisk, uden at gjøre
sig Rede for, hvad Gyldighed den menneskelige Viden har, og hvorledes de
Videns Kategorier, som idelig anvendes, oprindelig forholde sig til Ting
og Bevidsthed. Heller ikke vil det lykkes forud at skaffe sig Vished om,
hvad Viden egenlig er, og hvor vidt der gives en sand Viden, før man
indlader sig med Videnskaben selv; thi dersom Undersøgelserne om Videns
Værd ikke gjennemføres paa stræng videnskabelig Maade, kommer man til
falske Resultater. Den blinde Stolen paa faste Grunde og den løse
Tvivlen, som udhuler alle Grunde, ere Exempler paa det Uvidenskabelige.

Viden er et Bevidsthedens Lys; det gjælder om at naae til det Indlysende.
Ligesom der i Tingenes Verden ere mørke og selvlysende Legemer, saaledes
er der i Bevidstheden noget Dunkelt, der maa belyses af Andet, og Noget,
som i sig selv er indlysende. Vi maae da begynde med det umiddelbart
Indlysende og derfra gaae over til det Nærmeste, som deraf kan belyses.
Sammenhængen mellem Hiint og Dette maa selv være umiddelbart indlysende.
Men allerede dette første Skridt er vanskeligt. For at undersøge, hvad
Viden er, maae vi begynde med Viden. Er nu vor Viden ved sig selv
indlysende, saa er her jo Intet at undersøge; skal en Undersøgelse om
Viden selv have nogen Betydning, saa maa der være et Andet, der i og ved
sig selv er mere indlysende end Viden; men hvad skulde dette vel være? Et
vidende Selv, en Videns Gjenstand og begges Forening i een Bevidsthed!

Disse Forudsætninger have ingenlunde den Evidens, som udvortes Ting, der
fremstille sig for Synet; ikke desto mindre vise de sig ved nærmere
Undersøgelse at høre til det første ved sig selv Indlysende. Man
overbeviser sig herom ved Forsøg paa at forklare dem af noget Andet og
mere Indlysende eller af hverandre indbyrdes. Forsøger man at aflede
Viden af Ikke-Viden, saa indseer man let Skuffelsen, thi Adskillelsen af
Viden og Ikke-Viden maa jo foregaae indenfor Viden. Ligesaa meningsløst
er det at ville aflede det Bevidste af det Ubevidste; man maa jo nødvendig
have det Ubevidste i Bevidstheden blot for at kunne sammenligne det med
det Bevidste. Prøver man paa at fastholde Viden for sig uden Indblanding
af nogen som helst anden Bestemmelse, saa ender det med en Viden om
Intet; men en Viden om Intet er Ikke-Viden, og Ikke-Viden bliver da en
negativ Form for Videns Gjenstand. Forsøger man at forklare Viden af den
Vidende, da har man allerede underskudt Viden; og lige saa umuligt er det
at aflede det vidende Selv af Viden; thi for at kunne undvære et Selv, der
har Viden, maatte Viden for sig være vidende om sig selv, og just derved
adskille Viden af sig som vidende fra Viden om alt Andet. Hvor tom og
formel en saadan oprindelig Skjelnen og Adskillen end synes at være, saa
er den dog ikke ørkesløs; den er tvertimod den første og eneste Vei til
at naae videnskabelig Indsigt i det ved sig selv Indlysende.

Hvad her er sagt om Bevidsthed og Viden i reen Almindelighed, gjælder
ogsaa om de særlige Former, hvori Bevidstheden modtager sit Indhold, om
de oprindelige Sandsningsformer, Rum og Tid, om Forstandens eller
Tænkningens Former, de logiske Kategorier o.s.v. Det er ikke for sig
bestaaende Vidensafdelinger, der lade sig ordne i visse Rubrikker, som
kunde man have Sandsningen i een Skuffe og Forstanden i en anden, det er
tvertimod virksomme Former, der betinge hverandre gjensidig, og under
indbyrdes Samvirken indgaae i saa mangfoldige Combinationer og derved
afgive saa særegne Bestemmelser, at Analysen overalt maa gaae fra
Opløsningen af det Sammensatte til det Enkelte: Forstaaelsen er
middelbar.

Desuagtet er der i enhver af disse særegne Almeenformer noget
Eiendommeligt, som ved sig selv er indlysende. Paa hvilke forskjellige
Maader man end søger at udfinde de charakteristiske Bestemmelser ved Rum
og Tid og paavise, hvorledes og under hvilke Betingelser disse
Formbestemmelser fremkomme og udvikle sig i Bevidstheden, saa er det dog
uimodsigelig vist, at det Eiendommelige ved enhver af dem kun kan indsees
af et Udviklingsresultat, derved, at det allerede i Udviklingens
Begyndelse er med som noget ved sig selv Indlysende. Ved Rum forklare vi
Rum, ved Tid Tid, ved Grund Grund, ved Aarsag Aarsag; Forklaringen gaaer
i en Cirkel.

Og dog er Forklaringen ingenlunde ørkesløs; thi under Cirkelbevægelsen
skeer det, at Kredsen for enhver af Grundformerne udvider sig, idet den
alene for sin egen Skyld optager Bestemmelser af de andre i sig. Saaledes
bestemmes Rummet ved Tiden, Tiden ved Rummet, Sandsning ved Tænkning, og
Tænkningen igjen ved Sandsning. Saaledes forklares Bevidsthed og Viden
ved Hjælp af særegne Former, og disse igjen ved Hjælp af Bevidsthed og
Viden. Saaledes forklares Bevidsthedens Grundformer ved et tilsvarende
Indhold, og Bestemmelserne ved det tilsvarende Indhold af Bevidsthedens
Grundformer. Den videnskabelige Fordring fyldestgjøres derved, at ingen
Form, intet Indholdsled, ingen Gjenstand faaer Lov til at blive staaende
isoleret eller blot paa udvortes, tilfældig Maade at føies til de andre.
Alle Bestemmelser, Formbestemmelser og Indholdsbestemmelser,
Sandsebestemmelser og Tankebestemmelser, oprindelige og afledede, maae
indgaae i Bevægelsen saaledes, at de Punkt for Punkt belyse og
gjennemskinne hverandre gjensidig.

\section{Sandsningens Andeel i Viden: udvortes Erfaring}
\label{sec:2}
% pp. 9--19
% The objects of sense must be both inside and outside consciousness
% (the ``both/and'' that resolves the inside/outside antinomy).
% Critique of sensualism: deriving consciousness from ganglion-
% consciousness, spinal-cord consciousness, etc., is incoherent ---
% you cannot derive the non-sensory from the sensory.
% ``Sensualismens Eensidighed: den drukner Bevidstheden i det
% Ubevidste og gjør det Sandselige til det Oprindelige.''

\noindent Endskjøndt det ved sig selv er indlysende, at Viden maa have
sin Gjenstand, saa følger dog ingenlunde heraf, at man af Videns almene
Krav kan aflede Gjenstandens Beskaffenhed. Det er langt fra, at Gjenstand
og Bevidsthed ligne hinanden; Modsætningen er saa fremtrædende, at det
for et umiddelbart Blik livagtig seer ud, som om Gjenstanden stod
selvstændig udenfor Bevidstheden. Er dette nu saaledes? Til Bekræftelse
for Gjenstandens udvortes Selvstændighed støtter man sig paa Sandsernes
Vidnesbyrd; men Vidnesbyrdet er tvetydigt. Opfatter man Bevidsthedens
Kreds som en Cirkel med en bestemt Radius og lader Tingene staae i et Rum
for sig udenfor Peripherien, hvorledes skulle de da kunne sandses? Det
seer ud, som om de speilede sig i Bevidstheden, og forholdt sig til den,
som Gjenstandene til Speilet. At dette er en feilagtig Opfattelse, indsees
let. For at vise sig som sandselig Gjenstand, maa Gjenstanden sandses og
derved optages i Bevidstheden selv. Hele den sandsede Gjenstand befinder
sig indenfor den sandsende Bevidstheds Kreds; er der Noget udenfor, saa
maa det her være det Usandselige eller dog Usandsede, men derom kan den
sandsende Bevidsthed jo Intet vide. Dersom den sandsende Subjectivitet
virkelig kunde see Gjenstandene udenfor sin Bevidsthed i Rummet, da maatte
den jo ogsaa kunne see deres forskjellige Afstande fra
Bevidsthedsgrændsen; Syvstjernen f.\ Ex.\ maatte da vise sig at være
længere borte --- ikke fra det sandsende Individ, thi det forstaaer sig
--- men fra Bevidstheden selv, end Maanen.

Ikke desto mindre maae vi dog indrømme den sunde Sands, at Forudsætningen
af Gjenstandenes Tilværen udenfor Bevidstheden har virkelig Betydning.
Antager man nemlig, at Gjenstandene selv befinde sig indenfor
Bevidstheden, da maa man jo ogsaa antage, at de ere den sandsende
Bevidstheds egne Frembringelser. Men hvis Tingene kun ere til i det
Øieblik vi sandse dem, saa maae de jo ophøre at være til, naar vi ikke
sandse dem; den, der anfaldes af en Røver og er bange for at blive skudt,
behøver da blot at lukke Øinene til, for at tilintetgjøre baade Røverens
og Pistolens Tilværelse. Barnagtigt. Vi komme altsaa til det Resultat, at
Tingene, der sandses, maae være baade indenfor og udenfor vor Bevidsthed.

For at forlige de to, som det synes, modstridende Bestemmelser, er det
nødvendigt at undersøge, hvori Sandsningen egenlig bestaaer. Den
sandsende Bevidsthed er paa een Gang baade activ og passiv, activ, for saa
vidt den maa frembringe Gjenstanden, passiv, for saa vidt den i Et og Alt
maa rette sin Frembringelse efter Gjenstanden. Ifølge Sandsningens
Umiddelbarhed har det Passive en saadan Overvægt, at det sandsende Subject
end ikke mærker sin egen Medvirken, men finder sig udelukkende fortabt i
den sandsede Gjenstands Objectivitet. Antagelsen af et Udenfor beroer paa
Rumsanskuelsen, paa den uafviselige Nødvendighed, hvormed Gjenstanden
paatrænger sig og den slaaende Umiddelbarhed, hvormed den viser sig.
Subjectiviteten kan vende sig bort fra Gjenstanden, kan lukke sine Øine og
tilstoppe sine Øren: hjælper ikke, naar han seer sig om, saa er den der
igjen. Forskjelligartede Sandser, Synet, Følelsen, Hørelsen, ja i visse
Tilfælde endog Lugten og Smagen, kunne forene sig til et samstemmigt
Vidnesbyrd om Gjenstandens Realitet. Viden om Gjenstandens Realitet er en
reel Viden, en Viden af udvortes Erfaring.

Vel forudsætter den udvortes Erfaring ikke alene Sandsning, men ogsaa
Tænkning, Skjelnen, Sammenligning, Dømmen o.s.v.; men Sandsning og
sandselig Iagttagelse er dog Hovedsagen, da Tingen maa tages ganske
saaledes, som den er given. Respekt for det Givne er Respekt for
Kjendsgjerninger. Det er ikke Gjenstanden, der retter sig efter
Bevidstheden, det er omvendt Bevidstheden, der maa rette sig efter
Gjenstanden. Fra et strængt empirisk Synspunkt er det ligefrem
meningsløst at spørge, om Tingene ogsaa ere saaledes, som vi opfatte dem.
Hvor skulde de vel kunne være anderledes, naar vi betænke, hvad Sandsning
og Erfaring egenlig er? Det er jo ikke Bevidsthedens Sandseformer, der
oprindelig bestemme det, der sandses, og raade for, hvorledes Tingen skal
see ud; det er tvertimod den sandselige Ting selv, der bestemmer
Sandseformerne; hvis den ikke formaaer dette, saa er den ikke til at
sandse, d.v.s.\ er ikke en sandselig Ting. Den Gjenstand, som efter sin
Naturbeskaffenhed ikke kan sees, hører ikke til de synlige Gjenstande;
den, som ikke kan høres, ikke til de hørlige. Det er netop de sandselige
Bestemmelser ved Gjenstanden, der gjør den sandselig. Men hvad er
sandseligt? Dog vel det, der lader sig sandse. Og hvad er sandseligt til
for os uden det, som nøiagtig svarer til vore Sandseformer og kan bestemme
disse efter sin Beskaffenhed? Enhver Ting maa bestemmes ved sine
Egenskaber. At antage Ting, som falde udenfor alle Egenskaber, er at
antage Ting, der falde udenfor alle Bestemmelser; at antage sandselige
Ting, som ikke have sandselige Egenskaber, er at antage Ting, hvis
Sandselighed bestaaer i at være usandselig. Meningsløst!

Men, siger man, Erfaringstingene kunne jo maaskee nok have sandselige
Egenskaber og for saa vidt være sandselige i og for sig, uden at vi med
vore Sandser ere i Stand til at opfatte dem saaledes, som de virkelig ere.
Hvad vi sandse som udstrakt, er maaskee slet ikke udstrakt, eller dog ikke
udstrakt i den Form, hvori vi sandse det; hvad vi antage for at være
haardt, kan maaskee være blødt; hvad vi ansee for uigjennemsigtigt, kan
maaskee være gjennemsigtigt. --- Vi svare: vel muligt, men hvad vedkommer
det os?

Den udvortes Erfarings Gjenstande maae vi nødvendig opfatte saaledes, som
vi erfare dem, at drage Erfaringens Gyldighed i Tvivl er at drage alle
Kjendsgjerninger i Tvivl og give Afkald paa reel Viden. Jo mere
Realitetssands den Vidende har, jo større vil derfor ogsaa hans Interesse
for en skarp Iagttagelse af Tingene i alle Enkeltheder være. Hvad der
ellers i Almindelighed antages for rimeligt eller urimeligt, bekymrer ham
ikke; er Tingen virkelig til, saa maa den tages, som Erfaringen giver den,
om dens Tilværelse end iøvrigt kan forekomme os nok saa urimelig. Videns
Grundlag er, som heraf sees, et Grundlag af Kjendsgjerninger; en
ufravigelig Betingelse for at opfatte og sikkre sig de udvortes
Kjendsgjerninger, er at faae dem fuldstændig sandsede: dette er
Sandsningens Andeel i Viden.

At man desuagtet ikke kan aflede al Viden af Sandsning og sandselige
Kjendsgjerninger, er indlysende. Til Sandsning hører der Bevidsthed, men
Bevidsthed er ingen Sandseting, ingen Gjenstand for udvortes Erfaring. Jo
mere eensidig man vænner sig til at anlægge den udvortes Erfarings
Maalestok paa al Viden, desto større Hang vil man ogsaa faae til at
oversee Bevidsthedens usandselige Væsen, og udlede Sandsningen selv af
empiriske Egenskaber ved Sandseorganerne. Istedenfor at forklare
Sandsningen af Bevidstheden, lægger man nu an paa at aflede den af
Sandseorganernes Virksomhed, og denne Virksomhed igjen af den
eiendommelige Bygning, som kan være Gjenstand for udvortes Erfaring. Saa
taler man om Gangliebevidsthed, og gjør Bevidstheden til en physisk
Egenskab ved Gangliet, om Rygmarvsbevidsthed, om Hjernebevidsthed, ikke
blot i den Betydning, at Ganglier, Rygmarv og Hjerne ere nødvendige
Betingelser for den sandsende Bevidsthed, men udtrykkelig i den Mening, at
disse Organer selv ere sandsende Subjecter. Det Urimelige heri falder i
Øinene, naar man skal til at gjøre Rede for, hvorledes de mange
forskjellige Bevidstheder kunne flyde sammen i een eneste
Bevidsthedseenhed. Ved at sætte det Forbindende i Bindevæv og Nervetraade
opnaaer man kun, hvad for Øvrigt er meget fortjenstligt, at paavise
betingende Mellemled; men Sagen selv faaer man ikke opklaret. Kan
Bevidsthed løbe igjennem en Nerve som Elektricitet igjennem en Traad? Idet
Elektriciteten strømmer igjennem Traaden, bliver Traaden elektrisk; gaaer
det nu an at sige, at en Nerve selv bliver sig bevidst, idet Bevidstheden
gjennemstrømmer Nerven? Det hjælper ikke her at bryde af og erklære, at vi
i Grunden lige saa lidt vide, hvad Bevidsthed er, som hvad Elektricitet
er. Man sætte Grændsen for den menneskelige Viden, hvor man vil, saa er
der dog Et, som Videnskaben med Strænghed kræver af den Vidende, dette
nemlig, at han maa vide, hvad det er, han selv siger. Vide vi ikke, hvad
der menes med Ordet Bevidsthed, hvorledes kan man da paastaae, at der er
Bevidsthed i Ganglie, Rygmarv og Hjerne, ja ovenikjøbet gjøre
Bevidstheden til en Egenskab ved disse Organer? Tingen er, at man stirrer
sig blind paa sandselige Kjendsgjerninger, vil aflede det Ikke-Sandselige
af det Sandselige; men jo mere man forsøger at gjøre Alt tydeligt indtil
de mindste Enkeltheder, desto mørkere, gaadefuldere bliver det Hele. Det
er Sensualismens Eensidighed, at den drukner Bevidstheden i det Ubevidste
og gjør det Sandselige til det Oprindelige. Hvad der i Sensualismens Øine
skal have reel Værdi, maa være en Ting. »Selvet, Jeget, er ikke en Ting,
altsaa er Selvet, Jeget, egenlig ingen Ting«.

En Sætning, som den, at vi lige saa lidt vide, hvad Bevidsthed er, som
hvad Elektricitet er, viser sig, naar vi tage det Tvetydige bort, at være
grundfalsk. Saalænge Legemernes indre, molekulære Bygning, Smaadelenes
Omsætning, Leiring og Vexelvirkning endnu ikke er tilstrækkelig undersøgt,
kan det med Rette siges, at Elektriciteten har noget Gaadefuldt ved sig,
Noget, vi ikke kunne gjennemskue; men kan det Samme siges om Bevidstheden?
Ingenlunde! Vi have i foregaaende Paragraph seet, at Bevidstheden netop
afgiver den oprindelige Forudsætning for al vor Viden og hører til det
første, ved sig selv Indlysende. Hvor misligt det er at glemme eller dog
oversee dette, giver sig endog tilkjende paa Erfaringsvidenskabens eget
Omraade.

Det forstaaer sig af sig selv, at Erfaringslæren under ingen Form kan
undvære Tænkning; men for den sensualistiske Erfaringslære er det netop
charakteristisk, at den søger sin egenlige Forvisningsgrund i Sandsningen,
sætter Realiteten selv i det Sandselige, og lader Tænkningen være en blot
Formsag. Naar Sensualisten taler om Tingen og dens Egenskaber, det Hele og
dets Dele, da er det ingenlunde Tænkningen, der forvisser ham om, at hvor
der er Egenskaber, maa der ogsaa være Ting, at hvor der er et Hele, maa
der ogsaa være Dele; nei det er Sandsningen. Han mener ligefrem at kunne
sandse Tingene som Ting, Egenskaber som Egenskaber, det Hele for sig og
Delene for sig, og ligeledes ved Sandsningen forvisse sig om Egenskabernes
Forbindelse med Tingen, om Delenes Sammenhæng i det Hele. Men naar
Tænkningen begynder at røre sig frit, saa fremkommer Uvisheden.
Sensualisten bliver tvivlraadig, hvad skal han gjøre? Han troede nu saa
vist, at have sandselige Garantier for, at Guld er til som en Ting med
sine Egenskaber, et Hele med sine Dele, og at hans Viden om, at Guld er et
ædelt Metal, hvilede trygt paa udvortes Erfaring. Han veed jo, at Guldet
er guult, at det kan krystallisere i Octaedre og lade sig udvalse til
Blade af en $\frac{1}{20000}$ Linies Tykkelse. Sandselige Iagttagelser
have overbeviist ham om, at Guldet iblandt alle Metallerne har den mindste
Affinitet til Ilt, at det lader sig opløse i Kongevand o.s.v. Men hvad
Vished har han nu egenlig om alt dette? Han har villet sandse hver
Egenskab for sig; men den enkelte Egenskab lader sig ikke sandse
umiddelbart, og det af den simple Grund, at den ikke kan existere for sig
umiddelbart. Da han sandsede Egenskaben guult, maatte han henføre den
deels til Guldet, deels til Lyset, deels til sin egen Synsfornemmelse. Er
der nu nogen umiddelbar Sandsning, der kan lære ham, hvilken Andeel
Guldet, Lyset og Lysfornemmelsen hver især have i dette Guult? Han har
uden videre sat forskjellige Egenskaber sammen til en reel Eenhed i Ting
og Substans, men den umiddelbare Sands kan ikke lære ham Noget om en
saadan Eenheds Realitet. At være guult, at være smidigt, at krystallisere
i Octaedre, tyde i Virkeligheden hen paa saa uligeartede Ting, at ingen
Sandsning kan forvisse os om, at de absolut maae høre til een og samme
Ting. Tvertimod! den udvortes Erfaring lærer, at de samme eller dog
lignende Egenskaber kunne findes hos forskjellige Ting.

Sensualisten er for Erfaringens Skyld nødsaget til at antage en Delenes
Sammenhæng og en vis Orden, hvori Tingene maae følge med hverandre og
efter hverandre, men hvad Vished har han herom? Han er nødt til at betjene
sig af Ord som Nødvendigt og Tilfældigt, Grund og Følge, Aarsag og
Virkning, det Individuelle og det Almene, Love o.s.v. Men han kan ikke
have nogen egenlig Tillid til slige Bestemmelser, han kan ikke tillægge
Ordene reel Betydning. Den umiddelbare Sandsning opfatter det umiddelbart
Givne i dets Enkelthed og Tilfældighed og indeholder ingen Garanti for,
hverken at det fremdeles vil bestaae, ei heller at det som Led i en Række
formaaer at udøve nogen Indflydelse paa de andre Led. Ved at underlægge
samtidige Phænomener en nødvendig Forbindelse, ved at antage et
Forudgaaende for Aarsag og et Efterfølgende for Virkning, gaaer man ud
over det, som umiddelbart kan sandses, og da Sandsningen jo dog skal være
det Eneste, der kan give tilstrækkelig Vished, maa man enten reducere den
udvortes Erfarings videnskabelige Udbytte til en løs Mangfoldighed af
Phænomener eller underskyde Forudsætninger, som man dog ikke kan
vedkjende sig. Ved udelukkende at ville holde sig til en sandselig Vished
drives man over i Uvished, og Skepticismen staaer for Døren.

Den eensidige Empiris Tilhængere gjøre kun det Onde værre, naar de ved at
henvise os til Vanen, Gjentagelsen, Ideeassociationen og slige Palliativer
forsøge at tilsløre det Uvisse ved Hjælp af det Halvvisse, og give det
Skin af, at det er den sandselige Viden, som ved at indføre flere og flere
Mellembestemmelser lidt efter lidt arbeider sig frem og gjør
Tankebestemmelserne almeengyldige, medens dog Sandheden er, at den Punkt
for Punkt maa forudsætte noget Andet end Sandsningen. Sandsningen er et
Moment, et uundværligt Moment i al vor Viden om sandselige Ting; men naar
dette Vidensmoment skal udvides til at gjælde for Alt og udgjøre den hele
Viden, saa indsmugler man det Oprindelige, den vidende Bevidsthed, det
vidende Selv, medens man dog fornegter det Oprindelige; og saa fremkommer
det Uvidenskabelige.

\section{Indbildningens Andeel i Viden: indvortes Erfaring}
\label{sec:3}
% pp. 19--29

\noindent Ifølge sædvanlig Sprogbrug skulde man mene, at Indbildning og
Viden aldeles Intet havde med hinanden at gjøre. »En Indbildning er
Noget, som ingen Realitet har, og en indbildt Viden det Samme som en
falsk Viden«. Sagen vil dog vise sig i et andet Lys, naar vi besinde os
paa, at Indbildning væsenlig betyder en ideel, dog udad gaaende Dannelse,
en i vort Indre frembragt Dannelse, en Billeddannelse, der gjør
Bevidsthedens active Side kjendelig og svarer til, hvad vi kalde
Indbildningskraft. Indbildningskraftens Virksomhed bestaaer netop i at
danne Indbildninger; hvorledes disse Indbildninger forholde sig til den
reelle Væren, skal særlig undersøges; her bruge vi Indbildning og
Frembringelse af Indbildningskraft som eenstydige Udtryk.

At der imellem Sandsning og Indbildning maa være en nøie Forening, er let
at indsee. Naar vi efter at have sandset en Gjenstand, vende os bort fra
den, bliver dens Billede endnu staaende for Bevidsthedens Øie, om end
mattere og blegere. Dette Billede vidner om, at Indbildningskraften
allerede ved den umiddelbare Sandsning af Gjenstanden selv har været
virksom. Hvis Sandsningen af den virkelige Gjenstand ikke foregik ved en
billeddannende Virksomhed, vilde Huskning og Erindring være umulige.
Under Sandsningen selv er Billeddannelsen vistnok saa afhængig af Indtryk
fra Gjenstanden, at Indbildningen slet ikke mærkes; men dersom den
nærværende Gjenstand kunde opfattes som den viser sig for Bevidstheden
uden en ubevidst Billeddannelse, da vilde dens Efterbillede heller ikke
kunne danne sig; Bevidstheden beholdt Intet tilbage: Gjenstanden vilde
kun være til for den i Sandsningens Øieblik. Og hvad vilde saa Følgen
heraf være? At en Bevidsthed om flere forskjellige Gjenstande vilde være
umulig. Ja, Sandsningen selv vilde være umulig. Den sandselige Gjenstand
maa som et Hele bestaae af Dele. For at sandse et Hele, hvis Dele være a,
b, c o.s.v., maae vi begynde med at sandse en Deel, f.\ Ex.\ a; dersom vi
nu, idet vi sandse b, have tabt a, og, naar vi komme til c, baade a og b,
saa naae vi jo aldrig saa vidt, at vi faae Delene sammenfattede. Men naar
vi ikke kunne sammenfatte Delene i en Sandsnings Eenhed, hvorledes skulle
vi da kunne sandse det Hele?

Et slaaende Beviis paa Sandsningens og Indbildningens uopløselige Eenhed
ligger i Anskuelsen. Det forholder sig ganske vist med Anskuelsen som med
Rummet, vi maae forklare Anskuelse ved Anskuelse ligesom Rum ved Rum, det
vil sige: Anskuelsen hører med til det første ved sig selv Indlysende.
Ikke desto mindre er det Sandsningen af virkelige Gjenstande, der giver os
Betingelser for en Tydeliggjørelse af, hvad Anskuelsen egenlig bestemmer,
og hvad det er, der bestemmer Anskuelsen. Naar det hedder, at vi, for at
sandse et Hele, maae begynde stykkeviis med at sandse dets Dele, kunne vi
omvendt sige, at vi for at sandse en Deel, maae begynde med at sandse et
Hele; d.v.s.\ at vi for Sandsningens Skyld maae begynde med Anskuelse.
Hver mindste sandselige Deel har en Udstrækning, bestaaer af Dele, og er
altsaa et Hele: vi kunne ikke sandse Differentialer. Det er den
umiddelbare Forening af Anskuen og Sandsen, der bevirker, at vi paa een
Gang synes at sandse baade det Hele og dets Dele.

Uden Indbildning er det umuligt at sandse nogen Forandring i det Udvortes.
For at komme til Bevidsthed om, at en Ting forandrer sig, maae vi i det
Mindste sandse den to Gange og ved at sandse den anden Gang være i Stand
til at erindre, hvorledes den saae ud første Gang. Sammenligningen er
betinget af, at Sandsning omsættes i Indbildning, og Indbildning i
Sandsning. Paa Sandsningens Maade opfattes Forandringen derved, at vi
sammenligne Ting med Ting, hvorved det kommer til at see ud, som om
Tingene selv viste sig snart saaledes, snart anderledes; med andre Ord: vi
finde Tingene forandrede, fordi de virkelig forandre sig. Tiden, den Form,
i hvilken Forandringen foregaaer, gjælder da lige saa vel for Tingene, som
for vor Opfattelse af Tingene. Men kun ved Hjælp af Indbildningen kan
Forandringen opfattes saaledes, at ikke blot Gjenstandene sammenlignes med
deres Billeder, men at de forskjellige Gjenstandsbilleder sammenlignes med
hinanden indbyrdes. Under denne Sammenligning blive de frigjorte Billeder
igjen forandrede til Forestillinger; Forestillingernes Mangfoldighed er
Gjenstand for indvortes Erfaring.

Det Spørgsmaal, hvorfra saadanne Kategorier som: Indre, Mulighed, Aarsag
o.s.v. have deres Oprindelse, kan i denne Sammenhæng foreløbig besvares.
Naar der tales om Tingenes Indre, er det jo indlysende, at dette Indre
ikke kan henføres til nogen udvortes Erfaring. Tage vi det umiddelbare
Ydre bort, saa er der atter for Sandsningen et nyt Ydre. Hvorfra faae vi
da det Indre? Fra Indbildningen. Vi digte det ind i Tingen, og komme ved
Modsætningen til Bestemmelsen af et Ydre. Det Samme gjælder om
Muligheden. Hvad der skal paavirke den udvortes Sands maa være virkelig
til; det Mulige kan altsaa ikke komme udvortes, men maa komme indvortes
fra, d.v.s.\ fra Indbildningen. Ere nu Indre og Mulighed en Indbildningens
Tilgift til hvad der ved Sandserne maa bestemmes som Ydre og Virkelighed,
saa kommer det da an paa at indsee Forbindelsen mellem Sandsning og
Indbildning. Uden Indbildning ingen Sandsning. Vistnok er enhver Erfaring
af det Udvortes betinget af lutter Sandseindtryk; men intet Sandseindtryk
kan i reen, umiddelbar Enkelthed opfattes for sig alene; de mangfoldige
Enkeltindtryk, være sig samtidige eller successive, maae samles og
bestemmes ved en tilsvarende Mangfoldighed af Fornemmelser og
Forestillinger, der udspringe af Indbildningens Eenhed i Selvet. Betegne
vi Sandseindtrykket ved $i$, Forestillingen ved $f$, da giver Rækken $i$,
$i$, $i$ os ingen Erfaring, Rækken $f$, $f$, $f$ heller ingen. Rækken $i$,
$f$, $i$, $f$ derimod gjør en Erfaring mulig; thi nu kan $i$ bestemmes ved
$f$, og omvendt, en Sammenligning kan finde Sted, og Forskjellighederne
vise sig: $i'$, $f'$, $i''$, $f''$. Sammenfatningen beroer paa den
anskuende, forestillende Bevidstheds Eenhed.

Sandsningen af det Ene vækker uafbrudt en Forestilling om det Andet.
Sandsningen af A vækker Forestillingen om B, og denne forbereder igjen en
Opfattelse af C. Synet af et Blad vækker Forestillingen om en Stængel. I
Forestillingen ligger en Forventning, deels om Noget, der allerede er,
skjøndt det endnu ikke er sandset, som naar Synet af Forsiden vækker
Forestilling om Bagsiden; deels om Noget, som ikke endnu er, men vil
komme, som naar Synet af et Lyn vækker Forestilling om et Tordenskrald.
Forventningen er uvilkaarlig, og Indbildningen kan skuffe. At der maa
følge Noget, er vist, men om just det Forventede følger, er uvist; der kan
jo følge noget Andet. Nu er Muligheden bleven til i Indbildningen.
Saalænge Tingenes Omskiftelser blot sees at foregaae i Rum og Tid, er
Forestillingen om et Indre endnu ikke givet. Leddene A, B, C kunne
forandre deres Stilling til hverandre, uden at noget af dem forandrer sig;
men foregaaer der under Omsætningen en Forandring enten med eet af Leddene
eller med dem alle, saa faaer Indbildningen en ny Opgave. B forandres,
naar A træder til, B og C blive begge forandrede, naar de træffe sammen:
her maa forudsættes et Indre. Før og efter Forandringen er B det samme, og
dog ikke det samme. Var det ikke det samme, saa kunde man jo ikke sige, at
Forandringen er foregaaet med B, at det netop er B, der er blevet
forandret. Var B efter Forandringen endnu ganske det samme, saa var det jo
ikke undergaaet nogen Forandring. Det Anstødelige heri hæver Indbildningen
ved at henføre det forskjellige Ydre til et fælles Indre. Før havde B eet
Udseende, nu har det et andet, men i Grunden d.v.s.\ efter sit Indre er
det dog det samme. Den Sandsende bliver sig bevidst, at B efter
Omstændighederne kan vise sig med forskjelligt Ydre i forskjellige
Tilstande: Vand kan blive til Iis, og Iis til Vand, det maa ligge i
Vandets, i Isens Indre. Hvad det Indre er, bliver foreløbig bestemt ved en
Mulighed, men hvad det er for en Mulighed, maa den paafølgende Virkelighed
afgjøre. Synet af et Frø vækker Forestillingen om en Plante; Frøet er
Plantens Mulighed, men kun en Mulighed, thi det er jo ogsaa muligt, at en
Fugl kan æde Frøet; Synet af et Frø kan ogsaa vække Forestilling om en
Fugl. Mulighedsforestillingen er for saa vidt subjectiv, men ved at lægge
Muligheden ind i Frøet har Indbildningen fraskrevet sig sin Ret over den
og gjort Forventningen afhængig af det Sandsede; Indbildningen har, med
andre Ord, gjort Muligheden objectiv, d.v.s.\ lagt den ind i det
Virkelige selv.

Den ved Erfaring stadfæstede Forening af Virkeligt og Muligt, en Forening,
der selv er betingende for Erfaringens egen Mulighed, leder til
Indbildning om Aarsag. At A er Aarsag, og B Virkning, kan ingenlunde
sandses umiddelbart. Ved at gjøre Sandsningen til Erfaringens Et og Alt
kan man let forsvare den Paastand, at Aarsagsforestillingen er en
Indbildning, hvortil intet Virkeligt svarer. Vi kunne ikke see paa Ilden,
at den vil brænde os, og dog er der Ingen, end ikke Skeptikeren, der har
Lyst til at stikke sin Finger i Flammen. Hvorfor ikke? Synet af Ild vækker
Forestilling om Smerte. Hver Gang vi prøve paa at holde Fingeren i
Flammen, erfare vi, at det svier; hvorfor tør vi da, trods idelige
Gjentagelser, ikke ansee det for vist, at Ilden er Aarsag og Smerten
Virkning? Fordi, hedder det, man ikke kan see Aarsagen sidde i Ilden.
Vistnok kan man føle Smerten, men for at føle den som Virkning, maatte man
føle Ilden som Aarsag, og det er just det, man ikke kan.
Aarsagsforestillingen er en Indbildning. Denne Synsmaade er i reel
Forstand uigjendrivelig. Man kan lige saa lidt hævde Aarsagens Realitet
ved den blotte Sandsning, som ved den blotte Tænkning. Tænkningen kan
godtgjøre, at der maa være en Aarsag, saa sandt der er en Virkning; men,
hvis der nu ingen Virkning er, saa er der heller ingen Aarsag. Det
Videnskabelige beroer ikke paa at udelukke Indbildningen, men det beroer
paa exacte Bestemmelser af den indre Erfarings Andeel i den ydre, det vil
sige: paa exacte Bestemmelser af Indbildningens Andeel i Viden.

At Indbildningskraften, selv naar den støtter sig til udvortes Erfaringer,
ikke formaaer at grundlægge nogen virkelig sand eller selvgyldig Viden,
ligger i Sagens Natur og har vist sig i Kulturens Historie. Et Folks
Kultur begynder, som ovenfor antydet, ikke med Videnskab og videnskabelige
Undersøgelser, men med Myther og Mythologi. I Mythologien har
Indbildningskraften en fri Tumleplads uden endnu at fatte sin egen Frihed:
den digter Guderne og bilder sig dog ind, at det er udvortes, sandselige
Skikkelser, den beskriver; den digter det Overnaturlige og veed ikke af,
at det er lutter naturlige Bestanddele, den har benyttet. Lidt efter lidt
vækkes jo vistnok en Anelse om, at der er Forskjel mellem Digt og
Virkelighed, mellem Indbildningens Verden og den haandgribelige Tingenes
Verden, der er Gjenstand for udvortes Sandsning. Men den indre Erfaring af
Sindets Tilstande, af Trang og Længsel, af Haab og Frygt, er saa
overvældende stærk, at det udvortes Virkelige, farvet af Indbildningen,
overalt bliver omdannet til Middel for et bevidst Indvortes. Selv da
Oldtidens Bevidsthed var kommen til Besindelse og blevet opmærksom paa
Noget, som er anderledes tilforladeligt end enten sandselige Ting eller
Indbildninger og vexlende Forestillinger, nemlig Tanken, der gjennemskuer
Phænomenet, fastholder Væsenet og fatter det Almeengyldige, selv da var
man ikke i Stand til at sikkre sig en Viden om Tilværelsens virkelige
Indhold. Enten brød man Staven over hele den udvortes, sandselige Verden,
fornægtede al Mangfoldighed, Vorden og Bevægelse, og holdt sig, for at
undgaae Indbildninger, saa eensidig til det Tænkelige med Udelukkelse af
alt Forestillet, at de rene Tankebestemmelser, de tomme Almeenformer,
tilsidst opslugte hinanden i Modsigelsen, og det værende Ene, der blev
tilbage, forvandledes til et absolut Bestemt, der var blottet for alle
Bestemmelser. Eller, naar man indsaae, at det Sandselige, Foranderlige og
Forgjængelige dog maatte have en Art Væren, lod man denne
Tilsyneladelsernes Væren løse sig op i Begreber, selvstændiggjorde
Begreberne som i sig selv værende Ideer og opfattede Tingenes Verden som
en Ideernes Afskygning. Saa dannede der sig en Videnskab, men Videnskaben
blev speculativ, en Viden af væsenlige Indbildninger. Hvormeget end
Tænkningen anstrængte sig for at hævde det evigt Gyldige, kunde man dog
ikke naae Realiteten. Realiteten blev miskjendt, fordi Sandsningen og det
Sandseliges Betydning blev miskjendt. --- Speculation er Idealisme, en
Omdannelse af det for Bevidstheden Fremmede til Lighed med Bevidsthedens
eget Væsen, en Viden, hvori Gjenstanden og den Vidende ere Speilinger af
Et og det Samme, en Viden i Indbildning.

Og dog er Indbildningen lige saa nødvendig for Videnskaben som
Sandsningen, Speculationen lige saa nødvendig som Empirien, Idealiteten
lige saa nødvendig som Realiteten. Den Antagelse, at man ved udelukkende
at see paa Kjendsgjerninger og holde sig til det Reelle, kan frigjøre sig
for Indbildningen, er selv en Indbildning. Man kan f.\ Ex.\ ikke undvære
Sproget. Sproget er af idealistisk Oprindelse, der stikker en Indbildning
under ethvert Ord. Ordet skal svare til Gjenstanden, og dog er der jo
ingen Lighed mellem Ordet Træ og det virkelige Træ, Ordet Huus og det
virkelige Huus. Ordet er et Tegn, Baandet mellem Tegnet og det Betegnede
en Indbildning; men denne Indbildning er aldeles uundværlig. Hvad er
Intet? Et Ord! Ordet Intet og Intet selv have ingen Lighed med hinanden,
d.v.s.\ Ligheden er en Indbildning. Og dog er denne Indbildning aldeles
uundværlig for Tænkningen. Hvis Tænkningen ikke ved Indbildningens Hjælp
var i Stand til at fatte Intet, saa vilde den heller ikke være i Stand til
at fatte Negationen som væsenlig Grændsebestemmelse, ikke i Stand til at
hævde Forskjellen mellem Muligt og Umuligt, ikke i Stand til at opfatte
Lovene.

Indbildningen er en falsk Fører, en upaalidelig Veiviser; men netop i dens
Upaalidelighed maae vi søge dens Fortrin, dens Uundværlighed.
Indbildningens Upaalidelighed er nemlig Eet med dens Bøielighed, dens
Føielighed, dens uendelige Smidighed. Den gjør med Bevidstheden Alt, hvad
den vil, thi den er lutter Mulighed; den viser Bevidstheden ud over sig
selv og ind i sig selv, gjør den til Slave af alle Ting og til Herre over
alle Ting; den leger med Forskjellen og leer af Modsigelsen, sætter
Grændser overalt og udsletter alle Grændser, gjør det Sandselige
tænkeligt, det Tænkelige sandseligt; den er, med andre Ord, den vidende
Bevidstheds uendelige Frihed.

Den formelle Frihed er en ufravigelig Betingelse for at erkjende
Nødvendigheden; Viden maa, for Nødvendighedens Skyld, være fri. Det er til
videnskabelig Indsigt ikke nok, at den Vidende forefinder en Tingenes
Verden udenfor sin egen Bevidsthed, han maa vide, hvad dette Udenfor
egenlig betyder; men for at faae dette at vide, maa han ved Indbildningens
Hjælp kunne stige over Bevidsthedens Muur, maa kunne forestille sig,
hvorledes det vilde see ud, naar han stod paa den anden Side. Naar det
Urimelige, det Forvirrede da viser sig for ham, vender han tilbage til
Sandsningens Umiddelbarhed, lægger tilbørlig Vægt paa Sandsernes
Vidnesbyrd og anerkjender Tingenes Realitet. Det er i videnskabelig
Forstand ikke tilstrækkeligt, at den Vidende iagttager Tingene saaledes
som de i Kjendsgjerning følge med og efter hverandre, og mærker sig
Forskjellen mellem sædvanlige og usædvanlige Tilfælde; han maa paa
Indbildningens Vinger hæve sig høit over Erfaringens Tilfælde, og først
naar han har besindet sig paa, at det Hele vilde være fuldstændig
meningsløst, dersom der i alle disse udvortes Forandringer ikke var noget
Uforanderligt, en Lov, hvorefter det Ene fremgaaer af det Andet, og følger
paa det Andet, da først lærer han at bøie sig for Tankens Nødvendighed og
indrømme dens Bestemmelser af Love, Grunde og Aarsager oprindelig
Gyldighed.

\section{Apriorisk Viden}
\label{sec:4}
% pp. 29--40

\noindent Forskjellen mellem apriorisk og aposteriorisk Viden synes paa
den simpleste Maade at kunne betegnes saaledes: al den Viden, vi hente fra
vor egen Bevidsthed, og have forud for en særlig Viden om Tingen, er
apriorisk; hvorimod den Viden, der forudsætter en Sandsning af existerende
Gjenstande, er aposteriorisk. Grundforestillinger som Eenhed og Fleerhed,
Lighed og Ulighed, Storhed og Lidenhed, Skjønt og Uskjønt, Godt og Ondt
o.\ fl., ere aprioriske. Naar vi sige om een Handling, at den er god, om
en anden at den er ond, saa hente vi ikke vore Forestillinger om Godt og
Ondt fra Handlingerne, men vi henføre Handlingerne til Grundforestillinger,
vi bære i os selv. Naar vi kalde een Gjenstand skjøn, en anden hæslig, da
laane vi ikke Maalestokken for Skjønt og Hæsligt fra Gjenstandene, men fra
Forbilleder i vor egen Bevidsthed. At opfatte to Gjenstande deels som
lige, deels som ulige, forudsætter en Sammenligning; men en Sammenligning
vilde være umulig, dersom vi ikke havde Grundforestillingerne: Lighed og
Ulighed forud, d.v.s.\ a priori. Er Talen derimod om saadanne Gjenstande
som Stene, Jordlag, Planter, Dyr o.\ dsl., saa vender Forholdet sig om,
saa maa Gjenstanden gaae forud for vor Forestilling, vort Begreb om den.
De Forestillinger og Begreber, som saaledes hentes fra Tingene og
desaarsag komme bagefter Tingene, kalde vi aposterioriske, empiriske.

Men naar Grændsen mellem Apriorisk og Empirisk skal opfattes bestemt og
nøiagtigt, indtræde Vanskeligheder. At lægge Eftertrykket paa det
Aprioriske er charakteristisk for Idealismen, at gjøre det Empiriske til
Et og Alt er charakteristisk for Sensualismen. Ved at underkaste
Idealistens »medfødte Ideer« en skarp Kritik kan Sensualisten let bilde
sig ind, at det Aprioriske kun er et Skin, at al vor Viden lader sig
aflede af Sandseindtryk, Sandsninger, Afspeilinger af det Sandselige,
altsaa paa empirisk Maade.

»Skulde vor Viden«, siger Sensualisten, »virkelig stamme fra medfødte
Ideer, da maatte disse Ideer ikke blot være i Sjælen som Sjæl, men, for
Videns Skyld, udtrykkelig i Bevidstheden. Den Menneskesjæl, der er naaet
saa vidt, at den er bleven sig bevidst, maa da med det Samme være bleven
sig disse Ideer bevidst; thi at ubevidste Ideer ere i Bevidstheden, er jo
en Modsigelse. Men nu lærer vor indre Erfaring os, at de saakaldte første
Ideer have dannet sig lidt efter lidt, at en Mængde Sandseindtryk,
Fornemmelser og Enkeltforestillinger ere gaaede forud, før vi kunde danne
os Ideer, der, om de end synes enkelte, dog i Grunden ere meget
sammensatte. Hvor mange lige og ulige Indtryk, Fornemmelser af Behageligt
og Ubehageligt, af Hvidt og Sort, af Haardt og Blødt, maa der ikke være i
en Bevidsthed, før den kan naae til Fællesbenævnelserne: Lighed og
Ulighed, Haardhed og Blødhed. Fællesbenævnelser ere Abstractioner; for at
Bevidstheden kan abstrahere, maa den have Noget, hvorfra den kan
abstrahere. Det kan altsaa ikke være Abstractionen, men Sandsningen, der
er det Første. Hemmeligheden med de medfødte Ideer bestaaer simpelhen
deri, at den sandsende, for Sandseindtryk modtagelige Sjæl abstraherer det
Fælles eller Almene af det Individuelle, selvstændiggjør det Abstraherede
og bilder sig ind, at den i det Afledede besidder det Oprindelige. --- De
udvortes Erfaringer lære os det Samme. Gaves der medfødte Ideer, da maatte
vi finde dem hos Alle med nogen Bevidsthed; men at Børn og Vilde ingen
Anelse have om slige Ideer, er vitterligt. En Australneger f.\ Ex.\ kan
have et skarpt Øie for lige og ulige Ting i hans Omgivelser, men for
Lighed og Ulighed \textit{in abstracto} har han intet Blik. Siig ham, at
to Størrelser, som hver for sig ere lige med een og samme tredie, ere
indbyrdes lige: han vil ikke forstaae det. Vælge vi Ideer som Skjønt og
Hæsligt, Godt og Ondt, hvorom de Fleste dog have en Forestilling, saa
bliver det os ret indlysende, at disse Ideer ikke kunne være medfødte.
Vare de medfødte, saa maatte de jo være eens hos Alle, have samme
Gyldighed for Alle; der kunde da ikke opstaae nogen Strid om, hvad der i
og for sig er skjønt og hvad der er hæsligt, hvad der i og for sig er godt
og ondt. Men at der herom hersker stor Uenighed ikke blot iblandt Uvidende
og Halvvidende, men især blandt dem, der bedst skulde kjende Ideerne, er en
ubestridelig Kjendsgjerning.«

Har Sensualismen paa denne Maade berøvet Sjælen alle medfødte Ideer og
feiet det Aprioriske ud af Bevidstheden, maa Undersøgelsen begynde for
fra; og hvad bliver saa det Oprindelige? Erfaringen! Men for at erfare,
maa man jo sandse. Sandsningen! Men Sandsningen forudsætter en Sjæl, en
Bevidsthed, som kan sandse; Sandsning er, som ovenfor viist, uadskillelig
fra Indbildning, og den skjelnende Indbildning igjen uadskillelig fra
Tænkning. Men, naar al Erfaring nødvendig maa forudsætte en sandsende,
forestillende, tænkende Bevidsthed, kan Erfaringen jo umulig selv være det
Første. Vi faae det Aprioriske tilbage; dog ikke i Form af medfødte Ideer,
men som Evne, som Anlæg. Alt kommer nu an paa at udfinde, hvori Evnen
bestaaer, og at paavise de forskjellige Bevidsthedsmaader, paa hvilke
Anlæget giver sig tilkjende.

Den Halvidealisme, der kaldes den kritiske, lægger særlig an paa en skarp
Adskillelse mellem de Former, der nødvendig høre Bevidstheden til, og det
Indhold, som skyldes Tingen; Erfaringen kommer istand ved, at et Indhold
fra Tingen fylder Bevidsthedens Former. Det existentielt Concrete er altid
empirisk, men Formerne ere abstracte; det gjælder om at blive sig
Forskjellen mellem det Concrete og Abstracte klart bevidst, og derhos at
godtgjøre, hvilke og hvormange Bevidsthedsformer der uimodsigelig ere
aprioriske.

Følgende Fremgangsmaade kan gjælde for charakteristisk. »Vi opfatte
Rummet concret, idet vi see det opfyldt af en talløs Mængde
Erfaringsgjenstande. Men den concrete Opfattelse er tvetydig; vi maae
spørge: er Rummet som Form for Tingenes Tilværelse givet os i og med
Tingene, eller er det en Bevidsthedsform, der ligger i vor Opfattelse og
aldeles ikke vedkommer Tingen i sig selv? For at faae dette Spørgsmaal
besvaret, maae vi anvende Abstractionen. Mellemrummet mellem A og B
udfyldes f.\ Ex.\ af Gjenstandene a, b og c; tænke vi os nu, at disse
Gjenstande vare borte, saa vilde Rummet mellem A og B blive tomt; men det
vil ikke blive borte, ingen Deel vil gaae tabt: Afstanden fra A til B vil
være den samme som før. Hvis Rummet var en Bestemmelse, der tilhørte
Gjenstanden som Gjenstand, saa maatte jo hver Gjenstand tage sit Rum med,
naar den blev borte; der maatte, saa at sige, komme Huller i Rummet, et
Hul for hver af de forsvundne Gjenstande. Da dette nu umulig kan være
Tilfældet, efterdi Huller i Rum selv ere Rum med indbildte Grændser, saa
følger heraf, at alle Ting i Rummet kunne forsvinde, uden at Rummet
forsvinder: og heraf følger igjen, at det ikke er Tingene selv, men
Bevidstheden, der bærer Rumsformen i sig. Dette kan udtrykkes saaledes: vi
kunne abstrahere fra Tingene i Rummet, men fra Rummet som Rum kunne vi
ikke abstrahere; naar alle Ting ere forsvundne, beholder Bevidstheden sin
egen aprioriske Sandseform, det grændseløse, uendelige Rum tilbage. --- Vi
opfatte Tiden concret, naar vi betragte den i Eenhed med dens Indhold,
d.v.s.\ med de Forandringer og Begivenheder, som foregaae i den; ifølge
denne concrete Opfattelse hedder det: Tiden iler, Tiderne skifte o.s.v.
Men abstrahere vi fra Indholdet, saa staaer Tiden stille. Alt hvad der har
vist sig i Erfaringen, som et Forudgaaende, Daværende og Efterfølgende,
kan tænkes borte; men hvad Erfaringens Mulighed angaaer, kan Tiden som Tid
ikke tænkes borte. At der i den rene, fra alle virkelige eller indbildte
Forandringer adskilte Tid ingen Forandringer kan være, forstaaer sig af
sig selv; men at Tiden heller ikke kan ligge i Tingene eller hæfte ved
Tingene, er ligeledes indlysende. Vi dele jo Tiden i Nutid, Fortid og
Fremtid; skulde Tiden nu tilhøre det i sig selv Værende, saa maatte den
være en uendelig fortsat Nutid; thi Fortiden er forsvunden og svarer til
et Ikke-Værende, Fremtiden, som endnu ikke er kommen, viser ligeledes hen
til et Ikke-Værende. Men hvor lang er vel Nutiden? Den har, naar vi see
nøiere til, ingen Udstrækning; den svinder ind til et Nu, det Flygtigste
af Alt, og kan umulig angaae det i sig selv Værende. Dette Tvetydige, at
Tiden paa een Gang er stillestaaende og forsvindende, beviser, at Tiden
som Tid ikke er Realitet, men en Indbildningsform, en for alle
Forandringer nødvendig Form, en apriorisk Bevidsthedsform. Vi opfatte
Forstanden concret, idet vi opfatte de Ting, vi forstaae, i umiddelbar
Forening med vor Forstaaelse af dem. Det er en Erfaringssætning, at vor
Forstand maa rette sig efter den i Tingene udtalte Forstand: vi faae
Forstand paa Tingene, fordi der er Forstand i Tingene. Er her da to Slags
Forstand, en Forstand i os, som kan feile, og en Forstand i Tingene selv,
som ikke kan feile? Kunne Tingene igjennem Erfaring virkelig antages at
bringe den ufeilbare Forstand ind i vor Bevidsthed?

For at faae disse Spørgsmaal besvarede, maae vi ikke blot see bort fra
Forstandens mangeartede Indhold og ved Abstraction fremhæve de rene
Forstandsformer; men Abstractionen maa endog fordoble sig til Reflexion,
thi Forstandsformerne ere Reflexionsformer: Forstanden virker i Reflecteren
og Tænken. Der er Forstand i Tingen, for saa vidt Tingen har Egenskaber og
ved disse Egenskaber staaer i Forhold til andre Ting; Egenskaberne ere
abstracte Bestemmelser og blive kun forstaaelige derved, at den ene
Abstraction afspeiler sig, d.v.s.\ reflecterer sig i den anden. Lyset f.\
Ex.\ har sine Egenskaber, og de Legemer, som belyses, have, hvert for sig,
deres Egenskaber. Men i de abstracte Bestemmelser for sig er der ingen
Forstand. Det Forstandige viser sig først, idet Egenskaberne reflecteres i
hverandre gjensidig, Lysets i Legemernes og Legemernes i Lysets. Bladet er
grønt; det er en Erfaringssætning. Det Forstandige i denne Erfaringssætning
fremlyser deraf, at Sætningen kan forvandles til en Dom, hvori Tingen er
Subject og Egenskaben Prædicat. Det Samme gjælder om Sætningen: Skyen er
graa, og utallige andre.

Abstrahere vi nu fra Erfaringssætningernes forskjellige Indhold, da
overbevises vi om, at det Forstandige hverken ligger i Tingene, ei heller
i Egenskaberne, hverken i de særegne Subjecter, der kunne være
hvilkesomhelst, ei heller i Prædicaterne, hvorom det Samme gjælder, men
ene og alene i det reflecterede Forhold af Ting og Egenskab, Subject og
Prædicat: alle Forstandsbestemmelser ere Forholdsbestemmelser. Enhver
Forstaaelse af særegne Forhold maa nødvendig forudsætte et oprindeligt
Blik for Forhold. Det er heraf klart, at Tingene i deres mangeartede
Forhold umulig kunne lære os, hvad Forhold er, eller meddele os anden
Forstand end den, vi oprindelig have i os selv. Istedenfor at sige, at det
er Forstanden i Tingene, som ved Erfaring føres over i vor Bevidsthed,
maae vi omvendt sige at det er vor Bevidsthed, der lægger Forstanden ind i
Tingene. Forstandens almeengyldige Reflexionsbestemmelser og Kategorier
høre altsaa udelukkende til det Aprioriske.«

Det Grændseskjel, som den kritiske Halvidealisme paa denne Maade opstiller
mellem apriorisk og empirisk Viden, mellem det, der kommer fra
Bevidstheden, indenfra, og det, der kommer fra Tingene, udenfra, har
imidlertid det Mislige ved sig, at Adskillelsen kun er kritisk. Den
kritiske Adskillelse er altid kjendelig paa en fast, ubevægelig Grændse;
men en fast, ubevægelig Grændselinie mellem det Aprioriske og
Aposterioriske lader sig ikke drage. Dette fremlyser allerede af hvad der
(§ 1) er sagt om Viden og Gjenstand. Slaaer man Viden sammen med det
vidende Individ, da er det sandt nok, at Individet med sin Bevidsthed har
Gjenstandene udenfor sig, Individet er jo selv en udvortes Gjenstand; men
hvo er i Stand til at paavise en fast Grændse mellem Viden og Videns
Gjenstand? Imellem Lande, Provindser, Jordlodder o.\ dsl.\ kan der drages
bestemte Grændser. Disse Grændser kunne til en vis Tid være faste og saa
igjen flyttes; men saaledes forholder det sig ikke med Grændsen mellem
det, som ligger indenfor, og det, som ligger udenfor Bevidsthedens
oprindelige Eiendom. Det er ikke sandt, at man ved at abstrahere fra
Rummets og Tidens Indhold kan beholde det rene Rum og den rene Tid
tilbage; thi naar de sidste Skygger af Ting og Forandringer forsvinde,
forsvinde ogsaa Rum og Tid. Det er ikke sandt, at man ved at abstrahere
fra alle særlige Forhold og Forholdsled kan blive sig de almindelige
Forholdsbestemmelser bevidst som den rene Forstands aprioriske Former.
Forhold uden Forholdsled ere utænkelige; selv i $\frac{0}{0}$ maa der være
et Gjenskin af forsvundne Led, naar der skal være noget Forhold. Disse
Gjenskin, disse for Reflexionen nødvendige Reflexer, ere i sidste Instans
Reflexer af existerende Ting. Indbildningen kan have omdannet dem; men
Indbildningen kan Intet tage af sig selv, den maa hente sit Materiale fra
umiddelbar Erfaring; Abstractionen kan fortynde Reflexerne indtil
Gjennemsigtighed; men Abstractionen maa have Noget, hvorfra den kan
abstrahere.

Det Udvortes maa reflectere sig i et Indvortes, og omvendt; Gjenstandene
maae reflectere sig i Bevidstheden, og Bevidstheden ved Hjælp af
Gjenstandene; det Empiriske maa reflectere sig i det Aprioriske og dette i
hiint. Umiddelbare Tilværelsesgrændser, Existensgrændser, ere
Endelighedsgrændser, og slige Grændser betegne, at A falder paa den ene
Side, B paa den anden Side af Grændsen; men Reflexionsgrændserne have det
Særegne ved sig, at A og B kun udelukke hinanden ved at gjennemtrænge
hinanden i et uendeligt Vexelskin. Sjælen har ikke først, om end paa
ubevidst Maade, sine aprioriske Sandseformer og Forstandsformer, saa at
Indtryk fra Tingene kunne træde til, og Erfaringen begynde. Her er i Tiden
intet Først og intet Sidst. De aprioriske Former maae lige saa oprindelig
beroe paa Indtryk fra Tingene, som disse Indtryk paa Formerne. Det
Aprioriske indskrænker sig ikke til et vist Antal abstracte
Formbestemmelser, som, een Gang bragte til Bevidsthed, vise sig at være
færdige. Nei, af Grundforholdet følger, at der maa være lige saa mange
Ændringer og Forgreninger af de aprioriske Formbestemmelser, som der i
Erfaringsverdenen er Ændringer og Forgreninger af tilsvarende
Gjenstandsbestemmelser. Den voxende Mangfoldighed af erkjendte
Gjenstandsbestemmelser fremkalder efterhaanden en tilsvarende
Mangfoldighed af aprioriske Reflexionsbestemmelser. Al Viden og Videnskab
gaaer ud paa en Apriorisering af det udenfra Givne; thi Viden og Videnskab
tilstræber Indsigt, og al Indsigt bestaaer i en Aprioriseren. Den falske
Apriorisering, der er kjendelig paa, at man deels overseer, deels
forvansker det i Tilværelsen Givne, glemmer de empiriske
Gjenstandsbestemmelsers afgjørende Betydning for den bestemmende
Reflexion, miskjender det Aprioriske og maa desaarsag nøies med en
Indbildning istedenfor sand Indsigt.

Hvad Physiologien viser, at der i Aartusinders Løb kan spores en
fremadskridende Udvikling af Hjernen og Hjernevirksomheden, finder sin
Stadfæstelse i Videnskabernes Historie. Bevidsthed om det Aprioriske maa
erhverves lige saa vel som Bevidsthed om det Empiriske. En eensidig
Dyrkelse af aprioriske Videnskaber kan føre til Misforstaaelse af det
Aprioriske, og en eensidig Dyrkelse af empiriske Videnskaber til
Misforstaaelse af det Empiriske; thi Udviklingen beroer paa Vexelvirkning.
Sand Indsigt i det Empiriske er betinget af en tilsvarende Indsigt i det
Aprioriske, og omvendt.

Sand Viden er Indsigt; fremadskridende Indsigt beroer paa en stadig
fremadskridende Apriorisering: Viden er i sit Væsen apriorisk.


\chapter{Subjectiv og objectiv Viden}

\section{Objectiveringsproblemet}
\label{sec:5}
% pp. 40--49

\noindent
For saa vidt Vidensindholdet ligger i den Videndes eget Indre, er
Viden subjectiv, for saa vidt Indholdet hentes fra Tingenes Verden,
er Viden objectiv. Der kan dog ikke paa denne Forudsætning opstilles
nogen kritisk, d.v.s. fast, værende Grændse mellem subjectiv og
objectiv Viden. Al subjectiv Viden maa nødvendig have et Element af
det Objective, thi en Viden uden Gjenstand, en absolut gjenstandsløs
Viden, er lige saa meningsløs som et Cirkelcentrum uden Peripherie.
Ikke desto mindre er der en væsenlig Forskjel mellem ydre og indre
Objecter, og hvad de sidste angaaer, en væsenlig Forskjel imellem det
Subjective, der udelukkende tilhører den enkelte Individualitet, og
det for alle vidende Individualiteter Fælles, det ideelt Objective,
der er eenstydigt med det Almeengyldige. Viden af en Steen er udvortes
objectiv, Viden af Faldloven er væsenlig objectiv; men Individets Viden
af sit eget Befindende, sin Attraa og Længsel i ethvert givet Øieblik,
er væsenlig subjectiv, skjøndt Længsel og Attraa ere Videns ideelle
Objecter.

Grundbetingelsen for al Viden er, at det vidende Selv maa opfatte og
bestemme sin Gjenstand, hvad enten denne Gjenstand uvilkaarlig og med
blind Nødvendighed paatrænger sig udenfra, eller vilkaarlig fremhæves
indenfra. Først idet Videns Gjenstand opfattes og bestemmes af den
Vidende, d.v.s.\ objectiveres, bliver den et Videns Object. Intet
Object uden tilsvarende Objectivering; det er en apriorisk Lov, der
bærer al Empiri, en Grundlov, der i Videnskaben er, om muligt, endnu
mere urokkelig end Newtons Tyngdelov. Heraf indsees, at en kritisk
Grændse, en Grændselinie, paa hvis ene Side man fastholder den
objectiverende Subjectivitet for sig, medens Objectet bliver staaende
paa den anden Side, er forvirrende og meningsløs. Modsætningen af
Subjectivt og Objectivt maa bestemmes saaledes, at der i Objectet paa
hvert Punkt er et Moment af det Subjective, og omvendt: Objectiveringen
selv er Eenheden af det Subjective og det Objective, men en Eenhed af
modsatte Led, en Modsætningseenhed.

Først naar vi opfatte Modsætningen af subjectiv og objectiv Viden
\textit{in concreto}, er der Grund til at spørge om en kritisk Grændse.
Det Vidensindhold, der udelukkende hentes fra den vidende Individualitets
eiendommelige Indre, kan umulig antage nogen objectiv, i sig almeengyldig
Form eller begrundes videnskabelig. Hvad Subjectiviteten indenfor denne
Synskreds objectiverer, angaaer alene dens egen specifiske Væren, dens
Tanker, Følelser, indre Oplevelser. Om end Subjectet nok saa
uforbeholdent taler derom til Andre, taber det derved ikke sin
subjective Charakteer. Det saaledes Objectiverede kan høre til
naturlige Egenheder, Mangler, Fortrin, exceptionelle Særheder o.dsl.,
og det kan omfatte det for Subjectiviteten Allervæsenligste, dens
personlige Livsanskuelse, dens individuelle Opfattelse af sin
oprindelige Bestemmelse, sit sidste Maal og Betingelserne for dets
Opnaaelse.

Den objectivt Vidende maa være vidende om, at den menneskelige
Subjectivitet, at dømme efter objective Erfaringer, kan tillægge sig
ulige Værd, kan fatte sig blot som timelig, kan sætte sig selv som
evig og i det Hele opfatte sin Bestemmelse paa forskjellig Maade; men
han maa ogsaa vide, at Grunden til den forskjellige Selvvurdering ikke
kan hentes fra noget i sig Objektivt, men oprindelig ligger i
Subjectiviteten selv. Dersom den Paagjældende endnu ikke er klar over
sig selv og orienteret i sin egen Bevidsthed, er han idelig udsat for
at overskride Grændsen, jage efter objective Grunde, søge Visheden om
sig selv udenfor sig selv og forvexle den subjective Viden med den
objective.

Her kunne flere Tilfælde vise sig. Er Individualiteten svag, og har
Uklarheden sin Grund i Svaghed, saa faaer han ikke Ro, før han har
fundet en Autoritet, der kan sige ham, hvad han egenlig er og hvortil
han er bestemt; denne Autoritet er saa det Objective: nu mener han at
have en objectiv Viden om sin høieste Bestemmelse. Er Individualiteten
stærk, og har Uklarheden sin Rod i ubetvungen Egoisme, da finder han
ikke Fred, før han har faaet andre Individer sin individuelle Anskuelse
paatvungen, jo flere Individer, desto bedre: de Manges Underkastelse
er ham et Beviis paa, at han har forstaaet at gjøre sin subjective
Viden objectiv. Er Individualiteten dybsindig, d.v.s.\ saa aldeles
dybsindig, at det principielt Ueensartede løber i Eet for hans Syn,
da vil han overskride Grændsen med stor Anstand. Han vil blive sig
objectiv i følgende Indbildning: »Hvad der er Sandhed for mig
personligt, maa nødvendigviis være objectiv Sandhed, Sandhed for Alle,
thi mit Væsen og Menneskehedens Væsen er i Grunden det samme Væsen.
Naar desuagtet ikke Alle antage, hvad jeg antager, da er det enten af
Mangel paa god Villie eller af Mangel paa Indsigt i det Objective, thi
min subjective Viden er objectiv: jeg veed personlig den objective
Sandhed.« --- Betragte vi disse tre Individualiteter, den svage, den
egoistiske og den selsomt dybsindige, arbeidende sammen i Historien,
saa forstaae vi, hvilken Strid, hvilket Had, hvilken Forfølgelse og
Underkuelse det koster, at faae en saadan subjectiv Viden tilgavns sat
over i den objective.

Fra en modsat Yderlighed finde vi den objective Viden villig til, i
Videnskabens Navn, at tage den subjective under sit Formynderskab.
Først, hedder det, maa Videnskaben have den menneskelige Natur tilbunds
undersøgt, før Individet tør fælde nogen Dom om det Personliges egenlige
Værd; først naar man ved physiologiske Forskninger har udfundet, hvad
det vi kalde Sjæl egenlig er, kan den Enkelte faae at vide, om han
virkelig har en Sjæl eller ikke. At det ikke er Videnskaben, men en
Misforstaaelse af, hvad Videnskab er, som fra denne Side frister til at
overskride Grændsen, ligger i Sagens Natur. Vistnok er den
videnskabelige Bevidsthed vendt mod det Objective, selv naar der handles
om Subjectiviteten, thi kun det Almeengyldige har Gyldighed i
Videnskaben; men derfor er den dog langt fra at frakjende det Subjective
i Selvets Eiendommelighed dets oprindelige Ret og Betydning. For at
faae det Objective videnskabeligt belyst, maa den Vidende have
Objectivitetens Sol i Ryggen.

Hvorledes det Objective tager Magten fra den objectiverende
Subjectivitet, afklæder den alle individuelle Egenheder og nøder den
til at underkaste sig almeengyldige Bestemmelser og Love, sees
allertydeligt af saadanne Exempler, hvor det Objectiverede ikke afhænger
af udvortes Erfaring, men har sine Forudsætninger i Bevidstheden selv.
Vi forudsætte, at alle Mennesker ere fiirkantede; vi forudsætte, at
det omspurgte Individ, Cajus, er et Menneske. Det er vilkaarlige
Forudsætninger, og vi gjorde vel bedst i ikke at antage dem; men have
vi een Gang antaget dem, saa ere vi bundne ved Antagelsen og kunne
ikke unddrage os den Conseqvens, at Cajus maa være fiirkantet. Vi
sætte et Forhold som $a : b = c : x$, vi sætte det ganske vilkaarlig;
men, bundne til Forudsætningen, blive vi nødte til at antage
$x = \frac{bc}{a}$. Af subjective Forudsætninger fremgaae objective
Conseqvenser, og af objective Conseqvenser objectiv Viden. Vistnok er
denne Viden endnu kun formel; men den formelle Objectivering er dog
intellectuel, og intellectuel Objectivering en ufravigelig Betingelse
for al videnskabelig Viden. Vilde Nogen herved anmærke, at et subjectivt
Indhold jo ogsaa kan objectiveres intellectuelt, at det saaledes
Objectiverede maa antage videnskabelig Charakteer, saa sandt det er
den intellectuelle Objectivering, der giver Viden sit objective Præg:
da er det ikke vanskeligt at forvisse sig om, at Grændsen mellem
subjectiv og objectiv Viden endnu staaer fast.

Reel Viden, Viden af et positivt Indhold, udkræver ikke blot en formel,
men ogsaa en reel Objectivering. Et Par Exempler kan oplyse dette.
Dersom det menneskelige Selv er anlagt paa Udødelighed, da maa ethvert
Menneske berede sig paa at blive dømt efter Evighedens Love. ---
Dersom alle Legemer ere tunge, da maa ethvert Legeme, som, hævet op
over Jordens Overflade, berøves sin Understøtning, falde. --- I formel
Henseende ere begge Sætninger lige antagelige; men i Forhold til Viden
ere de uligeartede. Forsætningen i den første er en
Subjectivitetssætning, hvis almene Gyldighed vakler, fordi den opfattes
forskjellig af de forskjellige Individualiteter; dens Mangel paa
Objectivitet viser sig i en almindelig Uvished, en bestandig Svæven
mellem Realitet og Indbildning. Den Forsætning derimod er en
Objectivitetssætning. Forstanden, der støtter sig paa udvortes
Erfaring, kan gjøres indlysende for Alle, som have almindelig Sands og
Forstand. Subjectiviteternes indre Eiendommelighed kommer herved ikke
i Betragtning.

Da en Objectivering af udvortes Gjenstande lige saa vel maa forudsætte
Subjectivitetens Medvirkning, som en Objectivering af indvortes, saa
følger unegtelig heraf, at der i den første lige saa vel maa kunne
vise sig subjective Forskjelligheder, som i den sidste. Den Kyndige
objectiverer sig Tingene anderledes end den Ukyndige. Astronomen
objectiverer sig Verdensbyggingen paa en ganske anden Maade end den,
der uden Kjendskab til Astronomi, umiddelbart betragter Stjernehimlen.
Mineralogen seer i Mineralet et andet Object end den uvidende Liebhaver,
der samler paa sjældne Stene. Botanikerens Objectivering af en Plante
er væsenlig forskjellig fra Malerens og Digterens. Vi indrømme det.
Men disse Forskjelligheder udslette ikke Grændsen mellem subjectiv og
objectiv Viden; tvertimod, de tjene kun til at skjærpe den og gjøre
den kjendelig.

Den subjective Viden viser os hen til Forskjelligheder i
Subjectiviteten, som kun tilnærmelsesviis kunne hæves paa subjectiv
Maade, ved Optugtelse, Vækkelse, Formaning o.s.v. Den objective Viden
tager derimod kun Hensyn til saadanne Subjectivitetsforskjelligheder,
som den selv formaaer at ophæve eller dog udligne paa objectiv Maade,
ved nemlig at gjøre den Ukyndige kyndig, den Uvidende vidende.
Muligheden af en saadan Udligning beroer paa, at Individernes
Bevidsthed, hvor afvigende de foreløbige Synsmaader end kunne være,
indeholde Grundbetingelser, elementære Forudsætninger, der ere fælles
for Alle, nemlig Sandsning og Tænkning. Den Uvidende kan see, at Solen
bevæger sig, medens Jorden staaer stille, og det er netop en Betingelse
for, at han kan komme til at indsee, at det er Jorden, der bevæger sig,
medens Solen staaer stille. Den Ukyndige kan see, at Bladet er grønt,
og dette er just en Betingelse for, at han kan lære at indsee, at det
egenlig ikke er Bladet, men dets Chlorophyl, der er grønt.

Den subjective Viden om existerende Ting giver os kun en Mening; og
Meningerne ere forskjellige, ja stridende mod hverandre, efterdi de
beroe paa overfladisk Objectivering. Den objective Viden ophæver
Meningernes Forskjellighed, idet den giver Indsigt, og Indsigten
betinges af væsenlig, grundig Objectivering.

Indbildningen har imidlertid sin Andeel i den objective Viden lige saa
vel som i den subjective, og Indflydelsen heraf er i Virkeligheden af
saa betænkelig Art, at den truer med at berøve os Objectiviteten selv
og reducere al vor Viden til en blot subjectiv Viden. Objectet er,
ifølge Grundloven, det Objectiverede; for saa vidt bliver det jo
vistnok meningsløst at ville drage en kritisk Grændse mellem Objectet
selv og vor Opfattelse deraf. »Er f.Ex.\ dette Bord, denne Lampe i sig
selv et saadant Object som det, jeg kan objectivere?« Hvilket urimeligt
Spørgsmaal. Hvad er det for et Object, Du taler om? »Dette Bord!« Er
det det Bord, Du objectiverer, eller et andet, som Du ikke objectiverer?
»Ved Objectet forstaaer jeg naturligviis dette objectiverede Bord.« Og
saa kan Du spørge, om det Du objectiverer ogsaa virkelig er det, Du
objectiverer.

Men Objectiveringen indeslutter to uligeartede Momenter, det subjective
og det objective, og naar ethvert af disse fremhæves for sig, kan der
ganske vist være Spørgsmaal om, hvorvidt Objectet og det Objectiverede
i Et og Alt ere congruente. Vi kunne prøve det ved afvexlende at sætte
det Variable og Constante i Objectiveringens to Factorer.

Sætte vi det Constante i Subjectiviteten og lade det Variable falde i
det Objective, saa ville Objectiveringens forskjellige Resultater være
afhængige af det i Objecterne Objective: Resultaterne variere, fordi
Objecterne variere. Bordet viser sig fiirkantet, naar det er fiirkantet,
og rundt, naar det er rundt. Indbildningen gjør dette aabenbart ved at
lade hvert Object staae med en saadan Bestemthed og Selvstændighed
udenfor Bevidstheden, at der slet ingen Objectivering behøves: Stenen
er Steen og Bladet Blad, selv om Ingen sandser det; Solen skinner,
enten Nogen seer det eller ikke.

Sætte vi det Constante i Objectiviteten og lade det Variable falde i
det Subjective, da ville Objectiveringens forskjellige Resultater gjøre
enhver objectiv Bestemmelse ukjendelig. Indbildningen gjør det
anskueligt ved at lade det Objectiverede have Objectet bagved sig i
det Uendelige. Bagved den objectiverede Lampe staaer den egenlige
Lampe, Objectet, som det er i sig; hvis en anden Subjectivitet, en
Subjectivitet med fuldkomnere Betingelser, objectiverede Lampen, da
vilde det blive en anden Lampe. De objectiverede Lamper vilde være
forskjellige; men Lampen i sig, det egenlige Object, staaer altid
bagved; den objectiverede Lampe varierer med hver varierende
Objectivering, og Lampen i sig, det constante Object, viger bort i det
Uendelige.

Resultatet af disse Undersøgelser fører dog ikke til Brud paa
Grundloven, men til en nærmere Forstaaelse af de forskjellige Vendinger,
under hvilke Indbildningen kan skuffe den objectiverende Bevidsthed.
Idet der opstaaer Tvivl om Objectiveringens Gyldighed, har Indbildningen
strax to Objecter til Sammenligning, nemlig det objectiverede og det
ikke objectiverede. Den objectiverede Lampe, det Object, der svarer til
Ordet, har Lampen i sig, det Unævnelige, der ikke svarer til Ordet,
bagved sig. Hvorledes komme vi da til Bevidsthed om det ikke
objectiverede Object? Indbildningen har naturligviis selv objectiveret
det, men det skal være en stor Hemmelighed. Skal Indbildningen rykke
ud med Hemmeligheden, viser den sig strax som en Skuffelse. Her skulde
skjelnes mellem imaginære og reelle Objecter; men Adskillelsen lader
sig ikke gjennemføre. Først hedder det, at de objectiverede Objecter,
Bordet og Lampen, ere imaginære, og at de unævnelige Objecter, Bordet
i sig, Lampen i sig, paa hvilke Ordene kun passe som imaginære
Betegnelser, ere reelle. Men naar man saa besinder sig paa, at de
unævnelige Objecter, som Indbildningen paa eget Ansvar tilkjender
Objectivitet, netop ere de indbildte Objecter, altsaa de imaginære:
saa fristes vi jo til at antage de objectiverede for de reelle. Her er
en Uvished, som maa hæves.

\section{Erkjendelsesproblemet}
\label{sec:6}
% pp. 50--59

\noindent Objectiveringsloven siger, at kun det er Object, som bliver
objectiveret; men heraf følger endnu ingenlunde, at vor menneskelige
Objectivering er den eneste, der oprindelig bestemmer det Objective, eller
at Objecterne i Et og Alt ere saaledes, som vi formaae at objectivere dem.
Det Spørgsmaal, hvor vidt vore Evner række, naar Talen er om en sand, en
grundig, omfattende og udtømmende Objectivering af Existensbetingelser og
existerende Ting, er Erkjendelsens Grundspørgsmaal. I almindelig Forstand
kan Spørgsmaalet besvares saaledes: vi erkjende, hvad vi objectivere, i
Forhold til, som det Objectiverede kan blive os gjennemsigtigt; men det
bliver os kun gjennemsigtigt i Forhold til, som vor Bevidsthed formaaer at
frembringe det. Begrændsede Anskuelsesformer, Objecter af intellectuel
Objectivering, saasom: Linier, Trekanter, Cirkler, Tetraedre, Hexaedre
o.s.v.\ erkjende vi fuldt ud, for saa vidt vi ere i Stand til at frembringe
dem i Anskuelsens Form og ved tilsvarende Ligninger at henføre dem under
visse Love. At spørge, om der bag ved den objectiverede Cirkel ikke skulde
staae et unævneligt Object, en Cirkel i sig, som vi ikke kunne erkjende,
er ørkesløs Dybsindighed.

Anderledes forholder det sig med de Objecter, der i sig have Noget, vi
ikke kunne frembringe, et Element, der paanøder sig Anskuelsen og bestemmer
Forestillingen, men bliver staaende for Bevidstheden som noget Udvortes.
Dette er Tilfældet med Materien og det Materielle. Vi kjende Materien, vi
anerkjende dens Realitet; men dens Væsen, det, den er i sig selv, kunne vi
ikke erkjende. Vor Tænkning kan ved Anskuelsens Hjælp frembringe et
Cirkelelement, et Hexaederelement, men et Masseelement, en Moleküle, et
Atom formaae vi ikke at frembringe. Derfor tager det Objective i Materien
sig ud som noget Mørkt, Gaadefuldt, i sig Uerkjendeligt.

Det i Materien Materielle er sandseligt, men ikke tænkeligt, om det end
bestemmes ved Tænkning. Undersøge vi Sandsningens Betydning for
Objectiveringen og derved for Erkjendelsen, saa henvises vi til Sandsernes
Physiologi. Vi lære, hvorledes Sandseorganerne ere byggede; vi lære, hvor
forskjelligartet den Virksomhed er, som disse Organer udøve, hvor
forskjellige de ere saavel i Antal som i Form og Udvikling hos de lavere
og høiere Dyrearter, ja at et enkelt Organ hos Mennesket, Øiet f.\ Ex.,
kan være saa eensidig anlagt, at det, som hos de Farveblinde, gjør en
Undtagelse fra den almindelige Regel. Idet vi nu indsee, hvor afhængig den
sandselige Objectivering er af Sandsernes eiendommelige Virksomhed, hvor
afhængig denne Virksomhed er af Sandseorganernes Bygning, og hvor afhængig
vor Erkjendelse er af Sandsernes Vidnesbyrd, saa er det intet Under, at
der opstaaer Tvivl, om vi ogsaa formaae at erkjende Tingene som de i sig
selv ere. At sætte Aarsagen i Tingen, Virkningen i vor Erkjenden og
beraabe sig paa Ligheden mellem Aarsagen og Virkningen er yderst misligt.
Hvad Lighed er der vel imellem en Æthersvingning og en Lysfornemmelse,
imellem en Luftbølge og en Lyd? Hvilke Omdannelser maae de fra Tingene
kommende Paavirkninger ikke undergaae, inden de gjennem Hjernen naae til
Bevidstheden?

Antage vi nu, at der paa andre Kloder i Universet kan gives Væsner med
Sandser af anden og høiere Art end vore, da ville disse Væsner jo
objectivere sig Tingene paa en langt fuldkomnere Maade end vi, og i det
Materielle maaskee erkjende et Object, hvorom vi ikke have nogen Anelse.
Ikke nok hermed. Naar vi ei engang kunne stole paa vore Sandseformers
objective Gyldighed --- thi hvad forvisser os om at Rummet netop har tre
og kun tre Dimensioner? --- saa kunne vi jo heller ikke fuldtud stole paa
vore Forstandskategorier. Vistnok ere vi i al vor Erkjenden nødte til at
skjelne og adskille Eenhed og Fleerhed, Forudsætning og Conseqvens, Hele
og Dele, Aarsag og Virkning; men hvo kan vide, om en saadan Adskillelse
har Gyldighed for den fuldkomne Forstand? Det er jo muligt, at der gives en
intellectuel Anskuelse, som ved at see uendelig mange Ting i Rummet paa
een Gang, er ophøiet over Rummet, og ved at see utallige Bevægelser i eet
Syn, er hævet over Tiden; det er jo muligt, at en Forstand med en saadan
Anskuelse ikke erkjender den Skjelnen, Adskillen og Udstykning, der er
uundværlig for vor Erkjendelse, men seer Eenheden i Fleerheden,
Conseqvensen i Forudsætningen, Delene i det Hele, Virkningen i Aarsagen;
altsaa: det er jo muligt, at der gives en Forstand, saa qvalitativt
forskjellig fra den menneskelige, at vi end ikke kunne danne os nogen
Forestilling, end sige noget Begreb om den.

Men hvad følger saa heraf? At vor saakaldte objective Erkjendelse kun er en
Erkjendelse af Phænomener og tilsvarende Love, men ingenlunde nogen
Erkjendelse af det Objective selv, af Tingenes Væsen. Det er en vulgær
Mening, at mange Kundskaber og dybtgaaende Forskninger skulde lede til saa
høie Forestillinger om Erkjendelsesevnens Ufeilbarhed, at Forskeren let
kan fristes til Hovmod. Dette er ingenlunde Tilfældet. Vi finde i Reglen,
at de meest udmærkede Forskere, de, der have beriget Videnskaben med de
største Opdagelser, tale i et meget beskedent Sprog om vor Erkjendelses
Rækkeevne: »hvad vi vide er for Intet at regne imod, hvad vi ikke vide.«

Denne Indrømmelse bliver imidlertid misforstaaet, naar den
Uvidenskabelige, der er tilfreds med et praktisk Skjøn, vil belære
Forskeren og føre ham paa rette Vei. »Hvorfor gjøre et saadant Væsen af
Videnskaben, naar den dog, efter Din egen Tilstaaelse, ikke formaaer at
lede os til Erkjendelse af Tingenes Væsen? Du staaer jo med Dine
Undersøgelsers Rigdomme, som en fattig Mand paa de objective Erkjendelsers
golde Hede; hver Gang Du mener at holde Objectet fast, forvandler det sig
til et Luftsyn, et Phænomen af Noget, Du ikke kjender. Du maaler Styrken af
en elektrisk Strøm, men veed ikke, hvad Elektricitet er; Du iagttager
Legemets Bygning, men veed ikke, hvad Livet er, eller hvorfra det kommer;
Du griber efter Erkjendelsens Frugter og vil slukke Din Tørst ved
Sandhedens Kilder, men Frugterne undvige Dig, Kilderne skuffe Dig: Du har
det jo næsten som en Tantalus i sin Underverden. Nei, da er det Subjective
anderledes vederkvægende. Betragt Tilværelsen i Frihedens Speil; brug
Livets Lys, lad Underet hjælpe Dig, og bind Dig ikke ved smaalig Beregning
til blinde Love.« Slige velmeente Raad kunne dog ikke hjælpe den redelige
Forsker. Han veed, at hvis han blot i eet eneste Punkt giver Slip paa
Lovenes Traad, er al hans Forskning forgjæves.

Videnskabelige Opgaver kræve videnskabelig Løsning; de, der mene at kunne
hjælpe paa Videnskaben ved at indføre Erkjendelser af overvidenskabelig
eller uvidenskabelig Art, forraade kun, at de aldrig have fattet et
videnskabeligt Problem og neppe ane, hvad Videnskab væsenlig er.
Videnskaben maa hjælpe sig selv.

De anførte Vanskeligheder ere da heller ikke saa afskrækkende, som en
Erkjendelseslære af eensidig Kritik vil udgive dem for. Vistnok er det,
hvad Sandserne og Sandseorganerne angaaer, unegteligt, at eet og samme
Substrat maa ved at objectiveres i forskjellige Systemer vise sig som
forskjellige Objecter; men dette stadfæster jo netop, hvad Grundloven
siger. Betegnes Substratet ved $X$, de deri liggende Bestemmelser ved
$X_0$, $X_1$, $X_2$, saa ville disse Bestemmelser, naar de sandses i
$A$-Systemet, give $a_0$, $a_1$, $a_2$, \ldots\ og det hele tilsvarende
Object være $O_a$; i $B$-Systemet vil det være $O_b$; i $C$-Systemet $O_c$
o.\ s.\ frmdl. Men er dette et Beviis paa, at det Objective ikke kan
erkjendes? Ingenlunde; det er just et Beviis paa Vexelvirkning mellem
Objectivt og Subjectivt. Kunde samme Substrat ved at objectiveres i
forskjellige Systemer derimod give Anledning til Dannelse af samme Object,
$O_a = O_b = O_c$, da først vilde det Objectives Medvirken være forvandlet
til en Indbildning. At forskjellige Objecter maae danne sig, naar
Substratet objectiveres i forskjellige Systemer, er jo lige saa
forstaaeligt, som at forskjellige Legemer maae danne sig, naar Kulstoffet
kommer i Forbindelse med Ilt, med Brint, med Kvælstof o.s.v. En
Subjectivitet, som var i Stand til at gjennemskue de forskjellige
Sandsesystemers forskjellige Virksomheder, vilde ingenlunde enten erklære
visse Objecter, $O_b$ og $O_c$, for falske og kun $O_a$ for sandt, eller
alle Objecterne for lige falske; men meget mere indskjærpe, at hvert
Object i sit System har fuld Gyldighed paa Grund af Vexelvirkningen. At
Objectiveringen i det ene System kan staae lavere end i det andet, gjør
for saa vidt kun Udslag, som Virkningen altid maa svare til Aarsagen.

Man har gjort gjældende, at Erkjendelsen umulig kan naae det Objective,
eftersom Aarsagen i Objectet maa være Virkningen i det Subjective høist
ulig. Denne Vanskelighed er mere fremkunstlet end virkelig. Ifølge
Vexelvirkningen kan Aarsagen ikke udelukkende tilhøre det Objective og
Virkningen det Subjective, men idet Betingelserne tages med, er der Aarsag
paa begge Sider, i det Subjective saa vel som i det Objective, og Virkning
fra begge Sider; der kan da, ifølge Vexelvirkningens Lov, ikke være Tvivl
om, at Virkningen svarer til Aarsagen. Kun Abstractionen i Forening med
Indbildningen kan fremkogle det Skin, at der er objective Bestemmelser for
sig uden alt Hensyn til det Subjective, og subjective Bestemmelser uden
Relation til det Objective; men man behøver blot at gjøre et Forsøg paa at
gjennemføre en saadan Adskillelse, og man vil i Kjendsgjerningen opdage,
hvor intetsigende den er. Naar hver objectiv Bestemmelse er en Bestemmelse
af Vexelvirkning, og hver subjectiv Bestemmelse ligesaa, hvor kan man saa
tvivle om, at den af Vexelvirkningen resulterende Erkjendelse maa give en
virkelig Eenhed, en væsenlig Lighed, af det Objective selv og det
Objectiverede.

Ja, indrømmer man, en Lighed er her jo vistnok; men Ligheden angaaer ikke
Objectets Væsen, men kun dets Phænomen. --- Dog, hvad vil dette egenlig
sige? At Indbildningen igjen driver Spil med Abstractionen! Man vil dog
vel ikke sætte Tingsobjecter i Række med Hallucinationsobjecter? Har
Phænomenet slet Intet med Væsnet at gjøre, saa er det jo aldeles
betydningsløst; saa kan Lampens Væsen maaskee være skjult under Phænomenet
Elephant, og Elephantens Væsen under Phænomenet Blækhuus. Hvad hjælper
det, at vi erkjende Tingenes Phænomener, naar Phænomenerne slet ikke
vedkomme Væsnet? Skal et bestemt Phænomen staae i et nødvendigt Forhold
til et bestemt Væsen, som Skin til Realitet, saa maae Realitet og Skin
nøiagtig svare til hinanden; det Skin, der svarer til Realiteten, er et
reelt Skin, og det Skin, hvortil ingen Realitet svarer, et falsk Skin.
Erkjendelsen maa altsaa gaae ud paa at skjelne mellem det falske og det
reelle Skin. Men paa det reelle Skin have vi kun eet Mærke, dette nemlig,
at det svarer til en bestemt Realitet. Er der nu nogen Kritik, der kan
paavise en reel Forskjel mellem Tingens Realitet og dens Væsen, mellem det
reelle Skin og det sanddrue Phænomen? Det reelle Skin er den bestemte
Tilsyneladelse af Væsnets Realitet, og Realiteten en Borgen for, at det
reelle Skin ikke bedrager. Videnskaben kan altsaa trøste sig med, at i
samme Forhold som den formaaer at erkjende Phænomenerne i deres
Sammenhæng, vil den ogsaa være i Stand til at erkjende de væsenlige
Bestemmelser og i disse have et relativt Udtryk for Tingenes Væsen.

I Undersøgelserne over Rum og Tid, over Forstanden og dens Kategorier
drives Abstractionerne undertiden til en saadan Yderlighed, at Illusionen
aabenbart har Kritiken til Bedste. Tiden slipper endda nogenlunde derfra;
men Rummet er idelig udsat for den kritiske Inqvisitions Pinebænk. At
Rummet netop har tre Dimensioner, hverken færre eller flere, er yderst
tvivlsomt; thi der skal jo være sandsende Væsner, hvis Opfattelse
indskrænker sig til een, høist til to Dimensioner, og der er i Videnskaben
allerede optraadt Erkjendende, som ikke kunne nøies med mindre end fire
Dimensioner. Det seer betænkeligt ud. Dersom der virkelig kan sandses i et
Rum med een eneste Dimension, hvorfor skulde man saa ikke ved behørig
Oekonomi kunne bringe det til at sandse i et Rum med slet ingen Dimension?
Er man allerede naaet til at opdage den fjerde Dimension, saa er det jo
ikke urimeligt, at man med Tiden vil kunne opdage en femte, en sjette,
maaske en ellevte Dimension, og hvad er saa Rummet?

Endnu betænkeligere seer det ud med Forstandskategorierne. At det ene
Menneskes Forstand kan overgaae det andets, er en bekjendt Sag; men herved
maa dog udtrykkelig fremhæves, at, endskjøndt den underlegne Forstand ikke
forstaaer den overlegne, saa er denne dog altid i Stand til at forstaae
hiin og gjøre det tydeligt, at Lovene og Kategorierne ere de selvsamme.
Analyser, Begrebsdannelser og Resultater kunne paa Videnskabens høiere
Trin være saa utilgængelige for dem, der endnu kun kjende de første
Begyndelsesgrunde, at det næsten seer ud, som om Forstanden selv ved hvert
Skridt opad forandrede sin Natur. Hvad den Kyndige seer med et eneste
Blik, kan volde den Ukyndige utroligt Besvær, inden han faaer det
udstykket, sammenstykket og forstaaet. Men naar man af en saadan Differens
lader sig forlede til at antage en Forstand, der er saa ophøiet, at de
Kategorier, Forstanden i os er nødt til at anvende, tabe deres Betydning,
saa befinder man sig midt i det Meningsløse. Vi kunne tænke os vore
Kategorier integrerede ved Mellembestemmelser, som vi endnu ikke have
bragt til Bevidsthed; vi kunne tænke os en Alsidighed i Overblik, en
Skarphed i Skjelnen, en Hastighed i Sammenfatning, en Samtidighed i
Heelhedserkjendelse, som uendelig overgaaer enhver menneskelig
Forstandsudvikling; men vi maae bestandig forudsætte, at Kategorier og
Grundlove oprindelig ere de samme.

Hvad Mening er der vel i den Antagelse, at en fuldkommen Forstand slet ikke
har Brug for vore Kategorier? Her er ingen videnskabelig Erkjendelse, men,
reent ud sagt, et Indbildningens \textit{Salto mortale}, der underholder
Kritiken. For at indsee, hvad en saadan fuldkommen Forstand har at betyde,
maae vi ikke blot tænke os Grændserne for vore Forstandsevner uendelig
udvidede; nei, vi maae gaae anderledes til Værks, vi maae tænke os vor
egen Forstand afskaffet og en heel ny Forstand, en Forstand med et andet
Væsen og andre Love sat i Stedet. Hvad hjælper os da til en saadan
Ombytning? Indbildningen! Thi kun Indbildningen kan falde paa et Raad som
dette: »Opgiv den Forstand, Du har, da først erkjender Du den Forstand, Du
ikke har.« Men Indbildningskritik er en upaalidelig Veileder; for at gribe
den Forstand, vi ikke have, maae vi slippe den, vi have; vi skulle altsaa
erkjende den fuldkomne Forstand ved at gaae fra Forstanden: det er meget
uvidenskabeligt. Kritiken forsikkrer, at det ikke er saaledes meent, og vi
troe det gjerne, thi Indbildningen tilslører Meningen; men naar Sløret
løftes, hvad saa? Kun ved at fastholde Objectiveringens Grundlov og
frigjøre sig fra Illusionerne, de kritiske saa vel som de speculative, vil
Videnskaben efterhaanden vide at bestemme sin væsenlige Grændse og
arbeide sig frem til Erkjendelsesproblemets Løsning.

\section{Fuldkommen og ufuldkommen Viden}
\label{sec:7}
% pp. 60--69

\noindent Alt som Erkjendelsernes Kredse udvide sig, bliver vor Viden
fuldkomnere; det kunde da synes, som maatte den fuldkomne Viden, der
omfatter Alt og gjennemskuer Alt, overskride alle Grændser, og være
aldeles ubegrændset. Det kan dog ikke forholde sig saaledes. Ingen Viden
er uden Gjenstand, ingen Objectivering uden Object, ingen Erkjenden uden
et Noget, som bliver erkjendt. Gjenstand, Object, det Noget, som
erkjendes, ere forskjellige Udtryk for Grændsebestemmelser, idet de
betegne det i Viden, som ikke er Viden, men hvoraf Viden bestemmes. En
Viden uden Modsætning, uden Grændse, en Viden af Viden i det Uendelige er
som en Cirkel med en uendelig Radius. Forskjellen imellem den fuldkomne og
ufuldkomne Viden beroer paa en væsenlig Forskjel mellem iboende Grændser
og Skranker.

En Viden, der har Skranker, er indskrænket, og indskrænket er al den
Viden, der betinges af Sandsning. Enhver sandselig Gjenstand er for den
sandsende Bevidsthed en Skranke, Maanen og Stjernen saa vel som Muren og
Døren. Det er en Sandsningens Grundlov, at Bevidstheden enten maa undvære
Realiteterne eller finde sig i, at de stille sig som Skranker; dens
Ufuldkommenhed viser sig just deri, at den ikke selv kan frembringe dem.
Gives der altsaa en fuldkommen Viden, saa kan den ikke være sandsende; den
maa da have Realiteterne paa en ganske særegen Maade, nemlig derved, at den
oprindelig frembringer dem. Vi støde her paa en gjennemgaaende Forskjel
mellem sandsende og reelt frembringende Objectivering. Den sandsende
Objectivering er en Objectivering i Afmagt, svarende til en afmægtig
Viden; thi Passiviteten i den Sandsende er netop et Afmagtens Udtryk. Den
reelt frembringende Objectivering er en Objectivering i Almagt;
Frembringelsen selv, den skabende Virksomhed, er Almagtens Udtryk. At en
skabende Virksomhed ikke kan gjennemskues af sandsende Væsner, om de end
staae tusinde Trin over Mennesket, følger af sig selv. Men ligger der i
Begrebet om en skabende Frembringelse ikke en besnærende Modsigelse, en
Slange, som Indbildningen har skjult i Græsset, og som Kritiken vil
opdage? Hvis saa var, maatte vi standse ved det Resultat, at kun en
indskrænket, afmægtig Viden kan være sand; medens en fuldkommen, almægtig
Viden nødvendig maa være noget i sig selv Usandt. Unegtelig er her en
Modsigelse, dog ligger den, naar vi see nøiere til, ikke i den fuldkomne
Viden, men i den ufuldkomne.

Vi have i det Foregaaende seet, at den sandsende Objectivering, om den end
tager Tænkningen til Hjælp og sætter alle baade physiologiske og
psychologiske Betingelser i Requisition, aldrig kan komme til Rette med
Objectet selv. Hver Gang vi ville lade det uobjectiverede Object staae, saa
forsvinder det, og naar vi have bragt det til at forsvinde, see, saa staaer
det der igjen. Det er virkelig en Modsigelse, og denne Modsigelse kan ikke
hæves, med mindre vi forudsætte en fuldkommen Viden. Intet er objectivt,
som ikke er objectiveret; og dog faaer den sandsende Bevidsthed et Substrat
til Rest, som er objectivt uden at være objectiveret: dette gaadefulde
Objectiveringssubstrat er Realiteten selv, den Realitet, som den mægtige
Viden sætter og hævder. Formedelst dette os fremmede Mellemværende, vor
oprindelige Skranke, faaer den indskrænkede Viden alt Objectivt begrundet,
vel ikke direct paa sig selv, men dog indirect paa Viden, saa at
Erkjendelsens Kredsløb kan blive sluttet.

Det er umuligt at komme til Bunds i den ufuldkomne Viden og gjøre alsidig
Rede for dens Beskaffenhed og Værd uden stadig at have den fuldkomne som et
Ideal \textit{in mente}. Den Afviisning, at Antagelsen af en fuldkommen
Viden er en blind Forestilling, en løs Hypothese, siger Intet, naar man kun
vil indrømme, at Glimt af det Fuldkomne ere nødvendige for at belyse det
Ufuldkomne.

Skjøndt det i vor vidende Bevidsthed undertiden synes at vrimle af
Billeder, Forestillinger og Tanker, næsten som af Bier i en Kube, saa
forholder det sig dog, naar vi see nøiere til, ingenlunde saaledes. Det er
tvertimod kun faa Billeder, yderst faa Forestillinger, som den bestemmende
Tænkning kan holde tilbage, medens Sværmen forsvinder. Af tre
Forestillinger staaer hver Gang kun den ene i fuld Belysning, medens de to
andre ere mere eller mindre fordunklede. Naar flere Ting paa een Gang
tildrage sig vor Opmærksomhed, og det er, som om vi samtidig i lige Grad
beskjæftigede os med dem alle, da er det en Indbildning, som den hurtig
skiftende Opmærksomhed fremkalder. Nøiagtig talt, er der kun eet Element,
der i det givne Øieblik kan faae Plads i Bevidsthedens Forgrund, thi
Pladsen er snæver. I Forhold til som det Hele er lyst, ere Delene dunkle,
og naar en enkelt Deel træder frem, maae de andre træde tilbage. Vor Viden
er indskrænket af Rum og Tid; det er en Ufuldkommenhed; men det er egenlig
ikke Rum og Tid, det er Skrankerne i Rum og Tid, der foraarsage
Ufuldkommenheden.

Det er en Indbildning, at Former som Rum og Tid skulde være udelukkede fra
den fuldkomne Viden, og at vor Viden vilde blive desto fuldkomnere, jo mere
intetsigende den fandt disse Former. Fuldkommengjørelsen er tvertimod
betinget af den Videndes Evne til at apriorisere og tilegne sig Rummets og
Tidens reale Indhold. Dette kan for sandselige Væsners Vedkommende
naturligviis kun skee i Tilnærmelse. Vi bevæge os i vor Bevidsthed fra den
ene Gjenstand til den anden, fra den ene Forestilling, den ene
Tankebestemmelse til den anden, og Bevægelsen kræver Tid. Jo langsommere og
besværligere det gaaer, desto ufuldkomnere benyttes Tiden; jo hurtigere og
lettere, desto fuldkomnere. Jo fyldigere Indholdstilegnelsen er, desto mere
vil Tiden forandre sig fra at være den tomme Form, der venter paa
Udfyldning, det uendelig Langtrukne, hvori det Forbigangne og det
Nærværende og Tilkommende befinde sig under lutter Udstykning, til et
udfyldt Øiebliks Continuitet, hvori de samlede Momenter af Udfyldningens
Indhold gjensidig belyse hinanden: Tiden vil, med andre Ord, nærme sig
Evigheden. Hvor nær den Vidende, der er fordybet i sine Opgaver, kan bringe
Tidens Uro til evig Ro, derom minder hiint Ord af Archimedes: \textit{noli
turbare circulos meos!} Den, der skal orientere sig i Himmelrummet, maa
begynde stykkeviis og tage hver Egn for sig. Udstykningen i Rum er ogsaa
Udstykning i Tid. Det tager Tid, inden man ved Hjælp af Planer, Axer,
Storcirkler og sphæriske Coordinater kan finde sig til Rette med Samtidigt
og Successivt; men Astronomen seer det Hele som med eet Blik: Rummet
indgaaer i Begrændsning, og dets umaalelige Udstrækning overvælder ikke den
Vidende. Her er i det Ufuldkomne et Gjenskin af det Fuldkomne. Tilnærmelsen
beviser den fuldkomne Videns Gyldighed; men for sandsende Væsner er her en
Skranke.

For saa vidt de sandsende Væsner tillige ere tænkende Væsner, har man meent
at kunne blive Endeligheden kvit, overstige Skranken og naae det Fuldkomne
ved at holde sig til Tænkningen som Tænkning. Med en fra sandselig Tyngsel
befriet Anskuen kan Tanken jo fare gjennem Tid og Rum hurtigere end Lynet,
hastigere end Lyset. I de rene Tankeformer skulde vi altsaa kunne opnaae en
uindskrænket, absolut fuldkommen Viden. Desværre en Illusion! Den »absolute
Viden«, der betinges af fortsat Abstraction fra alle Realiteter, er trods
sit rige Indhold af Formbestemmelser uden alt Virkelighedsindhold, uden
Realitet. Ved at tage endog blot et Glimt af Realitet med i en bestemt
Forestilling, vil man let overbevise sig om, at det med Tankens lovpriste
Hurtighed ikke ganske forholder sig, som man sædvanlig antager. Lad en Alen
f.\ Ex.\ gjælde som Maal for en umiddelbart anskuelig Veilængde. Man kan
nu, hedder det, let tænke sig en Miil. Men naar en dansk Miil virkelig skal
gjennemløbes i Tanken, da maa det jo skee ved at lægge en Alen 12,000 Gange
til sig selv i en ret Linie og lade Anskuelsen ophæve Udstykningen. Hvor
meget Tanken end skynder paa, saa mærker man dog, at det tager Tid. Man kan
behøve et Par Dage til at faae Opgaven indøvet, og saa vil Tanken med den
bestemte Forestilling dog kun tilnærmelsesviis gaae sammen i Anskuelsen. Og
endda skulde Tanken med Lethed følge Lysets Fart igjennem det uendelige
Rum?

Vor Videns Ufuldkommenhed skal, siger »Kritiken«, vise sig deri, at vi kun
erkjende Tingenes Love, men ikke deres Væsen; selvfølgelig have vi intet
Begreb om en fuldkommen Viden. --- Ja, svare vi, saa længe man kun
betragter Lovene som almeengjældende Bestemmelser for den Orden og
Sammenhæng, hvori Tingene træde os imøde, saa længe man endnu er uvis om,
hvorvidt disse Bestemmelser ligge i Tingene selv eller blot beroe paa vor
Opfattelse af Tingene, er det ganske naturligt, at en Erkjendelse af Love
og en Erkjendelse af Væsen maae falde vidt ud fra hinanden. Det er
unægtelig en Udvei, naar man vil være den endeløse Tvivlen om Lovenes
almene Gyldighed kvit, at lægge dem ind i vor egen Bevidsthed og bestemme
dem som Forstandsformer, der ere nødvendige for al Erfaring. Naar ingen
Erfaring er mulig uden Love, saa kan man være vis paa, at Erfaringstingene
nok skulle rette sig efter Lovene. Men Opfattelsen lider af betænkelige
Svagheder. Enten maae Lovene, ifølge denne Synsmaade, være yngre end
Tingene eller ældre end Bevidstheden; men begge Antagelser ere lige
urimelige. Man beraaber sig paa, at Erfaringstingene ikke ere Ting i sig
selv, men Phænomener for den Bevidsthed, som gjør Erfaringer. Heri ligger
ganske vist den Sandhed, at Phænomener og Love tabe al Betydning, naar der
abstraheres fra en objectiverende Viden; men hvorledes det gaaer til, at
Tingene som Ting, der dog maae have en vis Andeel i Bestemmelsen af
Phænomenerne, slet ingen Andeel skulde have i Bestemmelsen af Lovene,
bliver ikke opklaret. Have Tingene i deres Væsen maaske slet ingen Love,
fordi der ingen Forstand er i dem? eller ere Væsens-Lovene i sig selv
Udtryk af en saa fuldkommen Forstand, at de aldeles ikke passe for vor
Forstand? Begge Meninger ere lige meningsløse.

De, der antage, at Lovene hverken ligge i Tingenes Væsen ei heller i vor
Bevidsthed, men ere almindelige Formler for den Orden, vi abstrahere af
stedfindende Forhold, nedsætte vor Viden til noget saa Væsenløst, at den
kun kan bestaae i at paavise, samle og ordne en Vrimmel af Kjendsgjerninger.
Kjendsgjerningerne selv kunne ikke forstaaes, de maae tages umiddelbart som
de gives; den Orden og Sammenhæng, hvoraf Forstaaelsen afhænger, kan heller
ikke forstaaes; thi den forefundne Orden er selv en Kjendsgjerning:
saaledes bliver Alt uforstaaeligt. Vel kan man ved Iagttagelser,
Beregninger og Inductioner overbevise sig om, at Noget er sædvanligere end
Andet, Noget almindeligere end Andet, at Noget saa ofte gaaer forud, at det
kan kaldes Aarsag, og at Andet saa ofte, ja uden nogen os bekjendt
Undtagelse, følger efter, at det kan kaldes Virkning; men at Aarsager og
Virkninger ere virkelig til, eller at Lovene ere væsenlige, i sig
nødvendige, derom kan man ikke overbevise sig. Man bruger Forstanden uden
at kjende Forstanden, og nøies med en Forvisning uden Indsigt.

Der er ganske vist en uoverstigelig Kløft mellem den fuldkomne Viden og den
ufuldkomne, og som Følge heraf en uoverstigelig Kløft imellem den
fuldkomne og ufuldkomne Forstand; men Kløften gjør intet Brud paa Lovene
selv, de evige Love. Om Videnskaben stadig skred frem i Billioner Aar, og
vi oplevede de sidste Resultater: vi vilde dog aldrig naae til en simultan,
baade intuitiv og distinctiv Viden af Universets virkelige Indhold; men
allerede paa det Trin, hvor vi nu ere, kunne vi erkjende, at vor subjective
Forstand i bevidst Form maa være væsenlig Eet med den ubevidste Forstand,
der har sit objective Udtryk i Tingene, saa sandt vi kunne erkjende, at
vore Hjerner, vore Bevidsthedsorganer ere Naturlovene undergivne. Vi kunne
indsee, at den objective Forstand i Tingene maa have sin oprindelige Grund
i en vidende Forstand, i en fuldkommen Viden; thi en Forstand, som Ingen
forstaaer, og som ikke forstaaer sig selv, er umulig.

Den ufuldkomne Viden maa, saa at sige, gjøre sig til Discipel af den
fuldkomne, hvad, ifølge Tingslovenes og Tankelovenes væsenlige Eenhed, da
ogsaa er muligt. Ligesom ingen Tanke har Betydning for sig, men kun i
Vexelforhold til andre Tanker, saaledes har heller ingen existerende Ting
udelukkende Betydning for sig, men kun i Forhold til andre Ting.
Vexelbestemmelse af Led og Forhold har lige saa fuldt Gyldighed for den
fuldkomne Viden som for den ufuldkomne. At den ufuldkomne Viden maa stave
sig frem, medens den fuldkomne seer Alt som med et eneste Blik, gjør i
Forhold til Lovenes Væsen ingen Forskjel. Det Blik, der seer Conseqvensen i
Forudsætningen, det Afledede i det Oprindelige, Delene i det Hele, seer med
det Samme den væsenlige Forskjel; hvis Forskjellen forsvandt, løb Alt ud i
Eet, og dette Eet blev til Intet. Men at Eenheden har Mangfoldighed og
Mangfoldigheden Eenhed, at Identiteten har Forskjel og Forskjellen
Identitet, at Delene er det Heles og det Hele Delenes Bestaaen, er netop
væsenlige Sandheder, som den ufuldkomne Viden idelig maa indøve, for ikke
at gaae under i lutter Udstykning. Hver enkelt Lov, opfattet for sig selv,
er et Abstractum og kan for saa vidt siges at savne Realitet; men jo
fuldkomnere Viden bliver, desto tydeligere viser det sig, at de
forskjellige Love i Ting og Virkelighed integrere hverandre. For at sikkre
en enkelt Lov, er det nødvendigt at udhæve en fælles Egenskab ved flere
Ting og betragte deres Forhold med Hensyn til denne Egenskab; men ingen
Ting kan existere med een eneste Egenskab: Legemet, som falder, Magneten,
som tiltrækker, ere under Opfyldelsen af hver paagjældende Lov i Forhold
til Lys, Varme, Elektricitet o.s.v.\ andre Love undergivne. Det tunge
Legeme er mere end tungt; det har en Mængde andre Egenskaber og er ved
hver særegen Egenskab impliceret i et System af tilsvarende Love.

Videnskaben lægger an paa at reducere alle Love til een Grundlov, og
beviser just herved, at den har den fuldkomne Viden til Ideal. Det er
Idealet af en uindskrænket Forstand, en Indbefatten af alle Ting og alle
Love under en Alt gjennemskuende Videns Eenhed, der som stiltiende
Forudsætning giver Videnskaben sit ideelle, aprioriske Grundlag. Tages
dette Grundlag bort, saa maa Videnskaben med alle dens Lærebygninger tabe
det imponerende Høidemaal, og Forskningerne synke ned til at tjene en
Videns Begjær, der aldrig mættes; Gudinden med Hadeshjelmen og det
gjennemborende Blik maa da miste sin jomfruelige Reisning, bøie sit Hoved
og finde sig i at gaae i det Nyttiges Trædemølle.


\chapter{Systematisk Viden}

\section{System og Methode}
\label{sec:8}
% pp. 69--80

\noindent At al videnskabelig Viden, væsenlig seet, maa være systematisk,
følger ligefrem af, at Erkjendelsens enkelte Led kun faae Betydning, naar
de sees i indbyrdes Sammenhæng. De sammenhængende Led ville da danne et
Hele, en Mangfoldighedseenhed, bestaaende deri, at hvert enkelt Led
indtager sin Plads i Forhold saa vel til de øvrige Led som til det Hele
selv; et saadant Hele er et System. Hvor mange Led der nødvendig behøves
for at danne et System, kan let beregnes: deres Tal maa være større end
$1$ og mindre end $\infty$. At der maa være mere end eet Led, forstaaes af
sig selv, thi der maa være Forhold og Sammenhæng; at intet bestemt,
tilværende System kan have uendelig mange Led, indsees deraf, at et
bestemt og dog uendeligt Tal ikke gives. Naar der tales om et uendeligt
System, da kan der kun tænkes paa uoverskuelig Storhed eller paa uophørlig
Vorden.

Et System af flere Led kan opløses i mindre Systemer, og flere mindre
kunne igjen sammensættes til et større. Maaden, hvorpaa det skeer,
afhænger af Methoden, denne af Sammenhængens Beskaffenhed, og denne igjen
af Leddenes Beskaffenhed. Bestaaer Systemet af existerende Led, og staaer
dets Sammensætning i vor Magt, da vil Methoden være: deels at danne et
Hele af tilsvarende Dele, deels at opløse baade det Hele og dets Dele i en
real Eenhed, en Anlægseenhed, hvoraf baade Hele og Dele igjen maae
fremgaae; deels endelig at komme til Indsigt i de Dannelseslove, paa hvis
Opfyldelse Systemet væsenlig beroer. Jo mere Leddenes betinget-nødvendige
Sammenhæng har Charakteer af udvortes og endelige Forbindelser, desto mere
vil Systemet være mærket paa den ene Side af Tilfældighed, paa den anden
af Vilkaarlighed. Jo inderligere og væsenligere Sammenhængen er, desto
mere vil en Alt beherskende Nødvendighed i Delenes Dannelse ud af
Heelheden paatrykke Systemet sit eiendommelige Præg. Er Sammenhængen
endelig saa gjennemgaaende, at Delene blive Momenter i Heelheden, Leddene
Momenter i Forholdene, da vil Heelhedsudviklingen være det Afgjørende og
Dannelsens deelvise Resultater paa alle Punkter vise tilbage til Anlæget;
Nødvendigheden vil gaae i Eet med Friheden, og vi have da et teleologisk
System.

At et teleologisk System maa være Videns Ideal, kan sees paa den
videnskabelige Methode. Methode er Omvei, Objectiveringens, Erkjendelsens
Vei fra de første nødvendige Forudsætninger til det søgte Resultat.
Resultatets Opnaaelse er Maalet. Den Omvei og Fremgangsmaade, der
charakteriserer Methoden, er altsaa: Veien til Maalet. Enhver videnskabelig
Opgave er teleologisk, thi Opgaven sætter et Formaal, og Methoden angiver
Midlerne, d.v.s.\ Betingelserne for dens Løsning: Beviset er teleologisk,
thi det er methodisk. Men paa Endelighedens Trin ville de objective Forhold
i Sagen ingenlunde svare til den subjective Bevægelse i Methoden. Ved
Spørgsmaal om Forholdet mellem en udvendig Vinkel paa Trekantens forlængede
Side og to indvendig modstaaende har Fremgangsmaaden i Beviset vel
Resultatet til Formaal, men Vinklernes Størrelsesforhold retter sig ikke
efter Bevisets Gang; dette Forhold bevæger sig ikke, det vorder ikke, det
er. Det kan være en Opgave at undersøge, hvorledes Vandet fremkommer af
sine Bestanddele; men er end Fremgangsmaaden ved Undersøgelsen teleologisk,
saa kan Bevægelsen i Sagen dog ikke siges at være teleologisk; intet
chemisk Experiment kan godtgjøre, at Ilten forener sig med Brinten i det
Øiemed at danne Vand. Et genetisk System, som fra objectiv Side
fyldestgjør Methodens subjective Fordringer, kan ved formelle
Constructioner, aprioriske Frembringelser vistnok deelviis, fremkomme; men
i Realiteternes Verden gives det ikke. Fra dette Synspunkt kan
Forstandskritiken, idet den søger at bekæmpe den saakaldte speculative
eller immanente Methode, der identificerer Begrebsbevægelse og
Sagbevægelse, paa Forhaand siges at have Seiren ihænde. Schellings Ord:
»Ueber die Natur philosophiren heisst die Natur schaffen«, giver jo
vistnok en Mening, for saa vidt en indtrængende Forstaaelse af Tingen er
betinget af en Efterdannelse, en intellectuel Gjenfrembringelse, hvorunder
Objectet fra Grunden af opstaaer Led for Led og bliver til i Bevidstheden.
Men Tilblivelsen er kun ideel, ikke reel; den Schellingske Propheti, at
»den Tid vil komme, da Videnskaberne mere og mere ville ophøre, idet ikke
mere Mennesket seer, men den evige Seen selv er seende i Mennesket«, vil
ikke gaae i Opfyldelse: hvad han kalder »Videnskaben selv« kan umulig
overleve »Videnskaberne«.

Det tør imidlertid ikke oversees, at den kritiske Opfattelse ogsaa har
sine svage Sider. Ved at holde sig til Brudstykker kan man meget vel
efterforske Kjendsgjerninger og Love uden teleologiske Hensyn; naar man
kun forudsætter et Aarsagsforhold mellem Kjendsgjerningerne, saa har man
Alt, hvad man behøver. Dersom A (Forudsætningen med Aarsag og Betingelser)
er, saa maa B (Conseqvensen, Resultatet) ogsaa være. Dersom en Planetmasse
indeholder Vandets Bestanddele, saa maa der under indtrædende Betingelser
danne sig Vand. Man bevæger sig uophørlig mellem Forudsætning og Resultat,
Resultat og Forudsætning, og faaer ethvert Hele begrundet som Produkt af
Delenes Vexelvirkning.

I hvilket Omfang den videnskabelige Methode kan tilveiebringe et genetisk
System uden Teleologi, derpaa er den Kant-Laplace'ske Theori om Solsystemets
Dannelse et storartet Exempel. Men det genetiske System er kun
halvgenetisk, saa længe man indskrænker sig til at paavise det Heles
Tilblivelse uden at paavise Delenes Tilblivelse. Skal Delenes Tilblivelse
begrundes lige saa vel som Heelhedens, da maa Systemet nødvendig blive
teleologisk. Det maa da vise sig, at Delene saa vel som det Hele ere
udviklede af et fælles Anlæg, at Delenes reale Mulighed netop er bestemt af
Hensyn til det Hele, d.v.s.\ at Resultatet er præformeret i Forudsætningen,
og dets Opnaaelse Udviklingens Formaal.

Men, siger Kritiken, et saadant dobbelgenetisk System gives ikke, og kan
ikke gives. Naar vi spørge: hvorfor? faae vi to forskjellige Svar. Der maa
for det Første være Noget, der staaer fast, et Givet, hvortil man kan holde
sig og hvorfra man ved enhver Begrundelse kan gaae ud; men det er netop
hvad Teleologien miskjender. Først antager man Forudsætningen, for at faae
Resultatet begrundet, saa gjør man igjen Resultatet til Forudsætning, for
at faae Forudsætningen begrundet; det er en subjectiv Cirkelbevægelse.
Virkeligheden opnaaer man kun at gjøre Forudsætning og Resultat lige
hypothetiske. For at skjule det Hypothetiske og faae Bevægelsen til at see
ud som et Fremskridt, underskyder man en altbestemmende Viden og Villen,
d.v.s.\ man underskyder det Vilkaarlige; derfor maa al Teleologi forkastes
af Videnskaben.

Naar herimod gjøres gjældende, hvor uvidenskabeligt det er at adskille de
existerende Ting i to Klasser: saadanne, som maae have en Grund, nemlig de
sammensatte, og saadanne, som ingen Grund have, men ere fordi de ere,
nemlig de enkelte Elementer, Bestanddelene, og hvor inconseqvent det er at
gjøre Grunddelenes Existens tilfældig, for at gjøre Heelhedernes Dannelse
nødvendig, saa faaer man en anden Forklaring. Et teleologisk System er
umuligt for sandsende Væsner; al Forskning gaaer ud fra et faktisk Givet,
og Videnskaben maa indrette sin Methode derefter. Ved denne Opfattelse
kunne vi acqviescere, thi ved den er hele Synsmaaden forandret. --- Naar
man ikke vil indsmugle det Tilfældige for at undgaae det Vilkaarlige, maa
det teleologiske System nødvendig vedblive at være Videnskabens Ideal. Det
bliver da Kritikens Sag at vaage over, at vor Erkjendelses Vilkaar
respecteres, at man ikke ved Indbildningens Flugt ud over Skrankerne
forgriber sig paa Idealet.

Det dogmatiske System, hvis væsenlige Særkjende er en naiv Forvisning om
den reale Værens og den menneskelige Videns fuldkomne Overeensstemmelse,
har ikke sjældent gjort sig skyldigt i teleologiske Misgreb. For den
teleologiske Dogmatiker er det en afgjort Sag, at Delene ere til for det
Heles Skyld, at en Forstaaelse af deres Sammenhæng først kan opnaaes, naar
man har indseet, hvad Meningen er med det Hele. Idet Formaalserkjendelsen
bliver det Afgjørende, ligger det nær, at den vidende Bevidsthed søger
Maalestokken i sit eget subjective Væsen og, uden at vide deraf,
underlægger Tilværelsen en Mening, som bedst svarer til, hvad den
menneskelige Subjectivitet maa ansee for det Sande, det Gode, det Skjønne.
Dogmatikerens naive Tillid til sin »Totalanskuelse« sees tydelig af den
Maade, hvorpaa han behandler de enkelte Phænomener og Kjendsgjerninger.
Intet nok saa gjenstridigt Factum generer ham; han gaaer ikke af Veien for
noget mødende Tilfælde. Ved Omforklaring, Abstraction og Bortforklaring
faaer han Bugt med Alt, hvad der ellers ikke vil gaae op i
»Totalanskuelsen«.

Det dogmatiske System er i det Hele overveiende synthetisk. Ud fra et
almindeligt Princip, det være iøvrigt teleologisk eller ikke teleologisk,
skrider det frem gjennem særegne Bestemmelser, definerende, ræsonnerende,
demonstrerende, saa at hele Beviisførelsen faaer et mathematisk Anstrøg
(\textit{methodus mathematica}). Med Anstrøg af Uomstødelighed optræder det
dogmatiske System som et Autoritetssystem sædvanlig saaledes, at
Autoriteten er udvortes, for saa vidt den skyldes en herskende Overlevering,
men dog ogsaa saaledes, at den hviler paa den Overlegenhed, hvormed
Totalanskuelsen er gjennemført, og for saa vidt er indvortes.

Den gamle Ptolemæiske Verdensbygning havde den Aristoteliske Dogmatisme
til Grundlag. Hvor huult Grundlaget var, blev aabenbart, da Copernicus
brød med Autoriteten og fik Systemet nedsat til en Hypothese. Det var et
Vendepunkt ikke blot i Astronomiens, men i Videnskabernes Historie. Jo mere
Kritiken fik Magt over Autoriteten, jo skarpere Iagttagelserne af de
virkelige Kjendsgjerninger i alle Enkeltheder bleve, desto mere traadte
Analysen i Forgrunden. Kritiken er intet System, den er en Methode: den
kritiske Methode er væsenlig analytisk. Synthesen gaaer fra det Hele til
Delene, fra det Almene til det Concrete; Analysen adskiller det Hele for at
undersøge Delene. Ved en egen Forening af Induction og Abstraction baner
den sig Vei gjennem det Concrete til det Almene. Beskjæftiget med
Undersøgelser af det Enkelte og Lovene, af Elementerne og deres saa vel
virkelige som mulige Combinationer, har Kritiken ingen synderlig Sands for
Totalanskuelser; men den er aarvaagen paa sin Post, hvor det gjælder at
prøve Forudsætningernes Realitet, Fremgangsmaadernes Conseqvens og
Bevisernes Styrke.

Da Synthese og Analyse betinge hinanden gjensidig, saa indsees heraf, at
intet System kan være enten udelukkende synthetisk eller udelukkende
analytisk; Modsætningen beroer paa, at snart det ene, snart det andet af de
to Momenter danner Grundlaget. I Erkjendelse heraf har der udviklet sig et
System, som hverken vil være analytisk eller synthetisk, men
analytisk-synthetisk fra Først til Sidst, en Methode, som hverken vil være
dogmatisk eller kritisk, men dogmatisk-kritisk fra Begyndelsen til Enden,
det er det speculative System med den dialektiske Methode.

Den Strid, som heraf opstaaer mellem Kritik og Dialektik, er af
gjennemgribende Betydning; ingen af de stridende Parter kan enten absolut
seire eller absolut tabe, thi Seier og Nederlag afhænger af Stridsemnet.
Man kan med Rette sige, at det er en Strid om Grund og Existens, naar man,
vel at mærke, henfører alle Væsensbestemmelser under Grunden og ved
Existens ikke tænker paa Begrebet, men paa existerende Ting. I sidste
Tilfælde vil Kritiken altid beholde Overvægten; Kritikens Grændser ere
lutter Existensgrændser, Contradictionsprincipet hersker, og modsatte Led
udelukke hinanden. Hvad Grunden, det Ikke-Existerende, som bærer
Existensen, angaaer, vil Modsigelsens Grundsætning derimod ikke formaae at
hindre Forskjellen fra at falde sammen i Eenhed, for igjen at fremgaae af
Eenhed. Vi skulle for Methodernes Skyld oplyse Striden ved nogle Exempler.
At Tingene have Egenskaber, at de forandre sig, opstaae, forgaae o.s.v.,
derom ere Kritik og Dialektik indbyrdes enige; men naar det gjælder om at
bestemme, hvori Forandring egenlig bestaaer, hvorledes en Ting forholder
sig til sine Egenskaber, og hvad den Grund er, hvoraf det Existerende
opstaaer, og hvori det igjen løser sig op, saa reiser sig en Strid som paa
Liv og Død. Kritiken paastaaer, at enhver Ting er hvad den er, og kan
umulig være, hvad den ikke er; Dialektiken indvender, at dersom dette skal
gjælde ubetinget, saa kan ingen Ting forandre sig, Intet opstaae og Intet
forgaae. For at forstaae, hvad Forandring væsenlig er, maa man indsee, at
Tingen som Ting har sin Tilværelse i lutter Modsigelse. Tingen forandrer
sig, fordi den paa een Gang er baade Noget og Andet; den forandrer sig,
fordi den baade er Eet med og dog forskjellig fra sine Egenskaber; den
forandrer sig, fordi den har en Grund og er baade Eet med og dog
forskjellig fra sin Grund. At Tingen forandrer sig, vil sige, at dens
Tilværelse er dens Begreb, og Forandringen dens Methode, dens
Værensbestemmelser ere lutter Tankebestemmelser. Bevægelserne i dens
Forandring, Sagbevægelse og Begrebsbevægelse, ere i deres Forskjel
identiske; d.v.s.\ den dialektiske Methode er immanent.

Denne Identificeren af det Forskjellige, denne Forblanding af det
Uligeartede, dette aabenbare Brud paa Modsigelsens Grundsætning forarger og
oprører Kritiken. Er Noget det Samme som Andet, saa er Sort Hvidt, og Hvidt
er Sort, saa er Lampen Bord og Bordet Lampe, Stenen Menneske, og Mennesket
Steen! Saa paakaldes den sunde Forstand, og den sunde Forstand maa give
Kritiken Medhold. Desuagtet er Seiren ikke saa glimrende, som det ved
første Øiekast seer ud. Holde vi os til de rene, almene Tankebestemmelser,
da gjør Dialektiken ingenlunde Brud paa Modsigelsens Grundsætning;
tvertimod, den immanente Methode lægger netop an paa at faae Modsigelsen
grundig opløst. Dens Thesis er, at Noget ikke er Andet; dens Antithesis er,
at Noget er Andet; men den bliver lige saa lidt staaende ved den sidste
Sætning som ved den første, thi Synthesis er, at Noget idelig er i Overgang
til Andet, at Overgangen, Forandringen, netop er Modsigelsens Løsning. I
Analysen af rene Begreber kommer Kritiken altid tilkort.

Der kan gives en spidsfindig Dialektik, men der kan ogsaa gives en
spidsfindig Kritik. I de rene Begrebers Analyse er Kritiken altid
sophistisk. Den støder overalt paa Modsigelsen, men istedenfor at binde
alvorlig an med den og løse den op, gaaer den ved Hjælp af allehaande,
Sagen selv uvedkommende, Hensyn (\textit{diversi respectus}) uden om
Opgaven, og filtrer sig ind i Modsigelser. Sagen er, at Tankernes,
Begrebernes indbyrdes Forhold ere Uendelighedsforhold, de kritiske
Grændser derimod lutter Endelighedsgrændser. Dialektik er i Abstractionernes
og Uendelighedens Sphære slet ikke Andet end gjennemført Kritik.

Paa Existensens, d.v.s.\ paa Endelighedens Trin, har Kritiken derimod
fuldstændig Ret. Alle Realitetsgrændser ere Endelighedsgrændser, kritiske
Grændser. I det Øieblik den immanente Dialektik vover at betræde
Realitetens Grund, kan Kritiken tage den tilfange. Hvor Forholdene ere
endelige, maae de endelige Hensyn (\textit{diversi respectus}) gjælde.
Hvad der i een Henseende kan være et Høire, kan i en anden Henseende være
et Venstre; men derfor er Høire dog ikke Venstre, og Venstre ikke Høire.
Kritiken lærer os, at der til forskjelligartede Sphærer hører
forskjelligartet Erkjendelse; Kritiken lærer, at forskjelligartede
Videnskaber give forskjellige Systemer og kræve forskjellige Methoder.

\section{Videnskabernes systematiske Sammenhæng}
\label{sec:9}
% pp. 80--90

\noindent Med Hensyn til Vished, Begrundelse, Indsigt og sammenhængende
Erkjendelse er der noget for alle Videnskaber Fælles, Noget, som under
enhver Form gjør Viden videnskabelig. Paa Grund af det Fællesvidenskabelige
ere alle Videnskaber som een Videnskab; eet Blod rinder i deres Aarer, de
ere af een Familie, ligesom alle Muserne ere Søstre. Men med Hensyn til
Stof og Materiale, til Opgave og Fremgangsmaade, ere de høist forskjellige.
Denne Forskjellighed rører ikke blot Phænomenernes Overflade, den trænger
ind i Videns Væsen, den modificerer Indsigten. Undersøgelse af en Jordart
giver en anden Art Indsigt, end Undersøgelse af en historisk Begivenhed
eller af en geometrisk Curve. Man maa altsaa betragte Videnskabernes
Sammenhæng deels fra Fleerhedens og Forskjellighedens, deels fra Lighedens
og Eenhedens Synspunkt.

Begge Synsmaader ere traadte eensidig frem i Behandlingen af hvad man
kalder Videnskabernes Encyklopædi. Allerede i Oldtiden gjordes Forsøg paa
at sammenfatte den Kreds af Videnskab og Kunst, som enhver fribaaren Græker
og Romer maatte besidde, systematisk efter enkelte Discipliner; af den
græske Tvedeling: Musik og Gymnastik, fremgik siden de saakaldte syv frie
Kunster: »Grammatik, Dialektik, Rhetorik, Arithmetik, Geometri, Astronomi
og Musik«. Da Tænkningen stod over Erfaringen, optog Philosophien de
særskilte Videnskaber i sine Lærebygninger og endte med den stoiske
Tredeling: Logik, Physik, Ethik. I Middelalderen havde man Compendier,
mere eller mindre fuldstændige Oversigter over det samtlige Indbegreb af
menneskelig Viden paa den Tid under Titler som: \textit{Summa, Speculum,
Theatrum vitæ humanæ} o.\ dsl. Senere opgav man det systematiske Hensyn
endog i den Grad, at man betjente sig af en alphabetisk Anordning af
Stoffet og indrettede Encyklopædierne som Realordbøger. Den franske
Encyklopædi (der paabegyndtes 1749), redigeret af Diderot, var et
storartet litterært Arbeide, hvorved det attende Aarhundredes Aand med
kritisk Forstand i klare Former bekæmpede Alt, hvad den ansaae for blind
Overtro og uhjemlede Traditioner. Men saa værdifulde encyklopædiske
Ordbøger end kunne være, saa er det dog kun paa en indirect og temmelig
udvortes Maade, at de fremme nogen Indsigt i Videnskabernes væsenlige
Sammenhæng.

For en almindelig Videnskabslære ere de encyklopædiske Forsøg, som ledes
af almene, omfattende Principer, derimod af større Værdi; skjøndt det
Eensidige ogsaa her er vanskeligt at undgaae. Den første philosophisk
begrundede Oversigt over Videnskaberne har Baco af Verulam leveret i
\textit{Novum Organon} (1620) og \textit{De dignitate et augmentis
scientiarum} (1623); men trods den aandrige Sammensmeltning af Tanke og
Billede, Speculation og Empiri, formindskes dog dens videnskabelige
Vægtfylde betydelig, naar den underkastes en kritisk Prøvelse. Af dybt
indgribende Betydning ere derimod de Fremstillinger, der have enten Kants
eller Comte's eller Hegels Principer til Grundlag.

Der gives, ifølge den Kantiske Fornuftkritik, to forskjellige Kilder,
hvoraf vore Erkjendelser udspringe, Sandsen og Fornuften. Fra hiin kommer
Sandseerkjendelsen, fra denne de rene Fornufterkjendelser. De sidste falde
i Mathematik og Philosophi. Vi faae altsaa tre forskjellige Slags
Erkjendelser: empiriske, mathematiske, philosophiske. Enhver af disse
Erkjendelsesmaader har sin særegne Charakteer: Tilfældighed og Anskuelighed
tilkommer den empiriske, Anskuelighed og Nødvendighed den mathematiske,
Nødvendighed uden Anskuelighed den philosophiske.

»I de mathematiske Erkjendelser blive vi os Tingenes Sammensætningsformer,
i de philosophiske deres Sammenknytningsformer \textit{in abstracto}
bevidst. Kant kaldte hine: Former for den figurlige Synthesis, d.\ e.\
anskuelige Forbindelsesformer, disse: Former for den intellectuelle
Synthesis, d.\ e.\ kun tænkelige Forbindelsesformer, Love. Betragtningen
af de forskjellige Erkjendelsesarter fører til en videnskabelig
Architektonik, d.\ e.\ til en Lære om et System af alle menneskelige
Videnskaber. De tre Erkjendelsesmaader svare til tre forskjellige
Systemformer. Den rene philosophiske Erkjendelsesform hører under det
kategoriske System; den rene mathematiske Form under det hypothetiske
System, og den rene empiriske Gehalt under det conjunctive System.
Philosophien er en Erkjendelse af blotte Begreber. Med en saadan
Erkjendelsesmaade er ingen anden systematisk Afleden mulig end ved en
kategorisk Indordnen af særegne Tilfælde i Reglens Almindelighed.
Mathematiken er en Erkjendelse ved Begreber i reen Anskuelse. I denne
Anskuelse er det mindst Sammensatte det Oprindeligste og den Grund, hvoraf
nye Følger afledes ved bestandig nye Sammensætninger. Det Sammensattes
Afledning af det Usammensatte skeer her ved hypothetisk Tankeforbindelse.
Endelig har den rene Empiri kun med Opfattelsen af det enkelte ved
Sandseanskuelsen Givne at gjøre. Hver Kjendsgjerning beroer paa sig selv,
og den ene staaer ved Siden af den anden i et conjunctivt System, der
sammenordner Dele i det Hele.« (Apelt: \textit{Die Theorie der Induction.})

Det vil efter dette Schema neppe være muligt at danne sig noget Begreb om
de særlige Videnskabers eiendommelige Charakteer eller indbyrdes
Sammenhæng. Allerede den kritiske Adskillelse mellem de tre
Erkjendelsesmaader viser sig høist betænkelig. Hvorledes vil man kunne
tænke sig en Adskillelse mellem Sammensætnings- og Sammenknytningsformer
saaledes gjennemført, at der i Virkeligheden danner sig to forskjellige
Videnskaber, af hvilke den ene har lutter Sammensætningsformer, Former for
den figurlige Synthesis, lutter anskuelige Forbindelsesformer, medens den
anden har lutter Sammenknytningsformer, tænkelige Forbindelsesformer, Love?
Kan der da være lutter Love, Nødvendighedsforhold uden anskuelige Led i den
ene Videnskab, og lutter Sammensætninger af anskuelige Led uden Love i den
anden? Man prøve om et saa simpelt geometrisk Theorem som: To Nabovinkler
udgjøre tilsammen to Rette, lader sig godtgjøre ved anskuelige
Sammensætninger uden intellectuel Synthesis, og man vil øieblikkelig
opdage det Umulige. To bestemte og for saa vidt anskuelige Vinkler, man kan
afsætte, give jo kun et Exempel, et enkelt Tilfælde, der kan ombyttes med
hvilke som helst andre Tilfælde; kan man nu afsætte alle mulige Vinkler med
samme Toppunkt paa samme Side af en ret Linie? For at faae Theoremet
beviist, maa man nødvendig opløse den figurlige Synthesis i den
intellectuelle. Vælge vi omvendt et Exempel blandt Sammenknytningsformerne,
de blot tænkelige Forbindelsesformer, mon vi da kunne undvære det
Anskuelige? At enhver Virkning maa have sin Aarsag, tør vel gjælde for et
Exempel paa en intellectuel Synthesis; men hvad er Virkning, hvad er
Aarsag? Uforstaaelige Ord, naar dertil ikke knyttes bestemte Forestillinger.
Men hvorfra skulde Forestillingerne komme, dersom der intet Anskueligt var,
hvorfra de kunde hentes, intet Exempel, hvorpaa de kunne støttes? Sandheden
er, at der i enhver mathematisk Sætning maa være Noget, som hører til det
blot Tænkelige, ligesom der i enhver philosophisk Sætning maa være Noget,
som viser hen til det blot Anskuelige. Den eneste Forskjel, som
Fornuftkritiken her kan beraabe sig paa, er den, at det Anskuelige er langt
mere fremtrædende i visse Dele af Mathematiken, navnlig i den synthetiske
Geometri, end i Philosophiens abstractere Former. Dog, selv om vi lade den
Adskillelse gjælde for principiel, der dog kun grunder sig paa et:
\textit{a potiori fit denominatio}, hvorledes forholder det sig da med det
Empiriske? Dersom de empiriske Videnskaber ikke kunne drive det videre end
til en paa »Tilfældighed og Anskuelighed« støttet Combination af
Kjendsgjerninger, hvor den ene staaer ved Siden af den anden i et
conjunctivt System uden Nødvendighed --- thi Mathematik og Philosophi
lægge jo Beslag paa Alt, hvad der hører Nødvendigheden til --- hvad skal
man saa dømme om deres videnskabelige Værd?

Auguste Comte, Positivismens, d.\ e.\ »den positive Philosophis«,
Grundlægger, gaaer ud fra Umuligheden af at naae absolute Begreber, og
opgiver desaarsag ethvert Forsøg paa at finde Universets Oprindelse,
Tilværelsens Formaal og Phænomenernes inderste Aarsag, men søger derimod
ved Forbindelse af Iagttagelse og Forstandsvirksomhed at udfinde de
virkelige Love og uforanderlige Relationer i givne Phænomener.
Positivismen, som i alle Kjendsgjerninger seer Tilsyneladelser af
uforanderlige Naturlove, lægger an paa at reducere Lovenes Antal til det
mindst mulige, uden dog at gjøre sig Haab om at kunne føre dem tilbage til
en eneste Grundlov; den søger, saavidt muligt, at fremstille de
iagttagelige Phænomener som særlige Tilfælde af et ubekjendt
Grundphænomen. Af denne Synsmaade følger conseqvent en Bestræbelse efter at
samle alle Videnskaber til eet System. Først maa der da skjelnes mellem
Theori og Praxis: »den positive Philosophi beskjæftiger sig kun med den
theoretiske Kundskab om Naturen, altsaa kun med de fundamentale Begreber,
derimod ikke med deres praktiske Anvendelse. Indenfor Naturvidenskabernes
Omraade maa der endvidere skjelnes mellem de generelle eller fundamentale
og de particulære, afledede Videnskaber. Astronomien er fundamental,
Geologien afledet; Chemien er fundamental, Mineralogien afledet o.s.v. De
fundamentale Naturvidenskaber søger nu Positivismen at opstille i en ordnet
Rækkefølge, saaledes som deres indre Natur medfører det. Rækkefølgen
udvikles efter en encyklopædisk Lov saaledes, at Mathematiken paa een Gang
bliver Grundvold for og integrerende Deel af Philosophien; derefter følge
Astronomien, Physiken, Chemien, Physiologien og Sociologien eller den
sociale Physik. Disse sex fundamentale Videnskaber udgjøre »de positive
Videnskabers Hierarchi«, der bør lægges til Grund for al fremtidig
videnskabelig Udvikling, fordi deres indre Sammenhæng og Eenhed omfatter
den hele menneskelige Viden.

Det maa indrømmes, at Positivismen har et skarpt Blik for
Naturvidenskabernes eiendommelige Charakteer og indbyrdes Sammenhæng; men
Aandsvidenskabernes Charakteereiendommelighed kjender den ikke. Den
forkaster det teleologiske Princip, og maa allerede af den Grund blive
inconseqvent i sin Behandling af Sociologien.

Fra den absolute Videns Standpunkt er Hegels systematiske Overblik det
meest storartede, og dog i Deductionen selv det meest forfeilede. »I en
sædvanlig Encyklopædi«, siger han, »tages Videnskaberne empirisk, som de
forefindes. De skulle deri fuldstændig opregnes og ordnes derved, at det
Lignende, det under fælles Bestemmelser Sammentræffende, stilles sammen.
Den philosophiske Encyklopædi er derimod Videnskaben om den nødvendige, ved
Begrebet bestemte Sammenhæng, og om den philosophiske Tilbliven af
Videnskabernes Grundbegreber og Grundsætninger. Videnskaberne ere efter
deres Erkjendelsesmaade enten empiriske eller reent rationelle. Absolut
betragtet skulle begge have samme Indhold. Det er den videnskabelige
Stræbens Maal at hæve det, der vides empirisk, til det altid Sande, til
Begrebet, at gjøre det rationelt og derved indlemme det i den rationelle
Videnskab. Videnskabens Hele deler sig da i tre Hoveddele: 1) Logik. 2)
Naturvidenskab. 3) Aandsvidenskab. Naturens og Aandens Videnskaber kunne
da betragtes som de reale eller særskilte Videnskabers System til Forskjel
fra den rene Videnskab eller Logiken, efterdi de ere den rene Videnskabs
System i Forhold til Naturen og Aanden.« Hertil svarer saa følgende Schema:
Overgangen fra Logik til Naturvidenskab dannes af Mathematiken paa Grund af
Rummet og Tiden. Naturvidenskaben falder i tre Dele: a) Mechanik med
Hensyn til Materie og Bevægelse. b) Det Uorganiskes Physik med Hensyn paa
Lys, Magnetisme, Elektricitet, de physiske Legemer og den chemiske Proces.
c) Det Organiskes Physik, hvorunder da igjen hører Geologien, der behandler
Jorddannelsen; Botanik og Plantephysiologi, hvis Gjenstand er den vegetabile
Natur; Physiologi, comparativ Anatomi, Zoologi, der behandle den sunde,
animale Natur, ligesom Medicinen den syge. Ogsaa Aandsvidenskaben falder i
tre Afsnit: a) Aanden i sit Begreb, Anthropologi og Psychologi (Sproget
Philologi). b) den praktiske Aand, der yttrer sig i Daad og Handling,
hvorunder Retslæren, Statsvidenskaben, Historien. c) Aandens Fuldendelse i
Kunst, Religion og Videnskab, hvortil svarer Esthetik, Religionslære,
Philosophi.

Her behøves ikke vidtgaaende Undersøgelser over den Maade, hvorpaa
Naturvidenskaben i Systemet bliver deduceret af Logiken, og
Aandsvidenskaben igjen af Naturvidenskaben: en eneste Udtalelse af
Forfatteren er nok til at opklare den hele Misforstaaelse. »Encyklopædien«,
siger Hegel, »er egenlig en Fremstilling af Philosophiens almindelige
Indhold; thi hvad der i Videnskaberne er grundet paa Fornuft, afhænger af
Philosophien, hvad der derimod er i dem af det, der beroer paa vilkaarlige
og udvortes Bestemmelser, eller, som det hedder, er positivt og statutarisk,
ligger udenfor den.« Hvor rimelig og fornuftig en saadan Erklæring end ved
første Øiekast kan synes at være, saa viser den sig dog, ved nærmere
Eftersyn, at være baade tvetydig og misledende. At den er tvetydig, falder
i Øinene, naar vi sammenligne den med Positivismens Grundtanke. Comte og
Hegel synes jo i dette Stykke at være fuldkommen enige, og dog hvilken
Forskjel! Hos Comte er det i Fornuftens Navn Mathematiken, som baade er en
Videnskab for sig og Grundlag for alle Videnskaber, hos Hegel er det
Logiken. Hvad det Absolutes Philosophi i Logiken sætter som Indhold af
absolut Viden, maa Positivismen ifølge sit Princip forkaste som absolut
Nonsens, fra begge Sider af Hensyn til Fornuften og det i Videnskaberne
Rationelle. Hvor vildledende Udtalelsen er, indsees klart, naar vi gaae ind
paa det Concrete.

At enhver særegen Videnskab maa have sit særegne, af Stof og Emne betingede
Udgangspunkt, følger af Sagens Natur. En immanent Logik kan ved Abstraction
begynde med »den rene Væren«, men hverken Astronomi eller Mineralogi er i
Stand til at begynde med den rene Væren. Paa Udgangspunktet sees det strax,
hvorvidt en Videnskab er empirisk iagttagende eller logisk abstraherende.
Enhver forskende, iagttagende Videnskab maa nødvendig have sin egen
Methode, sit eget Maal. Vistnok have alle Videnskaber et høieste fælles
Maal, nemlig Sandhedens Erkjendelse; men den virkelige Sandhed er ikke i
det blot Ideelle og Almene, den er i det Reelle, i den concrete
Virkelighed. Der er i Virkeligheden mange Udgangspunkter og mange Veie
henimod det sidste Maal. Enhver særegen Videnskab maa desaarsag forfølge
sit Maal og vogte sig for at blive opløst i et abstract Fællesmaal.
Naturvidenskabens Ideal er empirisk-forskende Naturvidenskab, ikke
Naturphilosophi; Historiens Ideal er historisk, ikke philosophisk. Alle
Videnskaber ere villige til at tjene og fremhjælpe hverandre gjensidig; de
ere villige til at tjene, berige og fremhjælpe Philosophien paa bedste
Maade; men udmunde i Philosophien som Floder i Havet, gjøre sig til blotte
Hjælpevidenskaber, »Haandlangere« under Philosophiens Formynderskab ville
de ikke, kunne de ikke, og det af gode teleologiske Grunde.

\section{Videnskabernes Inddeling}
\label{sec:10}
% pp. 90--97

\noindent Enhver Videnskab har sin egen Undersøgelseskreds, sin egen
Methode og sit eget Formaal; dens høieste Opgave er at naae sit særegne
Maal, dens største Fuldkommenhed at være ret videnskabelig. Det er i denne
Henseende umuligt at indføre enten Kastevæsen eller Rangforordning i
Videnskabernes Republik. Spørger man: hvilke ere de vigtigste Videnskaber?
da ville de Kyndige vel i Almindelighed svare, at de alle ere lige
vigtige; men spørge vi hver enkelt Videnskabsdyrker, hvilken Videnskab han
for sit Vedkommende maa ansee for den vigtigste, vil Mathematikeren
upaatvivlelig svare: Mathematiken! Botanikeren: Botaniken! Numismatikeren:
Numismatiken! og Enhver af dem med lige Ret. Naar der desuagtet gjøres
Forskjel mellem høiere og lavere, mere og mindre omfattende Videnskaber,
mellem Hovedfag og Bifag o.s.v., da har Adskillelsen ikke sin Grund i en
høiere eller lavere Grad af Videnskabelighed, thi uvidenskabelige eller
halvvidenskabelige Discipliner kunne ikke regnes med til Videnskaberne.

Inddelingen i høiere og lavere, over- og underordnede Videnskaber grunder
sig væsenlig paa et Totalitetshensyn, idet Noget maa være centralt, medens
Andet er peripherisk. Videnskabernes Organisme maa have et Tyngdepunkt.
Men Tyngdepunktet flytter sig i Tidernes Løb: Oldtiden satte det i
Philosophien, Middelalderen i Theologien, Nutiden finder det i
Naturvidenskaben. Desuagtet er Sagen betænkelig. At gjøre en enkelt
toneangivende Videnskab til den høieste og herskende er vel omtrent lige
saa misligt, som at ville tillægge Kuglens Centralatom Kuglens virkelige
Tyngde og veie det for sig, til Beviis for, at hele Tyngden er concentreret
i Centrum. Det Høieste, Herskende i Videnskabernes Kreds er ikke en
Videnskab for sig, det er Videnskabernes Aand, det Videnskabelige selv, der
besjæler dem alle. Naar dette Høieste, dette Overordnede skal træde frem i
Form af en særegen Videnskab, en Videnskab for sig, maa det nødvendig faae
de øvrige Videnskaber udenfor sig, og kan ikke være det Altgjennemtrængende,
Altbesjælende.

Hvilke Fortrin og Fuldkommenheder maatte ikke den Videnskab besidde, der i
Sandhed skulde egne sig til at være ubetinget eneherskende og forvandle
Videnskabernes Republik til et absolut Monarchi? Det maatte jo være ikke
blot den meest omfattende, men tillige den allerconcreteste, ikke blot den
aandsfyldigste, men tillige den allervidenskabeligste af alle. Men hvilken
Videnskab tør vel paatage sig at fyldestgjøre disse Betingelser?
Philosophien kan det ikke være; thi Philosophien veed med sig selv, at der
i den rene Væsenserkjendelse altid maa gaae Virkelighedselementer til
Grunde; den er sig klart bevidst, at de Realmomenter, den assimilerer i
sit ideelle Indhold, kjøbes og betales med Tab af Realitet, saa at dens
concreteste Bestemmelser, selv naar den gjør sit Yderste, maae betinges af
Abstractioner og falde indenfor det Abstracte. Naturvidenskaben kan det
ikke være; thi Naturvidenskaben i dens Almindelighed er en Abstraction. I
Realiteten bestaaer Naturvidenskaben i mange Videnskaber. Den aandsfyldigste
og allervidenskabeligste af disse Videnskaber maatte da ved sine
Iagttagelser, Beregninger og Inductioner kunne belære de øvrige Videnskaber
om, hvad Natur oprindelig er, og hvorledes det principielle Forhold mellem
Tingene i Naturen og Naturen selv nødvendig maa bestemmes. Men Naturlæren,
i hvilket Hovedfag vi end møde den, er altfor ædruelig stræng, altfor
sikker paa sine Opgaver og sin Methode til at forvexle sig med Metaphysik
og indlade sig med et Problem, der ligger uden for dens Omraade. Hvor er
saa den høieste, styrende Videnskab? Det skulde dog vel aldrig være
Theologien! Ja, sæt at Theologien med sin urandsagelige Aandsfylde, opvakt
af middelalderlig Slummer og dogmatisk gjenfødt til nyt Liv i Speculationens
Daab, satte sig paa Kongestolen og begyndte at synge homerisk: οὐκ ἀγαθὸν
πολυκοιρανίη· εἷς κοίρανος ἔστω! det vilde være et oplivende Syn. Samtlige
Videnskaber vilde sikkerlig bøie sig dybt for den Benaadede; dog, vel at
mærke, paa den ufravigelige Betingelse, at den, istedenfor at synge, i
Gjerning beviste, at den blandt alle Videnskaber virkelig er --- den
allervidenskabeligste.

Den positivistiske Inddeling er ikke tilfredsstillende. Vel er Forskjellen
mellem fundamentale og afledede Videnskaber af Betydning, for saa vidt
hine afgive Grundlaget for disse; men allerede den Omstændighed, at her
foruden Mathematiken, det fælles Grundlag, behøves endnu fem andre
Grundlag, viser, hvor misligt det staaer sig med Deductionen. Skulle de fem
særlige Grundlag afledes af det fælles første, saa maa der jo skjelnes
mellem det oprindelige og de afledede Grundlag; men antager man afledede
Grundlag, saa kommer nok »Hierarchiet« til at bestaae af mere end sex
Videnskaber. Hvori bestaaer Deductionen? Det Afledede maa jo have sin Grund
og sine almene Betingelser i det Oprindelige, maa vise sig som det
Oprindeliges egen, ved synthetiske Tilsætninger modificerede, Udvikling;
saaledes ere Formlerne for de forskjellige Keglesnit modificerede
Udviklinger af deres Fundamentalformel, saaledes ere Krystalformerne i et
System afledede Former af Systemets Grundform. Hvorledes vil man nu kunne
begrunde den Paastand, at Chemi, Physiologi og Sociologi kun ere
Modificationer af Mathematiken? Man kan ikke engang aflede Astronomien af
Mathematiken. Vistnok maa det indrømmes, at hverken Astronomi eller Physik
vil kunne faae sin videnskabelige Begrundelse og fuldstændige Udvikling
uden Mathematikens Hjælp; men dette er jo noget ganske Andet. At Mathematik
er en uundværlig Hjælpevidenskab for Astronomi og Physik, er uimodsigeligt;
men i disse empiriske Videnskaber ere de mathematiske Sætninger, væsenlig
seet, kun Laanesætninger. I hvilket Omfang de end anvendes, og hvor
systematisk de end tilegnes, saa udgjøre de dog ikke Grundvolden selv. Det
er Iagttagelsen af de virkelige Phænomener, de virkelige Stjerner, de
virkelige Legemers physiske Egenskaber, og de virkelig stedfindende
Forhold, der give de nævnte Videnskaber deres eiendommelige Charakteer og
virkelige Grundlag. Kun paa høist ufuldkommen Maade kunne Videnskaberne
afledes af hverandre indbyrdes.

Det almene Grundlag, \textit{fundamentum divisionis}, maa søges i Videns
almindelige Forudsætninger: den Vidende, Videns Gjenstand og
Bevidsthedseenhed af begge. Principet vil da dele sig i forskjellige
Principer. Lægges Eftertrykket paa Videns Eenhed med den Vidende, saa faae
vi det Aprioriske; lægges Eftertrykket paa Videns Forhold til existerende
Gjenstande, altsaa paa Sandsning og Sandseindtryk, saa faae vi det
Empiriske. Videnskaberne maae da med Hensyn hertil dele sig i aprioriske og
empiriske. Det Aprioriske antager igjen forskjellig Charakteer, eftersom
Forestillingerne enten opløse sig i almene Tankebestemmelser eller afgive
exacte Forholdsled eller berige Anskuelsen med imaginære Masser og Kræfter;
hertil svare: Logik, Mathematik og rationel Mechanik. Ligeledes vil det
Empiriske antage forskjellig Charakteer; det vil modificeres eiendommelig,
ikke alene derved, at Gjenstandene frembyde forskjellige Hovedsider,
hvorfra deres lovbundne Vexelvirkning bliver at betragte, men principielt
derved, at den virkelige Loverkjendelse deels kan acqviescere i Forholdet
mellem givne Forudsætninger og deraf fremgaaende Resultater, deels maa
underforstaae en Selvvirksomhed, som igjennem Lovenes Opfyldelse stræber
mod et bestemt Maal. Med Hensyn hertil inddele vi de empiriske Videnskaber
i realempiriske --- mechanisk og chemisk Physik, Chemi og Geologi --- og
teleologisk-empiriske --- Botanik og Zoologi, Anthropologi og Historie.

Grundforholdet mellem det Aprioriske og det Empiriske viser i sidste
Instans hen til et Grundforhold mellem Nødvendighed og Frihed, der maa
behandles i Aandslæren. Dersom Friheden selv faldt udenfor Nødvendigheden,
vilde ethvert Erkjendelsessystem, der lægger an paa at hævde og belyse
Frihedens Ideal, uden videre være at erklære for uvidenskabeligt, thi
Nødvendigheden er, saa at sige, den eneste faste Grund, hvorpaa en
videnskabelig Bygning kan hvile. Men af det Aprioriske i Forholdet mellem
Videns Indhold og det vidende Selv maa det allerede kunne sees, at Friheden
boer i Nødvendigheden og arbeider sig frem af Grunden. At al Viden
oprindelig beroer paa en Selvvirken, for saa vidt den beroer paa Sandsen og
Tænken, er allerede viist i vort Foregaaende; at denne Selvvirken
ingenlunde afbrydes af det Nødvendige, det Objective, men tvertimod
betinges deraf, er ligeledes fremhævet. Hvis det Fuldkomne kunde naaes,
vilde det ende med, at den vidende Subjectivitet under Objectiveringen af
alt Andet i Virkeligheden objectiverede sig selv og sit eget Væsens
Fordringer: vi vilde da have en med Nødvendigheden identisk Frihed.

Men paa Ufuldkommenhedens Trin, paa Endelighedens Vilkaar maa der, som vi
have seet, forblive en Differens mellem den sandsende, objectiverende
Subjectivitet og Realiteten af de Objecter, der svare til dens
Objectivering. Heri ligger da Grunden til, at Videnskaben med sine
Nødvendighedsgrunde kun kan erkjende Frihedens Relativitet, og maa opløse
sig selv efterhaanden som den Vidende fordyber sig i den ideale Frihed.
Aandslæren skal da vise, hvorledes den objective Videnskab Trin for Trin
--- gjennem Esthetik, Ethik og Religionsphilosophi --- fuldbyrder sin egen
Opløsning i Idealet selv. Hvad endelig Forholdet mellem Fleerhed og Eenhed
angaaer, er der Grund til at dele Videnskaberne i Fagvidenskaber og
Aandsvidenskab, dog saaledes, at Inddelingsgrunden ikke fastsættes
kritisk, men opfattes dialektisk. Kritisk ligefrem gjælder Inddelingen
ikke; thi alle Videnskaber ere Aandsvidenskaber og maae gaae op i Eenheden,
saa sandt den videnskabelige Aand gjennemtrænger dem alle. Omvendt er
Aandsvidenskabens Eenhed kun en Abstraction, naar den ikke adskiller sig i
en Fleerhed af Aandsvidenskaber. Inddelingen motiveres derved, at den
fagvidenskabelige Grændse er kritisk, den aandsvidenskabelige dialektisk.
Enhver Videnskab, Philosophien, »Videnskabernes Videnskab«, ikke undtagen,
maa saa skarpt som muligt begrændse sit Omfang, følge sin Methode og
fastholde sit Maal; netop derved bliver den en Fagvidenskab, og den ene
Fagvidenskab kan ikke dialektiseres ud af den anden. Med Aanden og de
aandelige Momenter forholder det sig omvendt. Da Aandsudvikling er alle
Videnskabers fælles Maal, og Methoden afhænger af Maalet, saa maa alt det
særlige Indhold, Aandsvidenskaben optager fra Fagvidenskaberne i sine
ideelle Concretioner, med det Samme gaae ind under Eenheden. Selv idet
Aandsvidenskaben adskiller sig i særskilte Discipliner, vil det dog vise
sig, at den ene Disciplin kan forvandles til et Moment, der optages i den
anden, saa at Grændsen bliver dialektisk, og Totalitetseenheden afgjørende.


%%% ANDEN DEEL: FAGVIDENSKABER %%%

\part{Fagvidenskaber}

\chapter{Aprioriske Videnskaber}

\section{Logik}
\label{sec:11}
% pp. 98--117

\noindent Vi kjende ingen Tanke, som ikke er et Produkt af Tænkning, og
ingen Tænkning, som ikke er en Function af et tænkende Væsen. Da man
imidlertid kan finde Tanker udtalte, ikke blot i Ord, men i reelle
Frembringelser, saa bliver det muligt at sætte sig ind i et Tankesystem
uden at spørge, hvo Tænkeren er. Kun dette staaer fast, at Tanker ikke
kunne tænke sig selv, men maa forudsætte en tænkende Subjectivitet.
Logiken, Læren om Tænkningen, kan med Hensyn hertil fremstilles fra tre
Synspunkter. Man kan holde sig til Mennesket som det tænkende Subject og
af Tankehandlingerne abstrahere Love og Former; man kan uden at agte paa
det Selv, hvorfra Tankerne udgaae, analysere de almeennødvendige
Tankeformer og udvikle dem i et System; man kan opløse den endelige
Tænkning i det Uendelige saaledes, at den tænkende Selveenhed skinner
igjennem i tænkende Ideeformer. Herved bestemmes Forholdet mellem formel
eller subjectiv Logik, immanent eller objectiv Logik og Grundideernes
Logik.

Den formelle Logik deler sig i Læren om Begreb, Dom og Slutning.

Begrebet dannes paa menneskelig Viis, idet Tænkningen, som oprindelig, om
end ubevidst, er med under Sandsningen, særlig hæver sig frem og ved
Sammenligning og Abstraction opløser de umiddelbare Sandsebilleder i et
svævende Almeenbillede, der maa fastholdes i Ordet. Almeenbilledet af et
Træ f.\ Ex.\ kan ikke tegnes; de ubestemte Omrids vise, at det befinder
sig paa Overgangen fra Forestilling til Begreb. Indbegrebet af de til et
Begreb hørende Bestemmelser udgjøre Begrebets Indhold, Antallet af de
Gjenstande, hvorpaa det finder Anvendelse, dets Omfang. Ved at mærke sig
det omvendte Forhold mellem Indhold og Omfang kommer man til en
Adskillelse mellem over- og underordnede Begreber, Slægts- og
Artsbegreber, idet Slægtsbegreberne sees at have større Omfang men mindre
Indhold, medens det Omvendte gjælder om Artsbegreberne.

Dommen dannes paa menneskelig Viis, og Begrebet finder sit logiske Udtryk
i Sproget: Ord forbindes til Sætninger, og i Sætninger udtales Dommen.
Dommen maa da formelig angive: Subjectet, den Forestilling, der skal
bestemmes, Prædicatet, den bestemmende Forestilling, og Copula, det
bestemmende Forbindelsesord. Af Subjectets ved Copula bestemte Forbindelse
med Prædicatet følger, at Dommen kan adskilles i en Fleerhed af logisk
forskjellige Domme; og Dommene inddeles efter Qvalitet: i bekræftende,
benægtende og uendelige; efter Qvantitet: i singulære, particulære og
generelle; efter Relation: i kategoriske, hypothetiske og disjunctive;
efter Modalitet: i assertoriske, problematiske og apodiktiske. Fremgaaer
Prædicatet ved formel Udvikling af Subjectet, saa er Dommen analytisk; kan
Prædicatet kun ifølge en særegen, være sig apriorisk eller empirisk,
Anskuelse, føies til Subjectet, saa er Dommen synthetisk.

Slutningen dannes paa menneskelig Viis; thi Dommen maa begrundes: den
Tænkningsact, hvorved en Doms Gyldighed udledes af een eller flere andre
Domme, er en logisk Slutning. Har Slutningen kun een Præmis kaldes den
umiddelbar, har den flere Præmisser da middelbar. De middelbare Slutninger
inddeles i: de kategoriske, der betinges af, at det underordnede Begrebs
Kjendemærke (\textit{nota notæ}) subsumeres under det overordnede; de
hypothetiske, der beroe paa det Betingedes Afhængighed af sin Betingelse
(\textit{modus ponens} og \textit{modus tollens}); og de disjunctive, som
grunde sig paa, at Leddene af et Hele gjensidig forholde sig saaledes, at,
naar eet udtrykkelig er sat, saa ere de øvrige udelukkede (\textit{modus
ponendo tollens}), og at, naar de øvrige ere udelukkede, saa maa det
tilbageblevne sættes (\textit{modus tollendo ponens}). Den tænkende
Subjectivitets Indskrænkning er saavel i Begrebsdannelse som i Dom og
Slutning kjendelig paa, at Tænken og Væren overalt falde ud fra hinanden.

Den immanente Logik, grundlagt af Hegel, abstraherer fra den tænkende
Subjectivitet, identificerer de reale Værensbestemmelser med de almene
Tankebestemmelser og udvikler de concretere Begreber af de abstractere saa
dialektisk, at Tankebevægelsen fra »den rene Væren« til »den absolute
Idee« faaer Udseende af en continuerlig fremadskridende Sagbevægelse.
Systemet deler sig i Læren om Væren, Væsen og Begreb.

\emph{Væren.} Den almene Væren er Eet med den almene Tænken; men Tankens og
Værens Eenhed er i sin abstractere Form indholdsløs: den rene Væren er det
samme som det rene Intet. Og dog er Væren, det Positive, forskjellig fra
Intet, det Negative. At Væren baade er og ikke er Intet, vil sige, at Væren
uophørlig gaaer over i Intet, og Intet i Væren: denne Overgang er Vorden. I
Vorden er Intet en Grændse for Væren, og Væren for Intet. Den ved sit Intet
bestemte Væren er et Værende, et Noget, og det ved Væren bestemte Intet, et
Andet. Noget er forskjelligt fra Andet, thi Andet er dets Grændse; og dog
det Samme som Andet, thi Andet er selv Noget, og Noget igjen Andet; Noget
gaaer ifølge sin Grændse over i Andet, og Andet i Noget: denne Overgang er
Forandring. I Forandringen er Noget Eet med sit Andet og for saa vidt
uforanderligt. Det er altsaa det Uforanderlige, der forandres. Al den Uro,
som sees i Væren og Tilværen, er forstaaelig deraf, at Grændser som: Intet
mod Væren, Andet mod Noget, Flere mod Eet, Foranderligt mod Uforanderligt,
ere for Tanken rene Negationer, medens de dog i Forestillingen vise sig som
umiddelbare Led. Denne Modsigelse hæves, idet Negationen og det Umiddelbare
skinne igjennem hinanden i Reflexionen.

\emph{Væsen.} Den reflecterede Væren er paa een Gang Væsen og Skin. Det
Uforanderlige, der reflecterer sig i det Foranderlige, er Væsen; det
Foranderlige, Flygtige, der reflecterer det Uforanderlige, er Skin:
Vexelskin af Skin og Væsen er Reflexion. I Reflexionen har Identiteten
Forskjel, og Forskjellen Identitet. I sin Identitet med Væsnet er Skinnet
Phænomen, i sin Forskjel fra Phænomenet er Væsnet Grund. Phænomenet gaaer
til Grunde; men netop ved at faae Væsnet til Grund bliver det begrundet og
sat som Existens. Den i sig reflecterede Existens er et Existerende, en
Ting; men Tingen kan kun reflecteres i sig ved at reflectere sin Existens i
andre existerende Ting: de til disse Reflexioner svarende Bestemmelser ere
Egenskaber. Under den vexelsidige Reflecteren opløse Tingene sig i
Egenskaber, og Alt bliver ubestandigt. Men det Ubestandige forudsætter et
Bestaaende, der ikke kan opløses, fordi dets Bestaaen uendelig bekræfter
sig i det Ubestandiges Vexlen; dette Bestaaende er Substansen, og de i
Ubestandighed vexlende Fremtrædelsesmaader ere Accidenserne. Forholdet
mellem Substansen og dens Accidenser er et Causalitetsforhold: Substansen
er Aarsag, den oprindelige Sag; Accidenserne ere Virkninger. Men Aarsagen
har sit substantielle Udtryk i Virkningen: Accidenserne selv ere endelige
Substanser. Virkningen er saaledes en Virkning af den ene Substans paa den
anden, hvilket vil sige: en Aarsags Virkning paa en anden Aarsag, d.\ e.\
Vexelvirkning. I Vexelvirkningen er Virkningen Aarsag og Aarsagen Virkning;
derved at alle i Væsnet indeholdte Virkelighedsmomenter under alsidig
Vexelvirkning ere udviklede til fuld Totalitet, er det i Totaliteten
udviklede Væsen Begreb.

\emph{Begreb.} Den sig begribende Totalitet er i sin Identitet med sig selv
det Almene, i sin Forskjel fra sig selv det Særskilte, i sin totaliserede
Bestemthed ved sig selv det Enkelte. Begrebets Selvudvikling er dets
Grunddeling og Mediation, d.v.s.\ Dommen og Slutningen. Formedelst
Grunddelingen er Totaliteten i sin første umiddelbare Eenhed det Enkelte,
sat som Subjectet, og Totaliteten i den reflecterede Form, det Almene, sat
som Prædicatet: Subjectet er, hvad Prædicatet er, det Enkelte er det
Almene. Men derved, at det Enkelte begribes som det Almene, og det Almene
som det Enkelte, blive Begrebsadskillelsens modsatte Led gjensidig begrebne
i indbyrdes Vexelbestemmelse som det Særskilte. Grunddelingens Yderled,
Extremerne, have altsaa under Begrebets Selvudvikling afsat et Mellemled,
et Medium, hvori de kunne medieres. Denne Mediation er Slutningen.
Begrebets Selvmediation gjennem Dom og Slutning er dets Subjectivitet, den
ved Mediationen gjenvundne Umiddelbarhed er dets Objectivitet. Som
umiddelbar tilværende Totalitet er Begrebet Object, og ifølge
Grunddelingen, en uendelig Mangfoldighed af tilværende Objecter. Objecterne
i deres udvortes, indifferente Gjensidighed ere mechaniske, virkende paa
hverandre ved Tryk og Stød. Den herved indtrædende Adskillelse mellem
Bevægelse og Hvile er en relativ Tilsyneladelse af den absolute Mechanik,
hvori Bevægelsen er Hvile, og Hvilen Bevægelse. --- Under de mechaniske
Objecters Indifferens skjuler sig en Differens, der kommer tilsyne i den
chemiske Tiltrækning. Objecternes Differens medfører en Spænding, der
udlignes i det chemiske Produkt. Men Produktet, hvori Differensen er
neutraliseret, maa som det Neutrale igjen opløse sig i Differens, og
Processen begynde forfra. I den chemiske Proces gaaer Begrebet imidlertid
frem af sine egne Forudsætninger, ophæver Forudsætningerne i deres
Resultat, for igjen at ophæve Resultatet i Forudsætningerne. Det saaledes
udviklede Begreb er nu Herre over Stoffet: det i Forudsætningerne
anticiperede Resultat er Formaal, Forudsætningerne med deres
Mellembestemmelser ere Midler, Bevægelsen er teleologisk. Det teleologiske,
sig selv i sine Elementer begribende, Begreb er en Form, der giver sig
selv Indhold, et Indhold, der giver sig selv Form, er i absolut Eenhed af
Form og Indhold, af Subjectivt og Objectivt, af Formaal og Midler, med eet
Ord: den absolute Idee.

Grundideernes Logik lægger an paa en videre Udvikling og nærmere
Bestemmelse af den absolute Idees Indhold; de særlige Grundformer,
hvorunder Indholdet opfattes, ere Viden, Magt og Sandhed.

\subsubsection*{a) Videns Idee.} En Tanke, som Ingen har tænkt, et Begreb,
som Ingen begriber, en Idee, som Ingen fatter, en Viden, som Ingen veed,
ere absolut meningsløse Indbildninger. Forsøger man at tilsløre det
Meningsløse under Talemaader som disse: Værens Tanker tænke sig selv, det
objective Begreb begriber sig selv, Ideen bevæger sig selv med absolut
Viden, da skinner Modsigelsen igjennem; thi dette »sig selv« og »sig selv«
viser paa alle Punkter hen til en sig selv reflecterende Eenhed, en Eenhed,
som i alle Tanker tænker sig selv, i alle Begreber begriber sig selv, en
Selveenhed, som i uendelig Viden fatter, bærer og bevæger Ideen med hele
dens Indhold. Hvorfor en saadan absolut vidende Subjectivitet nødvendig maa
være, og hvorfor det ikke kan være den menneskelige Subjectivitet, er
logisk kjendeligt paa Reflexionens, Dommens og Slutningens Opløsning i
Videns Idee.

\emph{α) Reflexionens Opløsning.} For den endelige Subjectivitet er det
Sandselige et umiddelbart Andet for Sandsningen, ligesom det Tænkelige er
et Andet for Tænkningen. Opfattes nu dette sandselige og tænkelige Andet
som udvortes Grændse, som Skranke, saa faaer den sandsende-tænkende
Subjectivitet en Verden af sandselig-tænkelige Gjenstande udenfor sig.
Blændet af dette Umiddelbarhedens Udenfor, antager saa den endelige Viden,
at det Sandselige forud staaer fuldt færdigt og blot venter paa at blive
sandset. Denne Venten betyder da ikke nogen Mangel eller væsenligt Savn.
Kommer Sandsningen, saa er det godt; det Sandselige tager ingen Skade
deraf, lider ikke Noget derunder: det tillader gjerne den Sandsende at
anbringe sin Opfattelse, thi det er jo til at sandse, det er jo sandseligt.
Men kommer Sandsningen ikke, saa er det ogsaa godt; det Sandselige har
ligefuldt sine fuldstændige Værensbestemmelser, sine udprægede
Virkelighedsformer i sig selv. Som det gaaer med det i Tingen Sandselige,
gaaer det ogsaa med det Tænkelige. Derfor drømmer man om objective
Naturtanker, som ere uden at blive tænkte: et talende Beviis paa den
endelige Videns Indskrænkning. Hvad er herved at gjøre? Kritiken mærker
Uraad, men kommer ikke ud over det Endelige.

Den endelige Reflexions Opløsning i Videns Idee er Modsigelsens Løsning.
Sandseligt og Sandsning, Tænkeligt og Tænkning ere correlate Bestemmelser.
Hvad er vel det Sandselige for sig? Et Ubestemt, der maa lade sig bestemme
--- ved Sandsning. Holdes det Sandselige fast i al sin Ubestemthed, da er
det ikke et Tilværende, thi ethvert Tilværende er et bestemt Værende: det
endnu Ubestemte er ikke til i Virkelighed, det er kun en Mulighed. Det
Sandselige er en Mulighed, nemlig for Sandsningen, en Sandsningens
Mulighed, ligesom det Tænkelige er en Tænkningens Mulighed. Ogsaa det
Omvendte gjælder, at Sandsning for sig kun er en Mulighed for det
Sandselige, og Tænkning for sig kun en Mulighed for det Tænkelige. Hvor
intet Sandseligt er, kan der ingen Sandsning være, og hvor intet Tænkeligt
er, kan der ingen Tænkning være. At sandse et Intet er umuligt, og naar
Tænkningen tænker sit Intet, sin Negation, er dette Intet selv formedelst
Ordet et Tænkeligt og faaer Tilværelse for Tanken. I Videns Idee er den
endelige Adskillelse af Værensbestemmelser og Tankebestemmelser opløst,
saasandt der er en Viden, hvis Reflexion er uendelig.

\emph{β) Dommens Opløsning.} Fra den endelige Subjectivitets Synspunkt vil
Subjectet i en Dom trods alle Mediationer naae til fuld Identitet med sit
Prædicat. At Subjectet kategorisk er sit Prædicat, faaer sit fyldigste
Udtryk i den concrete Form: Individet er sin Art. Alle til Individet hørende
Bestemmelser ere sandselige, alle til Arten som Art svarende Mærker ere
tænkelige: Individet sandses, Arten tænkes; saaledes udtrykker Dommen paa
een Gang Forskjellen mellem Sandseligt og Tænkeligt, og tillige Eenheden.
Den endelige, af sit Udenfor indskrænkede Viden, kan imidlertid ikke komme
paa det Rene med sig selv angaaende Forholdet mellem Subject og Prædicat.
Det seer ud, som var Individet, det Sandselige, egenlig det Første og
Oprindelige, Arten, det Tænkelige, derimod det Afledede; thi den sandsende
og dømmende Subjectivitet maa jo gaae ud fra en sandselig Opfattelse af
Individerne og af Individerne abstrahere Arten. Saaledes er Individernes
Tilværen udenfor, det abstraherede Artsbegreb derimod indenfor Bevidstheden.
Sandsningen (den sandsende Bevidsthed) paastaaer, at det er Individerne, der
bestemme Arten, medens Tænkningen fra sin Side fastholder, at det er Arten,
der bestemmer Individerne; Sandsningen paastaaer, at Individet er det
Oprindelige, Arten det Afledede; Tænkningen paastaaer, at Arten er det
Oprindelige, Individet det Afledede: her er Spænding og Strid mellem
Sandsning og Tænkning. Den endelige Doms Opløsning er Modsigelsens Løsning.
De to Factorer, den sandselige og den tænkelige, ere nemlig begge optagne
som Momenter i den uendelige Selvhedsact, hvorved den absolut vidende
Subjectivitet oprindelig sætter Dommenes Indhold identisk med Dommens Form.

\emph{γ) Slutningens Opløsning.} Kom det blot an paa at opstille et
Slutningschema som dette: Individet, E, er Slægten, A, formedelst Arten, S:
E\,--\,S\,--\,A; Arten er Slægten formedelst Individet (S\,--\,E\,--\,A),
Individet er Arten formedelst Slægten (E\,--\,A\,--\,S); da vilde
Idee-Erkjendelsen ikke alene høre til de letteste, men ogsaa til de meest
overfladiske af alle videnskabelige Erkjendelser. Den væsenlige Opgave
kommer først tilsyne, naar man besinder sig paa, hvad et saadant
Slutningschema egenlig betyder. Det betyder nemlig, at af Slutningens tre
Led: Individ, Art og Slægt, er intet det umiddelbart første, og intet det
sidste, men at alle tre maae have samme Oprindelighed. Det Indbegreb af
Sandsebestemmelser, der er samlet i Individets Enkelthed, E, er forud
bestemt ved den Totalitet af Tankebestemmelser, der indeholdes i Artens
Særskilthed, S, men dette Tankeindhold er igjen en ejendommelig Concretion
af Slægtsbegrebets almeen-concrete Indhold. Saa skulde Dannelsen af
Slægtsbegrebet altsaa være det ubetingede Første. Men da dette Begreb i
Ideen maa indeholde alle Bestemmelser for Arterne, saa maa Dannelsen af
Slægtsbegrebet nødvendig forudsætte Artsbegreberne. Saa skulde Artsbegrebet
da være det absolut Første; men det er umuligt, eftersom ethvert
Artsbegreb deels forudsætter sit Slægtsbegreb, deels Individerne, hvis Art
det er. Her er Modsigelsen. Individet kan ikke dannes, før Arten og Slægten
ere dannede, thi seet i Ideen, er Individet den umiddelbare Slutningseenhed
af Art og Slægt. Arten kan ikke dannes, før Individet og Slægten ere
dannede; thi seet i Ideen, er Arten den reflecterede Slutningseenhed af
Slægt og Individ, og maa udtrykkelig have de to Extremer til sine
Præmisser. Saa er det dog vel indlysende, at det Slutningssystem, der skal
udtrykke Ideens Indhold i Videns Form, maa udspringe af en eneste
trilogisk Slutningsact, og at en saadan Slutningsact alene kan fuldbringes
af en absolut Vidende.

\subsubsection*{b) Magtens Idee.} At den immanente Logik ikke formaaer at
begribeliggjøre Identiteten af Tænken og Væren, fordi dens absolute Idee
mangler en absolut vidende Subjectivitet, maa vel være klart af vort
Foregaaende; men at den, naar dens almeenlogiske Væren skal omsættes i
virkelig Tilværelse, endnu mindre formaaer det, fordi dens absolute Idee
mangler en absolut mægtig Subjectivitet, vil det Følgende vise. Medens
Positivismen umiddelbart holder sig til Kjendsgjerninger og bygger paa
Realiteter, er Realitetsproblemet end ikke blevet til i det Absolutes
Philosophi. Dette Problem indeholder saa haarde Modsætninger, ja saa
skjærende Modsigelser, at den blotte Tænkning, den være nok saa dialektisk,
umulig kan løse det i Tænkningens Form.

\emph{α) Ting og Begreb.} Modsigelsen er, at Ting og Begreb baade ere det
Samme og dog ikke det Samme. De ere det Samme; thi der kan ikke tænkes
nogen ved Tingen afgjørende Egenhed, som ikke skulde høre med til dens
Begreb, og der kan i Tingens Begreb ikke findes nogen Bestemmelse, som ikke
væsenlig skulde høre med til Tingen; og dog ere de ikke det Samme, thi
Tingen er sandselig, men Begrebet usandseligt. De ere det Samme, thi Begreb
og Ting have samme Indhold, og dog ere de ikke det Samme, thi Begrebet, det
Almene, har et uendeligt større Omfang end Tingen, det existerende Enkelte:
eet Begreb kan være i mange Ting. De ere det Samme, thi den existerende Ting
er Begrebets Existens, og dog ikke det Samme, thi den existerende Ting er
intet Begreb, og Begrebet ingen existerende Ting. En Steen er intet Begreb,
og Begrebet ingen Steen; man kan knuse en Steen, men ikke Stenens Begreb;
man kan koge en Plante, men ikke Plantens Begreb; man kan ride paa en Hest,
men ikke paa Hestens Begreb. Begrebet kan tænkes og der kan tænkes over
Tingen, men Tingen som existerende Ting kan ikke tænkes; Begrebet kan
udtales i Ord, men Tingens existentielle Bestemthed er uudsigelig: Begrebet
er logisk, Tingen er antilogisk.

\emph{β) Det Logiske og det Antilogiske.} Naar det uden videre antages, at
det umiddelbart er Tingen, der tænkes, og at Ordet nøiagtig udtrykker, hvad
Tingen er, da er dette en Skuffelse, som den logiske Dom foranlediger. Hvad
er nu Dette? Dette er --- Dette! Og dog kan den paapegende Betegnelse
udtrykkelig vise, at Dette ikke er Dette; thi Dette f.\ Ex.\ er et Huus,
men Dette er en Hest, og Dette er igjen et andet Dette. Ved at pege paa
Gjenstanden viser Ordet bort fra Ordet, Tanken bort fra Tanken: dette er
det Antilogiske. Den Umiddelbarhed, som adskiller Ting og Begreb, kan ikke
være logisk; thi om end det Logiske kan vise ud over sig selv, skeer det
dog altid indenfor Logikens Sphære. Umiddelbarheden kan heller ikke være
ulogisk; hvis Gjenstandsbestemmelserne vare ulogiske eller stode i
Modsigelse med det Logiske, maatte Gjenstanden enten være begrebsløs eller
staae i Modsigelse med sit Begreb. Heller ikke kan det i Gjenstanden
Umiddelbare være Noget, der adderes eller udefra føies til det Logiske, som
kunde man adskille Gjenstanden selv fra det i Gjenstanden udtrykte Begreb.
Hvad skulde vel det i Gjenstanden være, som blev tilbage, naar Begrebet gik
bort? Materialet, Materien? Men hvis Materien med alle sine Bestemmelser
falder udenfor Begrebet, hvorledes kan Begrebet da være udtrykt i den
materielle Gjenstand.

Søge vi den logiske Rod, hvoraf det Antilogiske udspringer, saa finde vi
den i Kategorien Magt. Magten i sin Umiddelbarhed er Realitetens Udtryk, og
Udtrykket for Realitetens Umiddelbarhed er igjen det Antilogiske. Grunden
til at vi ansee Gjenstandene for Realiteter og tillægge dem en umiddelbar
objectiv Selvstændighed er jo netop Indtrykket af den Magt, hvormed de hver
især træde os imøde. Hvis Tingenes Verden kun var en Objectivering af en
absolut vidende Subjectivitet, vilde den kunne opløses i et System af
Begreber, af Love; men at Tingene ere Realiteter, vil med andre Ord sige,
at de ere Objecter af en absolut mægtig Objectivering.

\emph{γ) Udslag af Magt.} Den til Objectiveringen svarende Overgang fra
Ideelt til Reelt, fra Begreb til Ting, har ingen Mellembestemmelser og kan
ingen have. Hvorledes skulde disse Mellembestemmelser vel være beskafne?
Logiske kunne de ikke være, thi man kan berige et Begreb med saa mange
logiske Bestemmelser man vil, det bliver som Begreb dog indenfor det
Logiske. Antilogiske kunne de heller ikke være, thi det Antilogiske bliver
jo først til i Overgangen selv. Overgangen maa da være en Magtbevægelse, en
Omsætning af Magtens ideelle Realitet i reel Magt, et Spring, hvorved Intet
overspringes, en Overgang fra Mulighed til Virkelighed, forstaaet saaledes,
at Muligheden ligger i Begrebet, altsaa i det Logiske, Virkeligheden
derimod i den antilogiske Realitet, altsaa ud over det Logiske. De
Betingelser, der ere mæglende ved Overgangen, maae selv høre til det
Virkelige; altsaa: enten det Ideelle eller det Reelle, enten Begreb eller
Realitet; et Mellemværende gives ikke. At ville fortætte Begreberne til
Realiteter eller sublimere Realiteterne til Begreber er et speculativt
Fuskeri, i lige Grad stridende mod Erfaring og sund Logik. Det mellem
Idealitet og Realitet uendelig mæglende Væsen er det i sin Viden absolut
mægtige Selvvæsen, den almene, Alt objectiverende Subjectivitet i Magtens
Idee.

\subsubsection*{c) Sandhedens Idee.} Hvad er Sandhed? I Reflexionens Form
er Sandheden Eenhed af Tænken og Væren: Tænken har sin Sandhed i Væren, og
Væren i Tænken; Usandhed er altsaa Splid mellem Tænken og Væren. I Dommens
Form er Sandheden Identitet af Subject og Prædicat: Gjenstanden, Subjectet,
har sin Sandhed i Begrebet, det fuldt udviklede Prædicat, og omvendt;
Usandheden, Modsigelsen i Dommen, er altsaa Strid mellem Subject og
Prædicat. I Slutningens Form er Sandheden Extremernes Midte: jo
selvstændigere Extremerne fremtræde, desto fyldigere er Midten, og desto
dybere er Sandheden. Den uendelig indholdsrige Midte, den evige Grund,
hvoraf Videns intellectuelle og Magtens reelle Verden fremgaae som
Extremer i det Endelige og sammensluttes til articuleret Eenhed i det
Uendelige, er Sandhedens Idee.

\emph{α) Idealisme.} Idealismens videnskabelige Charakteermærke er, at hele
Sandhedens Indhold gaaer op i Videnskabens Form. Falder Eftertrykket
methodisk paa Begrebernes Oprindelse, altsaa paa den begribende
Subjectivitet, et endeligt Jeg med Skranker (Fichte), saa er Idealismen
subjectiv; er Methoden derimod den, at lade det almeengyldige
Begrebsindhold udfolde sig selv i Begrebsbestemmelsernes Udvikling (Hegel),
saa er Idealismen objectiv. Usandheden i den subjective Idealisme er, at
den ikke formaaer at udvikle Tilværelsens objective Former af de subjective
Jeg-Former; dens Mangel paa Formudvikling er den væsenlige Grund til dens
Mangel paa Indhold. --- For Indholdets Skyld begynder den objective
Idealisme med de allerabstracteste Værensformer: Væren, Vorden, Tilværen
\ldots, integrerer det Abstracte ved stedse concretere Bestemmelser, og
see! det ideelle Indholds objective Selvudfoldning foregaaer som ved et
Trylleri. Men hvad er Sandheden i denne Bevægelse? En uophørlig Tænkningens
Kamp med Usandheden! Det Abstracte er i sin Eensidighed det Usande, men som
Moment i en tilsvarende Concretion har det altid sin Sandhed. Sandheden er
en Overgang fra den eensidige til den alsidige Sandhed; Bevægelsen selv er
Sandheden, eller nøiagtigere: en bestandig Kommen til Sandheden. Og dog er
heller ikke dette Sandheden. Bevægelsen er kun sand, saalænge den fulde
Sandhed ikke er naaet, altsaa kun sand paa Grund af en relativ Usandhed.
--- For at Idealismen kan være absolut, maa Sandheden, den absolute Videns
absolute Indhold, være i den absolut Vidende. Med denne Forudsætning viser
Idealismen ud over sig selv, thi den absolut Vidende er den absolut
Mægtige: den absolute Viden er i Sandhedens Idee Eet med den absolute Magt.

\emph{β) Realisme.} Det Sande og Usande i Realismen sees allertydeligst af
»den positive Philosophi«. Positivismen frembyder i denne Henseende tre
charakteristiske Hovedsider: den forkaster enhver gjennemført
Begrebsudvikling under Navn af Metaphysik; den nedsætter den reale
Causalitet til en blot subjectiv Betingelse for at erkjende Lovene; den
fornægter den ideelle Causalitet og bryder Staven over al Teleologi. ---
Fra Kjendsgjerningernes, de endelige Realiteters Synspunkt, er det ganske
sandt, at Alt er relativt, at absolute Begreber ikke existere; men at det
Relative desaarsag ikke skulde have sin Grund i det Absolute, og at det
ikke skulde ligge i Videnskabens Interesse at faae Begrebsforholdet mellem
det Relative og det Absolute tilstrækkelig belyst, er et vilkaarligt
Postulat, en Usandhed. Naar Begrebet af det Relative gjennemtænkes, saa
viser sig det Absolute, og naar det Absolutes Begreb gjennemtænkes, saa
viser sig det Relative; naar det Endelige analyseres, saa viser sig det
Uendelige, og naar det Uendelige analyseres, viser sig det Endelige.

Lignende Eensidigheder komme tilsyne i den positivistiske Behandling af
Causaliteten. Ved at nedsætte Kategorierne, Aarsag og Virkning, til
nominelle Værdier, subjective Betegnelser for Tingenes Sammenhæng, og
berøve den reelle Causalitet sin objective Sandhed, berøver man tillige
Lovene al objectiv Sandhed. Men har den reelle Causalitet objectiv Sandhed,
saa maa den ideelle Causalitet ogsaa have oprindelig, subjectiv Sandhed. At
enhver Ting virker efter sin Natur, maa dog vel betyde, at den virker i
Medfør af sit Realbegreb; Realbegrebet er altsaa ideel Aarsag. Vel er det
sandt, at Begrebet som Begreb ikke er virkende Aarsag: Tyngdens Begreb
tynger ikke, Ildens Begreb brænder ikke, Stenens Begreb slaaer ingen Hund
ihjel. Men Tingenes antilogiske Existens er et Magtens Udslag af Begrebets
Logik. Hvis ingen Realbegreber vare realiserede i de existerende Ting, vilde
alle objective Muligheder og dermed al Virken og Vexelvirken ophøre. I
begrebsløs Facticitet vilde Tingene existere ved Siden af hverandre,
udelukkende hverandre, men for Resten ikke have det Mindste med hverandre
at gjøre. Det er Realbegreber der bestemme de virkende Aarsager; men
Realbegrebernes System bestemmes igjen af den Alt begribende, absolut
mægtige Selveenhed, som i uendelig Bestemmen, Begriben og Virken er absolut
Causalitet.

\emph{γ) Identitetsdualisme.} Det Absolutes og det Relatives Philosophi,
Repræsentanterne for Idealisme og Realisme, stille sig som to Extremer hver
paa sin Side af Sandhedens Idee; at de begge ere behæftede med principiel
Usandhed, kan sees deraf, at den ene conseqvent maa tilintetgjøre den anden.
Idealismen forudsætter, at det Almene er det objectivt Altbegribende, der
begriber sig selv. Naar denne Sætning gjøres tydelig, viser sig strax en
Mystification. Skal Sætningen tages strengt objectivt, saa maa det jo være
den almene, af alt Værende fyldte Væren, med andre Ord: Universet, der
under Virkelighedens Form begriber sig selv. At nu det virkelige Universum
udtrykker et System af Begreber, maa indrømmes; men Et er, at det Objective
udtrykker et System af Begreber, et Andet, at det begriber sig selv. Det
sig begribende Almene maa nødvendig opfattes som Selvvæsen, men et saadant
Selvvæsen er ikke objectivt; det er tvertimod subjectivt.

Realismen gjør det Sande umiddelbart til Et med det Virkelige; det er en
Usandhed, thi det er en Modsigelse. Saa lidt som der kan være noget
Objectivt, som ikke er objectiveret, lige saa lidt kan der være noget i
Virkelighed Sandt, som en Vidende ikke har sandet. Det Virkelige, fastholdt
for sig, er hverken sandt eller usandt; Sandheden beroer paa en Dom, og
Dommens Sandhed beroer paa en Slutning. Sandheden fordrer Dualisme af Viden
og Væren og i denne Dualisme absolut Identitet. Dualismens Identitet er den
Selveenhed, der magter Extremerne, den Slutningseenhed, der fra Evighed af
er den absolute Begyndelseseenhed.

\section{Mathematik}
\label{sec:12}
% pp. 117--125

\noindent Skjøndt Mathematik og formel Logik ere nær beslægtede
Videnskaber, saa er det dog en Miskjendelse af deres
Charakteereiendommeligheder, naar man, hvad Nogle have forsøgt, vil
reducere Logik til Mathematik, eller, hvad der er endnu urimeligere,
deducere Mathematikens Væren af Qvantitetsbegreber i den immanente Logik.
Hvor forskjelligt de to Videnskabers Grundlag er, kan sees af det
oprindelige Forhold, hvori enhver af dem sætter det Enkelte til det Almene.
Saa megen Vægt Logiken end for Concretionens Skyld maa lægge paa
Bestemmelsen af de enkelte Leds eiendommelige Forhold, saa kan den under
sin Begrebsudvikling dog ikke undgaae at lade det Almene fortære det
Enkelte, ligesom Tanken fortærer Forestillingen, hvorved Alt i sidste
Instans kommer til at beroe paa Gjennemsigtighed i Reflexion. Anderledes i
Mathematiken. Den begynder med det abstract Almene, for saa vidt den
begynder med at reducere alle Led til Størrelser og Forholdene til rene
Størrelsesforhold; men den sikkrer sig tillige det Enkelte saa punktuelt,
articulerer Forholdsbestemmelserne saa specifisk, at den netop derved faaer
Charakteer af exact Videnskab. Den kan paa Grund heraf ikke nøies med
Reflexionens Almindelighed, den kræver Operationen. For Reflexionen er
Talesproget tilstrækkeligt, thi i Ordet har det Almindelige Overvægten; men
for Punktualiteten i den exacte Analyse er Ordet ikke tilstrækkeligt;
Mathematiken maa, for Operationens Skyld, have sit eget pasigraphiske
Tegnsprog.

Forholdet mellem Reflexion og Operation kan tydelig sees paa de
mathematiske Opgaver. Betingelsen for at løse et videnskabeligt Problem er
fremfor Alt at faae det skarpt bestemt og rigtig stillet. Adskillige
Problemer henstaae i Videnskaben saa dunkle, at man ikke ret veed, om de
ere løselige eller uløselige, fordi de ikke ere opfattede med behørig
Skarphed og seete i fuld Belysning. Den exacte Videnskab kan i denne
Henseende siges at gaae foran med et godt Exempel, naar den øver sine
Elever i at sætte i Ligning. Opgaven bliver i Almindelighed fremstillet i
Reflexionens Form; men ved nærmere Betragtning viser det sig, at den maa
stilles endnu engang, nemlig i Operationens Form, for at kunne løses:
Opgaven bestaaer væsenlig i at faae Opgaven stillet, nemlig som Opgave for
Operationen. Følgende Opgave, som giver to Ligninger af første Grad, kan
tjene til Exempel. I Reflexionens Form fremstilles den saaledes: »Naar et
Antal af $m$ Oxer i $n$ Dage afæde $p$ Tønder Land Græsgang, ligeledes $m'$
Oxer i $n'$ Dage $p'$ Tønder Land, hvormange Oxer ($m''$) ville da i $n''$
Dage afæde $p''$ Tønder, idet Græsset voxer eensformigt og har paa alle tre
Græsgange den samme Høide fra Begyndelsen?« Enhver, der forstaaer Sproget og
veed, at $m$, $n$, $p$ o.s.v.\ betyde forskjellige Tal, forstaaer for saa
vidt Opgaven, og kan derover anstille sine Reflexioner; men Opgaven er at
stille Opgaven, d.v.s.\ at stille den i Ligning.

Af Forholdet mellem Reflexion og Operation forstaaes det eiendommelige
Eftertryk, den exacte Videnskab lægger paa Beviset. Den kan bruge Billedet,
som i Geometrien, Forestillingen, som i Arithmetiken, Reflexionen, som i
det Uendeliges Analyse; den har faste Grundsætninger, men den ændser ingen
Autoritet og tager Intet paa Skjøn: den fordrer Beviser overalt, og dens
Beviser ere uimodsigelige. Man skulde dog mene, at en Læresætning som den
pythagoræiske, saa ældgammel og beviist saa mange tusinde Gange, fra
forskjellige Synspunkter paa forskjellig Maade, omsider kunde ophøies til
staaende Autoritetssætning. Ja, kom det blot an paa Vished, saa er denne
Læresætning ganske vist hævet over enhver Tvivl; men den exacte Videnskab
kræver endnu Beviset, thi Beviset er ikke blot for Vishedens Skyld, men
væsenlig for Begrundelsens, for Indsigtens Skyld. Heller ikke kan Noget
antages efter det blotte Øiesyn: at en Curve parallel med en Cirkel selv
er en Cirkel, er øiensynligt; men den, som ifølge umiddelbart Skjøn
antager, at en Curve parallel med en Ellipse nødvendig maa være en Ellipse,
tager mærkelig feil. Skulde man dømme efter hvad man ligefrem kan see, saa
maatte Asymptoten ved fortsat Forlængelse komme til at berøre sin Curve;
men Beviset siger Nei, og Anskuelsen maa bøie sig for Tanken. At en ret
Linie af lige Længde altid betegner det Samme, vil vist Ingen betvivle ---
undtagen Mathematikeren. Kun han veed, hvor forskjelligartede de Elementer
kunne være, hvoraf en Linie bestaaer; han veed, at en Ellipse eller Parabel
kan ligge opløst i Linien; han veed, at to hinanden skjærende rette Linier
kunne forestille en Hyperbel.

Den reflecterede Eenhed af Reflexion og Operation har sit almindeligste
Udtryk i den mathematiske Function. Uden nogensomhelst Indskrænkning i
Bestemmelsen af det Almene i de rene Størrelsers indbyrdes Afhængighed af
hverandre, kan den exacte Videnskab specificere Afhængighedsforholdet saa
punktuelt, at Articulationen gaaer i det Uendelige. I Ligningen $y = f(x) =
ax$ være $a$ den constante, $y$ og $x$ de variable Størrelser. Man seer nu
strax heraf, hvilken Forskjel der er mellem faste og flydende Grændser. Den
constante Størrelse er vilkaarlig, den kan være saa stor og saa lille som
man vil; men naar den een Gang er sat, saa ere dens Grændser faste.
Størrelserne $y$ og $x$ have derimod saa bevægelige Grændser, at disse ikke
kunne fastsættes, uden at de paagjældende Størrelser forandre deres
Charakteer. Man seer heraf Forskjellen mellem Størrelse og Tal; ethvert Tal
er en Størrelse, men ikke enhver Størrelse er et Tal: $x$ og $y$ ere
Størrelser, $a$ er et Tal. Afhængighedsforholdet mellem $y$ og $x$ er
gjensidigt; naar $y$'s Afhængighed af $x$ betegnes ved $y = f(x) = ax$, saa
kan $x$'s Afhængighed af $y$ betegnes ved $x = \Phi(y) = \dfrac{y}{a}$.
Betegnelserne $f$ og $\Phi$ gjælde ikke Størrelserne selv, men deres
Afhængighedsmaade, og i deres Afhængighedsmaade sees Loven.

Ved Hjælp af Functionen bliver den exacte Videnskab i Stand til at
articulere Forholdet mellem Endeligt og Uendeligt paa en saadan Maade, at
Vexelbestemmelsen bliver indlysende. Er $y = ax$, saa er ogsaa $dy = a\,dx$,
idet $dy$ og $dx$ betegne Elementer, forsvindende smaa Dele, d.e.\
Differentialer af $y$ og $x$. Men heraf følger igjen, at man maa faae
$\dfrac{dy}{dx} = a$. Dette Resultat, saa simpelt det end er, er uendeligt
frugtbart. Det beviser nemlig, at et constant Forhold kan bevares, om end
Forholdsleddene indsvinde i det Uendelige, og at det ellers Ubestemmelige
kan bestemmes. Ifølge den angivne Function maa $\dfrac{0}{0}$ være liig med
$a$; men samme Function maa da ogsaa kunne give $\dfrac{\infty}{\infty} =
a$; Forholdet bevares, idet Tilvæxten til $y$, den afhængige Størrelse,
stadig maales med Tilvæxten til $x$, som her er den uafhængige. Maaler man
$a$ med $dx$, saa er $\dfrac{a}{dx} = \infty$; maaler man $dx$ med $0$, saa
er ogsaa $\dfrac{dx}{0} = \infty$. Det Uendelige synes for saa vidt at være
det Umaalelige, det Uopnaaelige; men gaae vi ud fra Ligningen $dy = a\,dx$,
og betegne den uendelige Summation ved $\int$, Integraltegnet, saa giver
$\int_{x=c}^{x=b} a\,dx = y = ab - ac$ en endelig Størrelse, hvilket
beviser, at det Uendelige, saa vel det uendelig Lille, $dx$, som det
uendelig Store, det ved $\int$ betegnede Antal, falder inden for endelige
Grændser. Heri ligger igjen Muligheden til en articuleret Bestemmelse af
Forholdet mellem Discretion og Continuitet.

Logisk forstaaet, er Discretionen Endelighedens, Continuiteten
Uendelighedens Form. Existensen er discret, »den rene Væren« er
continuerlig; Atomerne ere i Discretion, men Rummet er et Continuum;
Begivenhederne ere discrete, Tiden er Continuitet. Hvorledes komme vi nu
til Discretionens Endelighed? Ved et Brud paa Continuiteten, altsaa ved en
Udstykning. Men hvorledes komme vi saa tilbage til Uendeligheden? Ikke ved
en Sammenstykning; thi Sammenstykningen tilhører Endeligheden lige saa vel
som Udstykningen. Ved et Spring da? Ganske vist, naar Forskjellen mellem
Endeligt og Uendeligt fastholdes qvalitativt. Men den exacte Videnskab, der
i alt det Qvalitative seer paa det Qvantitative, kjender endnu en Overgang,
Tilnærmelsen, den uendelige Tilnærmelse. Adskilt fra Operationsbegrebet kan
Reflexionen saare let paavise Modsigelsen i et Begreb som uendelig
Tilnærmelse. Tilnærmelse betyder jo, at man bestandig kommer det Søgte
nærmere, saaledes er $2$ nærmere ved $100$ end $1$, $10$ endnu nærmere; men
naar det Uendelige altid er uendelig langt borte fra det Endelige, saa
betyder jo Tilnærmelsen Intet. En uendelig Tilnærmelse, en Tilnærmelse til
det, som man dog aldrig kan komme nærmere, er jo en Tilnærmelse, som dog
ikke er en Tilnærmelse, altsaa en Modsigelse.

Den exacte Videnskab magter Opgaven, idet den forstaaer at give een og
samme Størrelse baade en endelig og en uendelig Form, og derved i
Operationen bevise, at det Uopnaaelige er opnaaeligt, og det Opnaaelige
uopnaaeligt. I Ligningen
\[
  \frac{1}{1-x} = 1 + x + x^2 + \cdots x^n + \frac{x^{n+1}}{1-x}
\]
have vi et Exempel. Er $x > 1$, saa er Rækken divergent; ved at tænke sig
den fortsat i det Uendelige, vilde man komme uendelig langt bort fra
Maalet; men føier man Resten til, er den uendelige Bortfjernelse hævet: man
er ved Maalet. Er $x < 1$, saa er Rækken convergent, og Tilnærmelsen viser,
at i hvor vel Rækken er uendelig, og Maalet for saa vidt uopnaaeligt, saa
bliver Resten dog mindre og mindre: den uendelige Tilnærmelse er altsaa dog
en virkelig Tilnærmelse. Giver man Tallet $2$ Formerne $\dfrac{1}{1-\frac12}$
og $1 + \frac12 + \left(\tfrac12\right)^2 + \cdots \left(\tfrac12\right)^n$,
da er det ganske vist, at $\left(\tfrac12\right)^n$ er uendelig langt fra
$0$, saa længe $n$ er endeligt; men Rækken har tillige den Egenskab at vise
os Overgangen fra endelig til uendelig Discretion. Jo mindre Differensen
mellem $\left(\tfrac12\right)^n$ og $\left(\tfrac12\right)^{n+1}$ bliver,
desto mere nærmer Discretionen sig til Continuitet, og i Continuiteten er
Maalet naaet, thi $0 + 0 = 0$.

Den uendelige Discretion er altsaa Mellembestemmelsen mellem Continuiteten
og Udstykningen. For Analysens Skyld er Udstykningen, den endelige
Discretion, nødvendig; men Mathematiken hæver Udstykningen og faaer den
uendelige Discretion articuleret i Conseqvens af Functionen. Den
sammenstykker ikke krumme Linier af Polygoner, den hæver Udstykningen ved
en uendelig Summation af Elementer. Er $y = \sqrt{a^2 - x^2}$ en Ligning
for Cirklen, idet $a$ betegner Radius, og Coordinaternes Skjæringspunkt
ligger i Centrum, saa er $dy = \dfrac{-x\,dx}{\sqrt{a^2 - x^2}}$ et
Perepheriens Element og $\displaystyle\int_0^{2\pi} \frac{-x}{\sqrt{a^2 -
x^2}}\,dx$ den hele Cirkelperipheri som flydende, continueret Curve.

Saa længe Reflexionen endnu ikke var methodisk forligt med Operationen, var
den bestandig udsat for at forvildes. Mange af de gamle Tænkeres Paradoxer,
der see ud som Sophismer, ere i Virkeligheden, naar vi see nærmere til,
videnskabelige Problemer; saaledes den eleatiske Sophisme Achilles. Det var
ikke Dialektik, hvortil Eleaterne trængte; en endelig Beregning kunde
heller ikke hjælpe dem. Det var en tydelig Indsigt i Forholdet mellem
Continuitet og Discretion, de savnede; men den udisciplinerede Forstand
kan ikke fatte det, og saa hedder det strax: det Absolute, det Uendelige,
er os aldeles ubegribeligt; at tænke derover er kun ørkesløst Grubleri.
Men, mærkeligt nok, medens man paa realistisk, positivistisk Viis absolut
vil forkaste det Absolute, har man dog Betænkeligheder ved at forkaste det
Uendelige. Dette er meget inconseqvent. Kan det Endelige ikke forstaaes
uden det Uendelige, saa kan det Relative heller ikke forstaaes uden det
Absolute. Men Forstaaelsen er forskjellig. Analysen af det Absolute er en
Idee-Analyse, der hører under Logiken; Analysen af det Uendelige er en
Operationsanalyse, der hører under Mathematiken.

\section{Rationel Mechanik}
\label{sec:13}
% pp. 125--135

\noindent Skjøndt den rationelle Mechanik paa en Maade kan betragtes som
en Green af Mathematiken, maa den dog i nærværende Sammenhæng fremhæves som
en særegen Videnskab, fordi den i sit Grundlag optager eiendommelige
Forudsætninger, nemlig Materien og dens Kræfter. At den ikke desto mindre
er en apriorisk Videnskab, ligger deri, at det ikke er det Virkelige, men
det Forestillede, forestillet Materie, forestillede Kræfter, hvorpaa
Eftertrykket ved de almene Loves Undersøgelse væsenlig falder. Kraften er
det Bevægende, Materien det, som bevæges; charakteristisk for Materien er
dens Inerti, den kan, naar den er i Hvile, ikke sætte sig selv i Bevægelse,
og naar den er i Bevægelse, ikke sætte sig selv i Hvile. Idet Materien
bevæges af Kræfter, bevæges den af udvortes Aarsager, thi det er altid
Kraften i den ene Deel, der har den anden til Angrebspunkt. Den Maade,
hvorpaa den rationelle Mechanik har sikkret sig Betingelserne for
Kræfternes Analyse, er betegnende for dens hele Methode.

Uagtet Kræfterne ere usynlige, kunne de dog som Bevægelsernes forestillede
Aarsager synliggjøres ved Linier; en ret Linie, Bevægelsens forestillede
Spor, angiver saaledes baade Styrken og Retningen af en enkelt Kraft.
Liniens Retning svarer til Kraftens Retning, og dens Længde til Kraftens
Styrke. Det er nemlig indlysende, at af to Kræfter, der skifteviis virke
paa den samme Masse under samme Betingelser i lige Tid, vil den være den
stærkeste, som bringer Massen til at gjennemløbe den længste Vei. Virke
Kræfterne $P$ og $Q$ paa samme Punkt, vil deres Styrkes Resultat,
Resultanten, $R$, være liig med deres Sum, naar de virke i samme Retning,
og i modsat Fald liig med deres Differens: $R = P + Q$. Ere Retningerne
modsatte, og $P = Q$, saa er $R = 0$; der er altsaa Ligevægt, og ingen
Bevægelse. Ere Retningerne modsatte, men $P > Q$, saa er $R = P - Q$ ikke
$0$, her er Bevægelse i Retning af $P$. Samvirkende Kræfter med fælles
Angrebspunkt, deres endelige Størrelsesforhold være iøvrigt
hvilketsomhelst, kunne danne uendelig mange forskjellige Vinkler med
hinanden, to og to; det bliver derfor en Opgave at bringe de utallige
Tilfælde ind under en omfattende, almeengjældende Form: denne Form er
Kræfternes Parallelogram. Betegnes den Vinkel, de to Kræfter $P$ og $Q$
danne med hinanden, ved $(P.Q)$, deres Resultant ved $R$, saa gjælder
Ligningen $R = \sqrt{P^2 + Q^2 + 2PQ\cos(P.Q)}$ for alle mulige Tilfælde.
Er Vinklen $(P.Q) = 0^\circ$, saa er $\cos = 1$, og $R = P + Q$; er Vinklen
$(P.Q) = 180^\circ$, saa er $\cos = -1$, og $R = P - Q$. Parallelogrammets
Diagonal er Resultanten af de to Sidekræfter, og disse dens Composanter;
men hver Resultant kan igjen blive Composant i et andet Parallelogram, og
hver Composant en Resultant af andre Composanter. Parallelogrammets Lov er
en Grundlov for Sammensætning og Opløsning af mechaniske Kræfter, hvoraf et
uendeligt System af særskilte Love lader sig udvikle.

I Kræfternes Parallelogram giver den rationelle Mechanik os et anskueligt
Billede paa Eenhed af Realbegreb og Lov. Men den exacte Videnskab udvikler
ikke Begrebet som Begreb; den maa, for Punktualitetens Skyld, holde fast
paa det Forestillede, paa Forholdsleddene i deres udvortes Relationer, og
desaarsag fremstille Begreberne i Lovenes Form.

Forskjellen mellem Lov og Realbegreb beroer paa, at et Indbegreb af
sammenhørende Bestemmelser kan opfattes deels saaledes, at Leddene
forestilles hvert for sig, og deres indbyrdes Afhængighed viser sig da i
Loven; deels saaledes, at Leddene forvandles til Reflexer i den bestemte
Totalitet, og den i sine Reflexer reflecterede Totalitet viser sig da som
Begrebet. Den i forrige Paragraph berørte Qvotient $\dfrac{y}{x} =
\dfrac{dy}{dx} = a$ staaer netop paa Overgangen; $y$ og $x$ ere endelige
Led, $dy$ og $dx$ ere Reflexer. Den exacte Videnskab maa fastholde det
Forestillede og kommer i sine Analyser først og sidst til at dvæle ved
Forholdsarticulationen, d.e.\ ved Loven. Selv i Variationsregningen, hvor
Maxima og Minima ikke blot afhænge af varierende Led, men hvor disse Led
selv afhænge af vilkaarlige Functioner, kunne Leddene dog ikke forvandles
til rene Reflexer, men maae for Lovenes Skyld bevares i Forestillingens
Punktualitet. Denne Punktualitet træder særlig frem i den rationelle
Mechanik.

Vi see det tydelig i Vexelbestemmelsen af Kraft og Aarsag. Aarsagen med
sine Betingelser hører væsenlig under Reflexionen og gaaer op i
Realbegrebet, men Kræfterne med de forestillede Masser og materielle Led
falde ind under Lovene. Mechaniken skjelner mellem accelererende,
bevægende og levende Kraft, og faaer, idet Materiens Energi paa det
Strængeste fastholdes, en Række Love udviklede; den exacte Videnskab
udstykker Begrebet og beviser ved denne Udstykning, at den netop i sin
Eensidighed har sin Styrke.

\emph{Accelererende Kraft.} Som Bevægelsens Aarsag kan Kraften enten virke
pludselig, som i Stødet, og faaer da Navn af instantan Kraft, eller
vedblivende forandre Hastigheden, og kaldes da, positivt eller negativt,
accelererende. Men i Begrebet selv er Forskjellen forsvindende. Den
instantane Kraft udkræver en, om end nok saa lille, endelig Tid, for at
bevirke en endelig Hastighed og er for saa vidt væsenlig accelererende;
omvendt kan den accelererende Kraft opløses i uendelig mange Stød og er i
hvert Stød at betragte som instantan. Spørge vi nu om Aarsagen, da er
Kraften i og for sig som Kraft ikke Accelerationens Aarsag, den er det kun
formedelst Inertien. Ved endelig Adskillen af den constante Kraft, den i
Materien værende Inerti og det af begge Resulterende faae vi det exacte
Udtryk f.\ Ex.\ for Faldloven. En constant Kraft $k$ virker paa et
materielt Punkt; i det relative Øieblik $dt$ er $k\,dt = dv$, idet $t$ er
Tiden, $v$ Hastigheden. Det materielle Punkt modtager i hvert Øieblik en
uendelig lille, om end constant Hastighedstilvæxt. Det, som gjør, at
Hastighedstilvæxten, som kun er $dv$, i næste Øieblik er $2\,dv$, i næste
igjen $3\,dv$ o.s.fremdl., saa at den i det $n$te Øieblik er voxet op til
$n\,dv = v$ for $n = \infty$, er netop Inertien, som bevarer enhver
Hastighedstilvæxt og derved gjør det muligt, at Hastigheden kan voxe med
Tiden, være en Function af Tiden: $v = f(t) = kt$. Men først naar Kraft og
Inerti sees som Reflexer i den samme Virken, faae vi Begrebet Aarsag.

\emph{Bevægende Kraft.} Her bliver Kraften opfattet som Aarsag i Relation
til den forestillede Masse. Er Massen $M$ større end Massen $m$, og $g$
Accelerationen, da har den bevægende Kraft i $Mg$ unegtelig en større
Styrke end i $mg$; men medens der imellem $Mg$ og $mg$ kun er en
qvantitativ Forskjel, eftersom de begge høre ind under samme
Aarsagsforestilling, er der derimod en qvalitativ Forskjel mellem $g$ og
$mg$, thi her er Aarsagsforestillingen forandret. For Bevægkraften som
forestillet Aarsag er det, strængt taget, ikke nødvendigt, at den
foraarsager en virkelig Bevægelse. Vel er Aarsagen ikke uden Virkning, men
Bevægelsen kan være ophævet ved en Modvirkning, et Tryk, og Aarsagen har
ligefuldt sin Realitet. Saaledes med Vægten, et Produkt af Masse og
Acceleration uden kjendelig Bevægelse: $Mg = P$. Af Vægten $P$ sees det, at
Massen kun har Realitet i Forhold til Tyngden. Forandrer Tyngden sin
Intensitet, saa faaer Massen en anden Vægt og, hvad Ligningen $M =
\dfrac{P}{g}$ viser, en anden Realitet. Men saa fyldestgjørende den
endelige Udstykning af Forholdsleddene end er ved en exact Bestemmelse af
Relationer og Love, saa ufyldestgjørende er den, naar der spørges om
Realbegreberne selv og de væsenlige Aarsager.

\emph{Den levende Kraft.} Ved at gaae fra Forestilling til Forestilling,
fra Forudsætning til Forudsætning, ændre de givne Betingelser og indføre
nye, kan Mechaniken systematisk skride frem fra simplere til mere
complicerede Relationer, til nye Love og nye Formler. Forestiller man sig
Bevægkraftens Modstand i Trykket hævet, saa indtræder Bevægelsen, Tiden
faaer sin Betydning, og Bevægelsen lægger sig for Dagen i Produktet af
Masse og Hastighed, d.e.\ i Bevægelsesmængden. Kalde vi Bevægkraften $mk$,
Tiden $t$ og Hastigheden $v$, saa udviser Ligningen $mkt = mv$, altsaa $mk
= \dfrac{mv}{t}$, at en hvilkensomhelst konstant Kraft kan maales med den
Bevægelsesmængde, den frembringer i en Tidseenhed. Ikke desto mindre er
Tidens Mellemkomst ubekvem for Maalet. Naar to Legemer have samme
Hastighed, men det ene har erholdt den ved et pludseligt Stød, det andet
ved at falde i længere Tid; vil man saa paa Grund af Tidens Forskjellighed
kunne sige, at deres Kræfter ere forskjellige? Det vilde jo, for at bruge
et Udtryk af Leibnitz, være det Samme som at sige, at en Mand var rigere,
fordi han havde været længere Tid om at vinde sine Penge. Kraften $mk$ kan
i Tiden $t$ vække en Bevægelsesmængde i $m$, der er $mV$, og i $M$, der er
$Mv$: altsaa $mV = Mv$. Spørge vi nu, hvad en saadan Kraft kan udrette,
idet Tiden anvendes til at bevæge Massen eller rettere Vægten frem i
Rummet, saa komme vi til den levende Kraft. Sættes $mkt = mV$, vil
Virksomheden i det første Tidselement være $mk\,dt = m\,dV$. Hvad udrettes
nu i den løbende Tid? Kraften bevæger Byrden fremad i Rummet, men idet den
slæber paa Byrden, aftager Hastigheden, og den aftagende Hastighed er et
Udtryk for den aftagende Kraft. Hvorlænge Kraften kan holde ud, maa da sees
paa den Veilængde, $s$, hvorigjennem den formaaer at føre sin Byrde.
Ligningen $m\,dV = mk\,dt$ maa altsaa forvandles til $m\,dV =
mk\dfrac{ds}{V}$, saa at det kjendes paa Rummet, hvad der udrettes i Tiden.
Det gamle Ord, at Tiden er lige lang, men ikke lige nyttig, finder her en
egen Bekræftelse; $dt$, det relative Øieblik, er constant, men $V$ varierer,
hvoraf følger, at Rumselementet $ds$ ogsaa varierer: $ds > ds^1 > ds^2
\ldots$ efterdi $V > V^1 > V^2 \ldots$ Af $m\,dV = mk\dfrac{ds}{V}$ følger
$mV\,dV = mk\,ds$. Ved at integrere denne Ligning mellem Grændser erholde
vi $mV^2 = 2mks$ og med det Samme et Maal for Arbeidet. Men naar der til
Bevægelsesmængden $mV$ svarer en levende Kraft $mV^2$, saa svarer der til
Bevægelsesmængden $Mv$ en levende Kraft $Mv^2$. Da nu $mV^2$, ifølge
Forudsætningen $mV = Mv$, har en anden Værdi end $Mv^2$, saa følger heraf,
at Kraften i $mV$, naar den sættes i Arbeide, udretter noget Andet, end
Kraften i $Mv$. Den levende Kraft, der maales med Arbeidet, er altsaa
qvalitativt forskjellig fra den i Tid bevægende Kraft, der maales med
Bevægelsesmængden.

Stille vi os nu paa Logikens Standpunkt, altsaa udenfor Mechaniken, da
maae vi erkjende, at den exacte Videnskab ved sine Analyser netop
forskaffer os de Reqvisiter, der høre Realbegrebet til, naar Reflexionen
gjennemføres. Aarsagsbegrebet er et Realbegreb, et Virkelighedsbegreb, og
et Virkelighedsbegreb har altid Reflexer af virkelige Ting i sit Indhold.
Tingsreflexer ere Endelighedsreflexer, Reflexer af Udstykning; men i
Begrebet maa Udstykningen hæves. Reflexionen maa altsaa paa een Gang baade
sætte og hæve det Forestillede. Udstykningsreflexer af den forestillede
Aarsag ere: Betingelsen og det Betingede. Herunder indgaae alle
Bestemmelser. End ikke Kraften kan i denne Belysning holdes fast, thi Kraft
for sig er kun en Forestilling; det Forestillede, vi kalde Kraft, synker
ned til at være en Betingelse. Det Betingede, det Afhængige, indeslutter i
Grunden Betingelsen, det Uafhængige, og dog stille de sig imod hinanden som
Led imod Led. Betingelsesreflexionen er en Reflexion af forestillede Led,
en udvortes, mellem Væsen og Existens, Begreb og Virkelighed svævende, for
den mechaniske Virksomhed betegnende Reflexion. Saaledes ere Kraften og
Inertien Betingelser for Accelerationen; de forestilles hver for sig, og
vise dog uophørlig hen til hinanden. Massen og Accelerationen ere igjen
Betingelser for Vægten; den samme Masse kan undergives en anden
Acceleration, og den samme Acceleration appliceres paa en anden Masse; for
saa vidt maa Betingelserne adskilles. Og dog kunne de ikke adskilles, thi
det forandrede Produkt af Masse og Acceleration forandrer Vægten. Bevæges
Vægten, drages Tiden med ind i Betingelsernes Kreds, bliver den heraf
betingede Hastighed medbetingende, og den bevægende Kraft viser sig som
Bevægelsesmængde. Dette er igjen en Betingelse for, at Bevægelsesmængden
kan maales med det Rum, Modstanden gjennemløber. Saaledes have vi den
levende Kraft; dog endnu ikke som Aarsag, thi i Udstykningsreflexerne er
Aarsagen opløst: den er forvandlet til en Sum af Betingelser.

Realbegrebet danner sig nu logisk derved, at et Indbegreb af Betingelser
reflecteres i hverandre; Udstykningen forsvinder, og den i sine Reflexer
reflecterede Totalitetseenhed, Begrebet, viser sig som det Bestemmende.
Denne Uendelighedsreflexion har fristet »den immanente Logik« til at hæve
sig over Mathematiken og bilde sig ind, at den uden den rationelle
Mechaniks Mellemkomst selv kunde magte Lovene og de endelige Betingelser.
En i Videnskaben beklagelig Misforstaaelse! Immanenslogiken har ingenlunde
Uret i at ville ophæve Udstykningen; dens Uret bestaaer tverimod deri, at
den med al sin Reflecteren ikke tilstrækkelig gjennemskuer Reflexionens
dobbeltsidige Væsen. I Ophævelse af Udstykning maae, vel at mærke, alle
Momenter af det Udstykkede tydelig skinne igjennem; hvis det ikke skeer,
saa træder Indbildningen i Kjendsgjerningens Sted, d.v.s.\ Speculationen
bliver indbildt. Man naaer ikke det Absolute ved at oversee det Relative;
man begriber ikke det Ubetingede ved en almindelig Opløsning af alle
Betingelser; man maa udtrykkelig kunne gjenkjende de relative Led i deres
Reflexer, de endelige Betingelser i den opløsende Form: heraf sees Logikens
Afhængighed af de exacte Videnskaber.

Den rationelle Mechanik, der har sit Mesterskab i systematisk Behandling af
lovbestemte Forhold mellem det Betingede og dets Betingelser, tydeliggjør
da ogsaa, hvorledes særskilte Betingelser kunne henføres under een
Grundbetingelse, og denne igjen udvikle sig i en Mangfoldighed af
Enkeltbetingelser. For at Kræfter, der virke paa et Punkt, kunne være i
Ligevægt, maa deres Resultat være Nul; dette er Ligevægtens almene
Betingelse, dens Grundbetingelse. Men hvilken Vrimmel af særskilte
Betingelser og varierende Enkeltbetingelser kan ikke heri indgaae!
Kræfternes Ligevægt er Hvilens Aarsag; Bevægelsens Aarsag er Ligevægtens
Ophævelse: det er i Grunden samme Aarsag og dog under Betingelsernes Vexlen
en forskjellig Aarsag. Reflectere vi paa Forskjellen, saa have Statiken og
Dynamiken hver sin Kreds af Opgaver, og falde for saa vidt ud fra
hinanden; reflectere vi paa Eenheden, saa bliver den statiske Opgave
dynamisk, og den dynamiske statisk. Vexelbestemmelsen grunder sig paa
Reflexion af to til hinanden svarende Principer, de virtuelle Hastigheders
Princip og D'Alemberts Princip for de tabte Kræfters Ligevægt.

Den dybe Kløft mellem det Logiske i de almene Realbegreber og det
Antilogiske i Tingenes Verden kan vistnok ingen menneskelig Viden udfylde,
men den rationelle Mechanik kan bygge en Bro derover.


\chapter{Realempiriske Videnskaber}

\section{Mechanik og chemisk Physik}
\label{sec:14}
% pp. 135--151

\noindent De i Naturen virkeliggjorte Begreber ere mere end Begreber, det
er existerende Ting; her begynder Physiken, og Logiken kan som formel
Logik, særlig som Inductionslogik, blot tjene som veiledende. Vel ere
Tingene undergivne almindelige Love; men i deres virkelige Opfyldelse
afvige Naturlovene fra abstracte Love; Planeterne bevæge sig dog ikke i
Ellipser, den ballistiske Curve er ingen Parabel, og Loven for det
legemlige Pendul har et andet Udseende end den for det mathematiske. Den
rationelle Mechanik betræder en fremmed Grund og bliver Hjælpevidenskab for
mechanisk Physik. Her er en ny Methode, fordi her er et nyt Grundlag, d.e.\
de ydre Erfaringers Grundlag. Hverdagserfaringer ere for usammenhængende og
overfladiske. Den sandselige Ting med sine Egenskaber staaer i saa
mangesidige og forskjelligartede Forhold til andre Ting med tilsvarende
Egenskaber, at der behøves en egen Kunst for at løse Relationernes Knude.
Hvad der under Forviklingerne maa fra et valgt Synspunkt fastholdes som
væsenligt, kan, naar Synspunktet forandres, være at betragte som noget
uvæsenligt, ja som noget Tilfældigt, der griber forstyrrende ind i
Forstaaelsen af det, hvorpaa det kommer an. Analysen er mere end logisk,
den er reel; det Forstyrrende kan ikke fjernes ved den blotte Abstraction,
det maa, saa vidt det er muligt, fjernes ved Experiment. Forskningens
Opgave er at rette hvert Spørgsmaal til den virkelige Natur i en saa exact
Form, at Kjendsgjerningen selv nødvendig maa svare enten Ja eller Nei.

Som et lysende Exempel paa Vexelvirkning mellem mechanisk Physik og
rationel Mechanik staaer Newtons Opdagelse af den almindelige Tyngdelov.
Det Øie, der saae i det faldende Legeme en Yttring af den Kraft, der ogsaa
drager Maanen ned mod Jorden, tilhører ikke Sandsningen alene, men ogsaa
Tænkningen. Det var Tænkerens Anelse om en almindelig, alle Legemer iboende
Væsenhed, der gav Impulsen til den store Opdagelse. At Maanens Fald maa
forholde sig til Æblets, som den Kraft, der drager Maanen, til den Kraft,
der drager Æblet, er en Selvfølge; men i hvilket Forhold staae nu disse
Kræfter til hinanden? Den Kraft, der drager Æblet, er bekjendt ved
Faldloven $15\tfrac58$ Fod i 1ste Secund; den Kraft, der drager Maanen kan,
da Maanens Middelafstand og Middelomløbstid ere bekjendte Størrelser,
ligeledes bestemmes. Betegnes Jordens Afstand fra Maanens Midtpunkt ved
$R$, den til Buen (i et Tidssecund) svarende Vinkel ved $\Theta$, saa er
$\rho$, det Stykke, som betegner Maanens Fald mod Jorden, udtrykt ved
Ligningen $\rho = 2R\sin^2\tfrac12\Theta$. Kræfternes Forhold er altsaa
bestemt ved Qvotienten $\dfrac{2R\sin^2\tfrac12\Theta}{15\tfrac58}$; men,
for at $R$ kan bestemmes, maa, vel at mærke, tillige Jordens Radius $r$
være bekjendt. Her er en Kreds af forudsatte Iagttagelser. Og dog kunde
Newton ikke uden særlig Forberedelse faae den beriget ved nye
Iagttagelser. Vel laa det nær at søge Grunden til Tiltrækningens
Forskjellighed i Afstandens Forskjellighed, men med denne almindelige
Formodning kunde endnu ingen Erfaring gjøres; først ved den Antagelse, at
de tiltrækkende Kræfter staae i omvendt Forhold af Afstandens Qvadrater,
blev Ligningen bestemt:
\[
  \frac{2R\sin^2\tfrac12\Theta}{15\tfrac58} = \frac{\dfrac{\mu}{R^2}}{\dfrac{\mu}{r^2}} = \frac{r^2}{R^2}.
\]
Er Forholdet $\dfrac{\mu}{r^2}$ det rette? Det var Newtons store
Spørgsmaal. Svaret lod vente paa sig, men det udeblev ikke.

At de samme mechaniske Grundlove, der gjælde for Kloderne i Himmelrummet,
ogsaa maae gjælde for Legemerne paa Jorden, er i Videnskaben en uomstødelig
Forudsætning; men da Opfyldelsen varierer med Betingelserne, da Legemerne
paa Jorden ligge os anderledes nær og paa en ganske anden Maade kunne
gjøres til Gjenstand for Iagttagelser, Forsøg og Efterforskninger end hine
fjerne Legemer, forstaaer det sig af sig selv, at særlige Tilfælde med
særlige Opgaver og specifiske Betingelser her anderledes krydse og afbryde
den systematiske Continuitet, end ved Himmellegemernes Iagttagelse. Intet
Under, at Astronomien paa Grund af sin abstractere Charakteer, skjøndt
endnu ikke fuldendt, dog staaer med et Præg af Fuldendelse, der udmærker
den fremfor enhver anden empirisk Videnskab. Hvor inderlig det Rationelle
kan smelte sammen med det Empiriske, har Laplace beviist i sin
\textit{Mécanique céleste}.

I den mechaniske Physik bevare de virksomme Led, Atomer, Moleküler, Masser
under deres indbyrdes Vexelvirkning en vis udvortes Selvstændighed, saa at
Forandringerne væsenlig beroe paa en Skiften Plads i Rummet. Det tomme Rum
er imellem dem, for saa vidt de da ikke berøre hinanden umiddelbart. De
have alle den fælles Egenskab at være tunge. Tyngdekraften virker paa
Afstand, og Afstanden er det Tomme. Saaledes er Indtrykket, naar vi
eensidig holde os til det blot Mechaniske. Men indtrængende Forskninger
angaaende Lysets, Varmens, Magnetismens og Elektricitetens mærkværdige
Eiendommeligheder have ført til den Erkjendelse, at Modsætningen mellem de
tunge Masser og det tomme Rum ikke er tilstrækkelig til at forklare
Phænomenerne, at der foruden den veielige Materie maa være en uveielig, og
at der er tiltrækkende Kræfter, forskjellige fra Tyngden. Physiken deler
sig da med Hensyn hertil i mechanisk og chemisk Physik.

Da den uveielige Materie ikke kan være Gjenstand for umiddelbar Erfaring,
er dens Antagelse naturligviis kun en Hypothese. Men med Physikens
Hypotheser forholder det sig paa en egen Maade, vi see det klart af de to
bekjendte Hypotheser om Lyset, Corpuscularhypothesen og Svingningshypothesen.

Newton antog, at det lysende Stof bestod af smaa Legemer, som Lysgiverne
udkastede med en overordenlig stor Hastighed, og som ved at træffe
Nethinden frembragte Lysindtrykket. Huyghens antog derimod, at
Verdensrummet var opfyldt af et Stof, der med en overordenlig ringe Tæthed
forener en overordenlig høi Grad af Elasticitet. Dette Stof er Ætheren, som
af de lysende Legemer sættes i Bølgebevægelse; det er Ætherbølger, som naar
de træffe Nethinden, frembringe Lysindtryk. Hverken Lysgran eller
Ætherbølger paavirkes af Tyngden. --- Skulde Striden om slige Hypotheser
afgjøres ved almindelige Ræsonnementsgrunde, da vilde det see betænkeligt
ud med Physiken. De, som angribe Theorien om Lysgranene, indvende, at man
hvert Øieblik maatte faae vildfarende Smaalegemer ind i Øiet og see Glimt
fra alle Kanter. De, som hylde Theorien, svare, at denne Vanskelighed let
undgaaes. Naar nemlig et lysende Punkt for hvert Secund sendte 100 Lysgran
ind i Øiet, vilde Øiet modtage et stadigt Lysindtryk; nu gjennemløber Lyset
40,000 Mile i Secundet, der vil da være 400 Mile mellem to paa hinanden
følgende Lysgran, saa at der altsaa ikke er Fare for Sammenstød. Derimod
maatte Ætherbølgerne, efter deres Mening svækkes underveis, da Ætheren
umulig kan være fuldkommen elastisk. Modstanderne af Lysgrantheorien kunne
igjen indvende, at Stødet af Smaalegemerne maatte slaae Øiet ud. Et Lysgran
af en Masse som $\frac{1}{2\,000\,000\,000}$ Gran vilde give et Stød som en
Kugle paa 25 Grans Vægt udskudt af en Riffel o.s.v. En saadan Strid for og
imod vilde blive lige saa endeløs som en Strid om Tilstanden i den anden
Verden. Men her viser det sig, hvilken Forskjel det gjør at stride om et
Overnaturligt, hvis Forbindelse med det Naturlige er uvist, og at stride om
noget Usandseligt, der dog hører under det Sandselige, og hvis Betydning
altsaa maa kunne prøves paa Kjendsgjerninger. Saaledes kom der omsider et
Vendepunkt, der blev afgjørende for Newtons Hypothese. I Conseqvens af
Hypothesen maa nemlig Lyset gaae $\tfrac43$ Gange hurtigere i Vand end i
Luft; da det nu ved Forsøg er godtgjort, at Forholdet netop er det
Omvendte, saa følger heraf, at Corpuscularhypothesen maa opgives. Derimod
har Svingningshypothesen, ifølge hvilken Lysstyrken beroer paa Udsvinget,
Farven paa Svingningstiden, hidtil viist sig fyldestgjørende ved
Forklaringen saavel af Lysets Udbredelse som af særlige til Reflexion,
Refraction, Interferens, Plansætning o.s.v.\ svarende Phænomener.

I logisk Forstand ere Ting og Begreb commensurable; men i empirisk Forstand
ere de incommensurable, thi Tingen er antilogisk. Heraf følger, at ikke
blot Hypothesen, men ogsaa Analogien maa ved Inductionens Anvendelse spille
en fremragende Rolle ved de physiske Iagttagelser. Analogien kan veilede;
thi to Phænomener, der i flere Henseender ere lige, vække Formodning om
Lighed i alle Henseender; men Analogien kan ogsaa vildlede, thi der kan
under den tilsyneladende Lighed skjule sig en væsenlig Ulighed, og omvendt.
I Læren om Varmens Forhold til Lyset have vi et Exempel. Det er en bekjendt
Sag, at Varmen kan udbrede sig i Straaler ligesom Lyset, at Varmestraalerne
brydes, og brydes forskjelligt ligesom Lysstraalerne, at de tilbagekastes
efter samme Lov: Udfaldsvinklen $=$ Indfaldsvinklen o.s.frmdl. Men Varme er
ikke Lys; Uligheden kommer tilsyne, blandt Andet deri, at mørke Legemer,
Legemer, som ikke udstraale Lys, kunne udstraale Varme. Analogien optager
imidlertid Ligheden i Uligheden, og den hypothetiske Forklaring er, at de
blotte Varmestraaler, om de end kunne gaae igjennem Øiets Vandmedium, dog
ikke kunne paavirke Synsnerven, idet man antager, at Bølgerne ere for brede
og Svingningerne for langsomme. Uligheden voxer, naar vi erindre, at Varmen
udbreder sig, ikke blot ved Straaling, men ogsaa ved Ledning, idet den ved
umiddelbar Berøring gaaer over fra det ene Legeme til det andet og breder
sig gjennem hele Massen. Til den høist charakteristiske Forskjel mellem
gode og slette Varmeledere svarer ingen saadan Forskjel mellem gode og
slette Lysledere. Men Ligheden beseirer igjen Uligheden; thi en Forklaring
er mulig, naar man antager, at Varmebølgerne, de svingende Ætherdele, der i
Legemet frembringe Varme, tillige bringe Legemets egne Smaadele i
Svingning. Antagelsen finder Bekræftelse i Læren om Varmefylden og
Atomvarmen. Den gjennemgaaende Prøvelse af en Analogi falder her sammen med
Prøvelsen af en Hypothese.

I Physiken er Loven høieste Instans, den sidste Grund, hvortil Alt maa
henføres; hvad er vel Grunden til, at Legemerne tiltrække hverandre paa
Afstand? Tyngden. Men hvad er Tyngden? En Lov! Hvad er Aarsagen til Lys og
Varme? Æthervirksomhed. Men hvad er Æthervirksomhedens Grund? En Lov! Hvad
er Aarsagen til de polære Forhold i de magnetiske og elektriske
Phænomener? Magnetismen og Elektriciteten. Men hvad er Magnetismens og
Elektricitetens Grund? En Lov! Fra dette Synspunkt viser sig en mærkelig
Overeensstemmelse og tillige en paafaldende Modsætning mellem Physik og
rationel Mechanik. Begge vise eenstemmig hen til Loven som Grundlag, begge
have den hypothetiske Dom: dersom A er B, saa maa C være D, til Ledetraad.
Men Brugen af den givne Ledetraad maa blive forskjellig, efter som
Erkjendelsen gaaer i apriorisk eller empirisk Retning. Apriorisk forstaaet,
ere Forholdsleddene selv Momenter i Forholdet; man kan være forvisset om,
at der i Forudsætningens hypothetiske Antagelse ikke indeholdes det
Allermindste andet end hvad Eftersætningens logiske Conseqvens udsiger.
Dersom det antages, hvad Aristoteles lærer, at umiddelbar Berøring er
ufravigelig Betingelse for Legemernes Indvirkning paa hinanden, saa maa
Forkastelsen af Newtons »Principia« deraf være en uafviselig Conseqvens.
Paa Physikens Omraade have slige Argumentationer derimod ingen Gyldighed.
Og hvorfor ikke? Fordi Forholdet mellem Led og Forhold er i det Empiriske
et ganske andet end i det Aprioriske. I Physiken, navnlig i den mechaniske
og chemiske Physik, have Leddene en saadan Selvstændighed, at de
Relationer, hvorpaa Lovene kjendes, maae sikkres ved omhyggelig Induction,
inden den hypothetiske Dom kan komme istand, end sige føre til en paalidelig
Slutning. Dersom Stenen falder hver Gang den berøves sin Understøttelse,
saa --- Ja, hvad saa? Saa maa det være, fordi den er passiv og dvask og
ikke kan bære sig selv. Nei, saa maa det være, fordi den hører til de
Legemer, hvis Natur det er at søge nedad mod Midtpunktet, mod Verdens evige
Midtpunkt. Nei, saa maa det være, fordi den i Vexelvirkning med Jorden
tiltrækkes af Jorden. Maa der allerede, for at sikkre sig Præmissen,
forudsættes et System af Inductioner, saa forstaaes heraf, at kun omhyggelig
Forskning kan lede til en paalidelig Slutning. De tilsyneladende analoge
Tilfælde ere ofte saa forskjellige, at Slutningen kan bedrage, hvad enten
man anvender \textit{modus ponens} eller \textit{modus tollens}. Den
ufuldstændige Induction er nemlig en Sandsynlighedsslutning, hvori
Sandsynligheden for, at Prædicatet gjælder for det almene Tilfælde
(Slægten) voxer, jo større det Antal særlige Tilfælde (Arter) er, hvorom
det gjælder. Da en saadan Induction imidlertid kan have sine »Instanser«,
saa kunne de ved Inductionen fundne Love ogsaa have deres Undtagelser;
saaledes har den Lov, at Varmen udvider Legemerne, en Undtagelse; saaledes
har Mariottes Lov en Undtagelse o.s.v. Vistnok er det stridende mod Loven
som Lov at have Undtagelser; men det heri Stridende hidrører fra, at det
empiriske Udtryk for Loven i det Hele er ufuldkomment.

Det er en Selvfølge, at Erkjendelsen af særlige Love maa være desto
fuldkomnere, jo mere man ved Betingelser og Mellembestemmelser er i Stand
til at udlede dem af eller henføre dem under en tilsvarende Grundlov;
hvorimod de antage en blot empirisk Charakteer, naar ingen saadan Deduction
eller Reduction er mulig. Ved Reduction og Deduction blive de særlige Love
aprioriserede. Jo almindeligere d.v.s.\ simplere den Betingelse er, som
Grundloven medfører, desto lettere gaaer det med de særlige Loves
Apriorisering; jo mere forviklede og skjulte de Betingelser ere, som
Grundloven selv indeholder, desto vanskeligere blive Forskningens Opgaver.
Hvor klar og simpel er ikke den mechaniske Physiks omfattende Grundlov,
Tyngdeloven, i Sammenligning med den chemiske Physiks, navnlig
Magnetismens og Elektricitetens, specifiske Grundlove! Medens hiin beroer
paa en eneste, for al veielig Materie fælles Egenskab, ere disse derimod
bundne til et Complex af afledede Egenskaber og skjulte Betingelser, som
vanskelig lade sig gjennemskue. Lovene i den mechaniske Physik ere da ogsaa
blevne aprioriserede efter en ganske anden Maalestok end Lovene i den
chemiske.

Hvad dette betyder, kan sees paa Forholdet mellem Keplers og Newtons
forskjellige Standpunkt. Den Keplerske Reflexion abstraherer Lovene af
foreliggende Tilfælde og sikkrer sig dem ved empirisk Induction; den
Newtonske Methode gaaer derimod ud paa at erkjende Aarsagerne af deres
Virkninger, bestemme det Særlige paa Grundlag af det Almene og begribe
Lovene i deres Opfyldelse. Paa Objectiveringens Vei et storartet Fremskridt!
Kepler taler kun om de Arealer, som beskrives i Planetbanerne; men Newton
beviser, at Loven for Arealerne maa gjælde for en hvilken som helst
Bevægelse om et fast Centrum, da den i Almindelighed er bestemmende for
Centralbevægelsens objective Begreb. Kepler holder sig umiddelbart til det
Factiske, at Planetbanerne ere Ellipser, men, bunden til umiddelbare
Kjendsgjerninger, er han ikke i Stand til at angive Aarsagen; af Newtons
Principer kan derimod Aarsagen sees i Virkningen, for saa vidt det kan
indsees, at ikke blot Ellipsen, men i Almindelighed Keglesnittet er en
naturnødvendig Conseqvens, en udtrykkelig Aarsagsconseqvens af
Tiltrækningens eiendommelige Lov $T = \dfrac{\mu}{r^2}$, og at omvendt denne
Lov er en nødvendig Conseqvens af Bevægelsen i et Keglesnit. Kepler
indskrænker sig til at opstille de tre Love ved siden af hverandre uden
nogensomhelst Deduction; men ifølge Newtons Principer kan den ene Lov, naar
man gaaer ud fra det Almindelige og indfører Betingelser, deduceres af den
anden. At Deductionen bestandig maa knyttes til givne Betingelser, minder
om, at Beviserne, hvad det Physiske angaaer, hverken ere blot aprioriske ei
heller blot empiriske, men at de, ifølge Sagens Natur, beroe paa en exact
articuleret Vexelbestemmelse af Erfaring og Tænkning. Det er ikke en
speculativ Construeren a priori; det er en grundig videnskabelig
Apriorisering, hvorom her tales.

Anderledes i Læren om Magnetisme og Elektricitet. Selv om man, hvad
tidligere er skeet, gik ud fra den Forudsætning, at nogle Legemer ere
magnetiske, andre elektriske, saa vilde man dog umulig ved at gaae ud fra
de til magnetisk og elektrisk Virksomhed svarende Grundphænomener kunne
opfatte disse saaledes, at man af dem ved at indsætte Betingelser kunde
aflede og forklare de særskilte Phænomener. Hvad er Magnetisme? Et Udtryk
for leddelende, polære Functioner. En Tiltrækning efter Tyngdens Lov
forudsætter i det Mindste to Led, to Masser eller Massedele; desuagtet er
heri ingen egenlig Leddeling, den ene Masse er som den anden, og enhver
Forskjel i Egenskaber, selv i Størrelse, er Tyngdens Væsen, almindelig
taget, uvedkommende. Anderledes med Tiltrækningen mellem to Magneter, hver
af dem har sine Poler, de uligeartede Poler tiltrække, de ligeartede
frastøde hinanden; Tiltrækning betinges af Frastøden og omvendt. I Polernes
Forhold er den ene Masse ikke blot udenfor den anden, men Andetheden er
accentueret to Gange, den ene Pol er den anden qvalitativt modsat; Sydpolen
har sit bestemte Andet i Nordpolen, og just derfor er Tiltrækningen
gjensidig; derimod har Sydpolen kun et formelt Andet i en anden Sydpol;
Eensformigheden er i Strid med Leddelingens Princip, og Frastødningen
gjensidig. Hvor oprindelig de modsatte Led forudsætte hinanden, sees i
Kjendsgjerningen deraf, at man vel kan skaffe sig Magneter med mange Poler,
men ikke en Magnet, som kun har een Pol. Hvad er Elektricitet? Analogien
med Magnetismen er iøinefaldende paa Grund af den polære Modsætning; men
hvad det indre Virksomme angaaer, maa Ligheden snarere synes at tilhøre det
Phænomenale end Væsnet. Magnetismen er nemlig som »Kraft« og »Egenskab«
bunden til det magnetiske Legeme, medens Elektriciteten »frit kan gaae over
fra det ene Legeme til det andet«, og ved denne Overgang viser sig som
analog med Varmen. Der er, hvad Ledningen angaaer, en saa mærkelig
Forbindelse mellem Elektricitet og Varme, at man saa vel for Metallerne som
for deres Legeringer har fundet de samme Forhold mellem deres Ledningsevne
for Varmen som for deres elektriske Ledningsevne. Dersom Varmen havde
Poler, vilde det ligge nær at opfatte Elektriciteten som polariseret Varme,
men da ingen Kjendsgjerning tyder hen paa, at der skulde være en positiv
Varmepol, der tiltrak en negativ og frastødte en anden positiv, da der ved
bunden og frigjort Varme forstaaes noget ganske Andet end ved bunden og fri
Elektricitet, da Varmen, med andre Ord, kun er qvantitativt, ikke
qvalitativt, forskjellig i sig selv, saa indsees heraf, at der maa være en
Grundforskjel mellem Elektricitet og Varme.

Hertil kommer, at de magnetiske og elektriske Egenskaber ere meddelelige og
variable, saa at Legemer under forskjellige Betingelser kunne blive snart
magnetiske, snart diamagnetiske, snart igjen ikkemagnetiske, eller snart
positivt, snart negativt elektriske, snart igjen ikkeelektriske. Det er
mærkelige Indblik, som herved aabnes os i de moleculære Forandringer.
Umiddelbart seet, tager det sig jo ud, som sad den magnetiske Kraft i
enhver af Magnetens Poler, medens der i det saakaldte Ligevægtsrum ingen
saadan Kraft var; men Leddelingen er gjennemgaaende. Ved at bryde en
Stangmagnet midt over, faaer man ikke to halve Magneter, hver med sin Pol,
men to hele, hver med to Poler. Delen er et fuldstændigt Hele, og Delingen
kan fortsættes indtil de elementære Enkeltheder: enhver Magnet er en
Totalitet af Elementærmagneter. Grunden til, at de virksomme Poler danne
sig henimod Endefladerne ligger deri, at Elementærmagneterne i
Ligevægtsrummet udligne og ophæve hverandres Virkninger, hvorved der
fremkommer en egen, af Forholdene betinget Modsætning mellem bunden og fri
Magnetisme. --- Hvor dybt Elektriciteten er impliceret i de forskjellige
Legemers Eiendommelighed, kan allerede sees af den Maade, hvorpaa den
vækkes. Den vækkes ved Gnidning, men Gnidningen gjør netop de legemlige
Forskjelligheder kjendelig; det er et betegnende Udtryk for Forskjellen
mellem Glasset og Lakket, at de to Legemer ved at gnides med det samme
Stof, Uld eller Silke, blive, det første positivt, det sidste negativt. Den
vækkes ved Fordeling. Er et Legeme A ladet f.\ Ex.\ med positiv
Elektricitet, vil det gjennem et isolerende Mellemrum virke paa et andet
Legeme B saaledes, at der ved Fordeling i B fremkommer negativ Elektricitet
paa de Dele af Overfladen, der ere A nærmest (Influenselektricitet af
første Art), men i den fra A fjærneste Overflade derimod positiv
Elektricitet (Influenselektricitet af anden Art). Hertil svarer saa igjen
den stærkere eller svagere Tiltrækning, thi da Fordelingen er størst i de
gode Ledere, maa Tiltrækningen til disse ogsaa være stærkest. Den
eiendommelige, ved Elektriciteten bestemte, Forskjel mellem gode, halvgode
og slette Ledere, Isolatorer, der er saa betegnende for Leddelingen, vil
Varmen ikke blot ændre, men ved høiere Grader endog udslette; man kjender
intet Legeme, som i Glødhede forbliver Isolator. Men naar Varmen ikke
virker hæmmende, men betingende, er Leddelingen saa udpræget, at endog den
blotte Berøring er nok til at vække Elektriciteten.

At Magnetismens og Elektricitetens specifiske Phænomener ere lutter
Lovopfyldelser, derom kan Ingen tvivle; men Forholdet mellem det
Individuelle og det Almene er saa articuleret og desaarsag saa indviklet,
at Forskeren med al sin Kunst hvert Øieblik taber Traaden i den store
Labyrinth. Hvorpaa beroer egenlig Vexelvirkningen mellem det Elektriske og
det Magnetiske i Elektromagnetismen? Ørsteds Opdagelse sprang frem af en
Kjendsgjerning som et Blus i Mørket; men hvor er Phænomenets Grund, i hvilke
Betingelser ligger det Virksomme selv? Hvad er det Strømmende i den
elektriske Traad? Hvorfor vil Magnetnaalens Nordende gjøre et Udslag mod
Øst, naar den befinder sig over en Strøm fra Syd til Nord, og derimod et
Udslag mod Vest, naar den selv befinder sig under Strømmen? Hvad er Grunden
til, at to parallele Strømledere tiltrække hinanden, naar Strømmene gaae i
samme Retning, men frastøde hinanden, naar de gaae i modsat Retning?
Ampères Theori, at enhver Magnet kan betragtes som et System af elektriske
Strømkredse, Solenoider, og Magnetismen altsaa betragtes som et
eiendommeligt Phænomen af Elektriciteten, synes, mærkeligt nok, at passe
paa Jordmagnetismen, men ikke paa de legemlige Magneter, eftersom det ikke
kan antages, at elektriske Strømme skulde kunne vedligeholde sig, endog i
gode Ledere, uden enten at afgive eller modtage noget Arbeide.

Da en Apriorisering kan, ifølge Sagens Natur, gaae langt raskere frem i den
mechaniske Physik end i den chemiske, saa kunde det ved første Øiekast see
ud, som om Realbegrebernes System i hiin laa Realiteterne langt nærmere end
i denne; men Forholdet er netop det Omvendte. Den mechaniske Physiks Love
have et langt abstractere Præg, end Tilfældet er med Lovene i den chemiske
Physik. Jo dybere Lovene selv indgaae i Concretionerne, d.v.s.\ jo rigere
Leddenes specifiske Qvalitet og Forholdenes Specification reflectere
hinanden, desto klarere vil det vise sig, at Lovene selv ere
Realbegrebernes almeennødvendige Udtryk. Naar det desuagtet seer ud, som
var det Physisk-mechaniske i sig selv mere begribeligt end det
Physisk-chemiske, da er Grunden simpelhen denne, at de physiske
Concretioner ere antilogiske, og at det Antilogiske altid maa blive desto
mere fremtrædende, jo mere indviklede Concretionerne ere, medens den
empiriske Erkjendelse, for saa vidt den i Sandhed er videnskabelig, om dens
Interesse for Kjendsgjerninger end er nok saa stor, altid er henvist til
det mere Abstracte og derved til det Logiske.

\section{Chemi og Mineralogi}
\label{sec:15}
% pp. 151--160

\noindent Ved alle hidtil betragtede Phænomener af Leddeling og polær
Modsætning holde de legemlige Led sig dog endnu udenfor hverandre:
Energierne ombyttes, Spænding omsættes i levende Kraft, og levende Kraft i
Spænding, Varme i Elektricitet, og Elektricitet i Arbeide; men under alt
dette hævde de paagjældende Legemer dog deres relative Selvstændighed.
Først i Chemismen bliver den leddelende Virksomhed saa gjennemgribende, at
Leddene sætte Selvstændigheden til i deres Affinitet: den concret-legemlige
Existens opløses for at frembringes, og frembringes for igjen at opløses.
Ligesom den mechaniske Physik danner Overgang til den chemiske, saaledes
danner den chemiske Physik igjen Overgang til Chemien.

\emph{Chemi.} Hvad er i en Voltastøtte den oprindelig virkende Aarsag? Er
det den elektriske Strøm, der er Aarsag til den chemiske Virksomhed, eller
er det omvendt den chemiske Virksomhed, der er Aarsag til Strømmen, eller
er her maaske en eiendommelig, ligelig Samvirken af begge Factorer? Det
Sidste er vistnok paa Videnskabens nuværende Trin det Sandsynligste. Ved
Voltastøtten er Vandet blevet adskilt. Ligesom Vandet, saaledes adskilles
alle sammensatte Stoffer, naar de vel at mærke lede Elektriciteten og i det
mindste tildeels ere i smeltet eller opløst Tilstand. Den elektriske Lov
siger, at Vægtene af de ved samme Elektricitetsmængde adskilte Stoffer
forholde sig som disse Stoffers Æqvivalenttal. Idet Physiken fra alle
Sider stræber at gjøre sig Rede for Legemernes molekulære Tilstande og de
mulige Forandringer, som disse undergaae, har man forklaret sig
Adskillelsen af Vandets Bestanddele, som jo hver for sig kunne opsamles og
maales i et Voltameter, paa følgende sindrige Maade: Naar Brint er
forbunden med Ilt til Vand, blive ved den inderlige Berøring mellem de
mindste Dele Iltatomerne negativt og Brintatomerne positivt elektriske;
men paa Grund af den eensformige Fordeling af Iltens og Brintens Smaadele
viser Forbindelsen ingen fri Elektricitet. Der ligger nu --- kan man
forestille sig --- i en given Linie et vist Antal Vanddele mellem begge
Poler af en galvanisk Kjæde. Hver Vanddeel bestaaer af eet Atom Ilt og et
Dobbelatom Brint. Den positive Pol virker paa den første Vanddeel saaledes,
at Iltatomet, der er negativt, bliver tiltrukket, Brintatomet, der er
positivt, derimod frastødt. Men den første Vanddeel virker igjen paa den
anden, denne paa den tredie og saa fremdeles. Naar nu Tiltrækningen, som
den positive Pol udøver paa Iltatomet i den første Vanddeel, er
tilstrækkelig stor, rives dette bort fra sit Brintatom; men det fraskilte
Brintatom forbinder sig strax derpaa med Ilten i den anden Vanddeel; det
derved udskilte Brintatom trækker endvidere Ilten fra den tredie Vanddeel
o.s.v. Hele Strækningen igjennem foregaaer der en vedvarende Adskillelse og
Sammensætning af Vandet, kun ved selve Polerne kunne dets Bestanddele blive
frie.

Den indre Forbindelse mellem Elektricitet og Chemi er tydelig viist i den
elektrochemiske Theori. Ligesom Zink og Kobber i Berøring med hinanden faae
modsatte Elektriciteter, saaledes faae, ifølge Theorien, Atomerne af to
Grundstoffer modsatte Elektriciteter. At Atomerne ikke i og for sig ere
elektriske, men først blive det ved gjensidig Berøring, er en
Forestillingens Forudsætning, som forklarer den Kjendsgjerning, at et og
samme Stof kan, eftersom Forbindelsens modsvarende Led er, være snart
positivt, snart negativt. Saaledes danner Svovl i Forbindelse med Ilt det
elektropositive, men i Forbindelse med Brint det elektronegative Led. Ved
Grunddelenes indbyrdes Berøring bliver Elektriciteten bunden, saa at
Legemet overalt i sin indre Sammensætning ligelig er baade positivt og
negativt; men Leddelingen gjentager sig i stedse concretere Former. Naar
det ved Forbindelsen frembragte Legeme under de fornødne Betingelser
bringes i chemisk Forhold til andre Legemer, ophører Neutraliteten.
Grundstofferne forene sig parviis til binære Forbindelser, Forbindelser af
første Orden; de binære Forbindelser forbinde sig igjen parviis til
Forbindelser af anden Orden og danne Salte, i hvilke den ene af
Forbindelserne, Syren, er det elektronegative, den anden, Basen, det
elektropositive Led; Salte forbindes igjen til Dobbelsalte o.s.frmdl.
Affiniteterne have forskjellig Styrke.

Det er imidlertid ikke nok at tale i Almindelighed om stærkere og svagere
Affiniteter; her kræves et Maal for Affiniteternes absolute Størrelse:
dette Maal er Varmeudviklingen. Saaledes komme vi til den thermochemiske
Theori. Ved samme chemiske Proces udvikles der altid lige stor
Varmemængde, naar de ydre Forhold ere de samme. Varmeudviklingen, som
ledsager Dannelsen, er stedse lige stor, enten Forbindelsen eller
Adskillelsen skeer paa een Gang eller gradviis. Ved Adskillelsen af en
chemisk Forbindelse frembringes der altid lige saa megen Varme, som der
udvikles ved Forbindelsens Dannelse. Naar man vil udrette et Arbeide, maa
Arbeidsmængden være større eller idetmindste lige saa stor som Arbeidet.
Vil man adskille en chemisk Forbindelse ved den Arbeidsmængde, her
Varmemængde, som Dannelsen af en anden chemisk Forbindelse viser sig at
udkræve, saa maa denne Arbeidsmængde være større end den, som udfordres til
Adskillelsen af den første Forbindelse.

Nærmere beseet, er det dog heller ikke nok at have et Maal for
Affiniteterne; thi hvad der i Sandhed maa vække Forundring --- Affinitetsstyrken
svarer ingenlunde ligefrem til det Forhold, hvori Stofferne mættes. Chlor
og Brint have en stor Affinitet til hinanden, Chlor og Kulstof en
forholdsviis ringe. Ikke desto mindre mættes eet Brintatom med eet
Chloratom, medens eet Kulatom først er fuldkommen mæt, naar det har forenet
sig med fire Chloratomer. Saa fremkommer Atomitetslæren, Læren om »Valens«,
»Affinivalens«, »Qvantivalens«; hvoraf sees, at Chemiens Methode maa
forgrene sig i forskjellige, af hverandre mere eller mindre afhængige
Bimethoder. Da eet Atom Chlor forener sig med eet Atom Brint til
Saltsyreluft, eet Atom Ilt med to Atomer Brint til Vand, eet Atom Kvælstof
med tre Atomer Brint til Ammoniak, eet Atom Kulstof med fire Atomer Brint
til let Kulbrinte, saa kaldes, ifølge heraf, Chloret ligesom ogsaa Brinten
eet-atomet (monovalent), Ilten to-atomet (divalent), Kvælstoffet tre-atomet
(trivalent) og Kulstoffet fire-atomet (tetravalent). I Forbindelsen
$\mathrm{CH}$ har Kulstoffet altsaa tre, i $\mathrm{CH_2}$ to, i
$\mathrm{CH_3}$ een Atomitet fri; dette kaster et Lys paa de isomere og
polymere Dannelser, der iblandt Kulstofforbindelserne ere saa hyppige.

De chemiske Processer ere legemdannende Virksomheder, saavel med Hensyn
til Adskillelse som til Sammensætning Processer af concret Legemlighed.
Concretionsanalysen maa desaarsag optage en saa stor Mangfoldighed af
uligeartede Elementæranalyser i sig, at man umulig med Comte kan betragte
Chemien som fundamental i Modsætning til Mineralogien. Mineralogi og Chemi
betinge hinanden gjensidig, om vi end maae indrømme, at den sidste har en
bredere Basis end den første. Ingen Legemdannelse er i physisk Forstand at
opfatte som en enkelt Virkning af en enkelt Aarsag, men ethvert Legeme
fremkommer, være sig som Educt eller som Product, af forskjellige
Aarsagers Samvirken. Som det gaaer med Dannelserne, saaledes ogsaa med
Omdannelserne. I den Vexelvirkning af Cohæsion og Varme, hvorved
Tilstandsformerne bestemmes, kan det Mechaniske, det Dynamiske, det
Magnetiske, Elektriske og Elektrochemiske paa de forskjelligste Maader
deeltage.

\emph{Mineralogi.} Medens Chemikeren nærmest seer paa den legemdannende
Virksomhed, begynder Mineralogen umiddelbart med Dannelsens Produkt,
hvorhos tillige Tilstandsformen og den af samme betingede Krystalform
kommer i Betragtning. Krystalformen er et physisk Resultat af et Legemes
Overgang fra flydende til fast Tilstand, idet Smaadelene, ifølge Legemets
chemiske Charakteer, afleires paa bestemt Maade. At der i denne Henseende
ikke kan være Tale om et særeget Krystallisationsprincip, sees deraf, at
der ved Overgangen lige saa vel kan fremkomme en amorph eller blot
krystallinsk Dannelse, som en udformet Krystal; i denne Henseende herske
Betingelserne. For at bringe Legemerne til at krystallisere, kan man
anvende forskjellige Fremgangsmaader --- Opløsning f.\ Ex.\ i Vand,
Smeltning og Afkjøling, Fordampning og Fortætning af Dampe efter de
paagjældende Legemers Beskaffenhed; men Betingelsernes Indflydelse kjendes
dog i det Hele paa den Regel, at Krystallerne blive desto større og
smukkere, jo langsommere og roligere Krystallisationen foregaaer.

Da Krystalformen dannes af plane Flader, kunne saadanne mathematiske
Legemer, som Terningen, Qvadratoctaedret, Rhomboedret o.s.v.\ vælges til
Grundformer; men hvad det ved en Systemdannelse især kommer an paa, er
Axernes Antal, Stilling og Størrelse; hvor vidt der i en Gruppe haves f.\
Ex.\ tre lige store Axler, enhver enkelt lodret paa de andres Plan; eller
kun to vandrette Axer ligestore, den lodrette, Hovedaxen, derimod større
eller mindre; eller ingen af de tre Axer ligestore o.s.v. Hvad det Physiske
angaaer, kunne to forskjelligartede Substanser krystallisere i samme Form
(Isomorphi) og een og samme Substans i forskjellige Former (Dimorphi,
Polymorphi).

Forstaae vi ved Moleculærconstitution den Maade, hvorpaa Smaadelene i et
Legeme ere begrændsede, leirede og forbundne, saa vil man ogsaa forstaae,
at her maa være en eiendommelig Sammenhæng mellem Atomconstitution og
Krystalform. Hvad de til Grundstofferne svarende Enkeltatomer angaaer, kan
Polymorphi jo fremkomme derved, at de samme Atomer under flere forskjellige
Krystalliseringer leire sig forskjelligt; paa den anden Side lader
Isomorphien sig forklare deels, hvad undertiden eensidig fremhæves,
ligefrem deraf, at Atomerne i to eller flere forskjellige Grundstoffer ere
lige store og lige dannede, deels deraf, at Atomerne, skjøndt ulige store,
dog i Figur og Størrelse forholde sig saaledes, at en Gruppe af de mindre
Atomer kan tænkes congruent med det enkelte større. Hvad dernæst de til
sammensatte Legemer svarende Atomsammensætninger angaaer, er Mulighedernes
Antal desto større, jo flere Alternativer de paagjældende Sammensætninger
indeholde. For at to Legemer skulle i Ordets egenlige Forstand kunne
krystallisere sammen, maa det vistnok forudsættes, at de have samme
Krystalform; men Lighed i Krystalform er dog ikke uden videre Grund til, at
de krystallisere sammen: hertil kræves i Reglen Lighed i chemisk
Sammensætning. En Salmiakkrystal f.\ Ex.\ har samme geometriske Form som en
Alunkrystal; desuagtet krystallisere de dog i een og samme Vædske adskilte
fra hinanden; men see vi paa Bestanddelene, saa har Salmiakmolekulen da
ogsaa kun to, Alunet derimod tredive sammensatte Atomer.

Henregner man nu til Mineralier ikke blot de af ligeartede Dele bestaaende
uorganiske Stoffer, der danne den faste Deel af Jordkloden, men endog
adskillige i Jorden forekommende Stoffer af organisk Oprindelse, saasom
Steenkul, Bruunkul, Rav og Infusoriekisel, saa kan man da med Grund sige,
at Mineralogi og Chemi ere Concretionsvidenskaber i særegen Forstand,
Videnskaber, som optage Resultater af de foregaaende og forene Alt til en
indtrængende Erkjendelse af den concrete Legemlighed. Som Mineral har jo
Legemet, foruden sit chemiske og krystallographiske Charakteermærke, baade
mechaniske, optiske, thermiske og elektriske Egenskaber; et Mineralsystem,
som blot holder sig til ydre, »naturhistoriske« Kjendemærker,
tilfredsstiller ikke; et System, der alene tager Hensyn til Mineralets
chemiske Sammensætning, er ogsaa eensidigt: her maa, som i det »blandede«
System, overalt tages Hensyn til det Concrete.

Hvad Concretionsmethodens videnskabelige Værdi angaaer, indstille sig to
eensidige Idealer: ifølge det ene gaaer Videnskaben ligefrem ud paa
væsenlig Indsigt; ifølge det andet gaaer den fornemmelig ud paa
Erhvervelsen af Kundskab og Kjendskab til virkelige Ting. I første
Henseende fristes man til at sætte den absolute Idealismes Logik høiere end
Astronomien, som med sine abstractere Forholdsled og almindeligere
Loverkjendelse dog er empirisk, ja langt høiere end Chemi og Mineralogi,
som med de mange Specialundersøgelser er ganske fortabt i det Empiriske; i
sidste Henseende vil man kaste Vrag paa Ideelæren, som ikke beriger os med
Kundskaber, og sætte Mineralogien, som med stedse dybere Indtrængen i
Legemernes indre Bygning skaffer os meget at vide om deres virkelige
Beskaffenhed, ikke blot over Logiken, men endog over Astronomi og
almindelig Mechanik. Slige Vurderingsmaader ere imidlertid vage og
vilkaarlige. Den eneste sande Maalestok ved Bedømmelsen af det i
Videnskaben Videnskabelige maa, som ovenfor viist, søges i Forholdet mellem
Begrebet, Loven, og Tingen, eller den concrete Legemlighed.

Dersom det var muligt at udlede Tilværelsens Love af rene Begreber, af
Lovene de existerende Tings almindelige Egenskaber og af disse igjen den
virkelige Existens, da vilde Ideelæren ganske vist optage de øvrige
Videnskaber i sig; men det lader sig ikke gjøre. Dersom det var muligt, at
sandselig Iagttagelse og empirisk Induction ved sig selv kunde føre til
væsenlig Indsigt i Tingenes Grund, da maatte vi give dem Ret, som med
Ringeagt for dybere gaaende Begrebsanalyser udelukkende holde sig til
Kjendsgjerninger; men det lader sig ikke gjøre. Kun fra de tre
Udgangspunkter, fra Realbegrebet, Loven og den concrete Legemlighed, kan
Videnskaben arbeide sig frem til begrundet Erkjendelse af den virkelige
Sandhed. En Begrebsudvikling, som hverken har Rødder i det Virkelige eller
sigter paa bestemt Virkelighedserkjendelse, er trods al
Reflexionsvirtuositet gold og ufrugtbar, og en empirisk Iagttagelseskunst
med endeløse Analyser mangler sand videnskabelig Bevidsthed. Positivismen
er paa sin Vis lige saa eensidig som den objective absolute Idealisme paa
sin; fra begge Sider driver man Smugleri med tilsnegne Vare: det er en Art
videnskabelig Toldsvig. I klar Erkjendelse af Vexelforholdet mellem det
Logiske i Realbegrebernes Analyse og det Empiriske, det Antilogiske, i
Analyserne af den concrete Legemlighed ligger det Videnskabelige.

\section{Geognosi og Geologi}
\label{sec:16}
% pp. 160--172

\noindent Som et Hele med sine Dele frembyder Jordlegemet to forskjellige
Synspunkter for videnskabelig Opfattelse. For saa vidt det gjælder om at
undersøge Delene, hvoraf det Hele er sammensat, er Synspunktet geognostisk,
for saa vidt der lægges an paa en Fremstilling af det Heles Udvikling i og
med Delenes, er Synspunktet geologisk. Vil man end ikke stille Geognosi og
Geologi mod hinanden som to særskilte Videnskaber, saa er det dog
umiskjendeligt, at der ved det Geognostiske lægges afgjørende Vægt paa den
egenlige Forskning, medens det Geologiske mindende om det Geogoniske
nærmest falder ind under Hypotheser og Theorier. Geognosiens Methode er
analytisk, Geologiens synthetisk.

Sammenfatte vi alle Dele af Naturlæren, der høre med ved Undersøgelsen af
Jordlegemets Skikkelse og Sammensætning, under Fællesbenævnelsen Geognosi,
saa er Geognosi med Astronomi, physisk Geographi, chemisk Physik, Chemi,
Mineralogi og Lithologi en saa vidt forgrenet Videnskab, at den vel kan
siges at omfatte hele den jordisk-elementære Natur. Kan denne Videnskab da
i egenlig Forstand lære os, hvad Natur er? Nei! Den undersøger
Naturtingene, men ikke Naturen. Vistnok lærer den, at ethvert eiendommeligt
Legeme har sin eiendommelige Natur, og at den Virksomhed, hvorunder
Legemerne dannes, grupperes, forandres, er en gjennemgaaende
Naturvirksomhed; men Naturen selv i de uligeartede Naturer undersøger den
ikke, og det af den gode Grund, at Naturen selv hverken er en Ting for os
med Egenskaber ei heller en Ting i sig uden Egenskaber, men et Realbegreb
eller rettere: det System af Realbegreber, der uafbrudt ved Magtens Udslag
i lovbestemt Sammenhæng (jvnfr.\ §\ 11) kommer til Existens i
Naturtingene.

Kritisk forstaaet, gaaer det Geognostiske umærkelig over i det Geologiske,
men, væsenlig seet, er Overgangen dialektisk. Her er nemlig den
gjennemgaaende Forskjel, at Geognosien, der begynder med Delene og bevæger
sig mellem Udstykning og Sammenstykning, paa alle Punkter strængt maa holde
sig til den physiske Opfattelse af Aarsag og Virkning, af Forudsætning og
Resultat; hvorimod Geologien, der gaaer ud fra Heelheden og seer paa den
fremadskridende Udvikling, har Vanskelighed ved at udelukke Forestillingen
om lavere og høiere Udviklingstrin, og desaarsag, i Fremskridtets Navn,
hvert Øieblik er fristet til at omsætte Resultater i Formaal. Det er
naturligviis kun almindelige og ubestemte Antydninger, Physiken kan give om
Jordklodens allerførste Tilstand; men allerede disse Antydninger ere
tilstrækkelige til at vise Forholdet mellem det Geognostiske og det
Geologiske.

Gaae vi ud fra den Antagelse, at Maanen maa være dannet af Jordmassen paa
tilsvarende Maade som denne af Solmassen, saa fører Betragtningen af
Maanens Dannelse som forestillet Resultat tilbage til Forestillingen om de
nødvendige Forudsætninger. Den med Maanen svangre Jord maa da have havt et
anderledes Rumfang end den nu har; medens Jordklodens Radius for Tiden kun
er 860 Miil, maa den i hiin Tid have været saa omtrent 60,000 Miil. Før
Maanen kom til Verden, maa den moderlige Jord da ogsaa have dreiet sig
temmelig langsomt om sin Axe: Døgnet, som nu kun er 24 Timer, maa da have
udgjort 27 Dage. Der kan endvidere anføres Adskilligt til Forklaring af
Jordens Fladtrykning, Fortætning, forstørrede Dreiningshastighed o.s.v. Med
alt dette bevæger Betragtningen sig i en Cirkel af forestillede Phænomener.
Hele Begrundelsen gaaer ud paa at afpasse forestillede Forudsætninger efter
et blot forestillet Resultat; thi Maanens Dannelse af Jorden er jo for
Betragteren kun en forestillet, ikke en virkelig Begivenhed. Men om vi
endog antage, at Begivenheden har en saa stor Sandsynlighed for sig, at den
kan betragtes som Kjendsgjerning, er der saa herved egenlig betegnet noget
Fremskridt? Ingenlunde, --- at sige i geognostisk Forstand. Hvis her i
denne Sammenhæng skal tales om et Dannelsens Fremskridt, maa Betragteren
underskyde sin egen subjective Forestilling. Det Eneste, der fra et reent
objectivt Synspunkt kan sees i alle Begivenheder, alle Dannelser og
Forandringer, er »de uforanderlige Loves« Opfyldelse. Men at der med Hensyn
paa Lovenes Opfyldelse ikke skeer noget Fremskridt, er fra Physikens Side
uimodsigeligt: i Jorden, den roterende Dampkugle, bleve Lovene lige saa
fuldstændig opfyldte før, som efter Maanens Dannelse.

Jorden, hedder det, var dernæst en ildflydende Masse, og alt som
Afkjølingen foregik, idet Varmen udstraalede i Verdensrummet, begyndte der
at danne sig en fast Skorpe paa dens Overflade, omtrent som et Isdække
danner sig paa Vandet; men da Ebbe og Flod, hidrørende fra Maanens og
Solens Tiltrækning, fremkaldte bestandige Bevægelser i den fuldkommen
flydende Jordmasse, kunde denne Skorpedannelse ikke gaae for sig uden
stadige og store Forstyrrelser. Mangen Gang blev det tynde Dække sprængt,
indtil Jordens bestandig fremskridende Afkjøling frembragte et saa mægtigt,
af sammenbagte Steenflager bestaaende Dække, at dette idetmindste ikke
oftere kunde sprænges i Stykker, men vel kunde briste nu paa et Sted, nu
paa et andet. Ved at følge Beskrivelsen i det Enkelte og særlig betragte,
hvorledes Jordskorpen som et første raat Udkast bestandig omdannes og
beriges, idet den gjennembrydes af Eruptivmasser --- Graniter, Porphyrer,
Grønstene o.s.v.\ --- der i Sammensætning og Vægtfylde ere forskjellige,
efter som de komme fra en forskjellig Dybde; ved at see, hvorledes Vandene
under deres ødelæggende Indflydelse paa Overfladens Bestanddele selv
arbeide med paa videre gaaende Dannelser; hvorledes de i Vandene opløste og
derefter afleirede Masser forøge Skorpens Tykkelse her og der, ved denne
Forøgelse hindre Varmeudstraalingen og bevirke, at en Deel af Skorpen fra
neden af atter smeltes, saa at der ved Ophedningen dannes krystallinske
Skiferstene, Gneus, Glimmerskifer, Hornblendeskifer o.\ dsl.; ved, som
sagt, at fordybe os i disse første Dannelser, faae vi uvilkaarlig Indtryk
af, at her ved alt dette tilstræbes Noget; det er for Indbildningen --- som
bleve vi Vidner til en begyndende Udførelse af en storartet Byggeplan; som
gjorde vi Indblik i dunkle Værksteder, hvor en usynlig Mester arbeider med
at tildanne Materialet, dække over Afgrundene, lægge Steenbjælker og bygge
Bjerge. Men geognostisk talt --- er det en Phantasi. Her forberedes Intet;
thi her tilsigtes Intet; det Hele er en Række Forandringer, hvori Lovene
opfyldes: det gaaer Altsammen for sig saa naturligt og ligefrem, som naar
et Isdække paa Vandet danner sig i Blæstveir.

Tager det geologiske Sprog undertiden et halvpoetisk Opsving, saa kommer
Geognosien strax med en prosaisk Afkjøling. Man aabner sig Urverdenens
Archiver, betragter de gamle Steenmonumenter, taler om Billedskrift, om
Hieroglypher, hvis Mening det gjælder at tyde. Hvad er da Meningen med disse
cyklopiske Muurværker, disse gigantiske Granitblokke, denne Opstabling af
Euriter, Porphyrer, Gneuser, Graavakker og hvad de alle hedde? »Spørg ikke
derom, men forsk! Naturen selv har ingen Mening om Sligt, og paa
Menneskenes Meninger kommer det slet ikke an.« --- Men dersom der ingen
Mening er i den dunkle Skrift, hvorledes vil man da kunne tyde den? »Forsk,
forsk, thi Nøglen ligger skjult i Sammenhængen.«

Undersøge vi nu, hvorledes Jordudviklingen tager sig ud, naar vi særlig
betragte den under Heelhedens Synspunkt, da komme vi til en egen Opfattelse
af Forholdet mellem det Videnskabelige og det Teleologiske. Det maa paa
Forhaand udtrykkelig fremhæves, at hvad vi her opfatte som teleologisk,
ikke har noget med supranaturalistiske Forestillinger om et udenfor
staaende Væsens personlige Hensigter eller vilkaarlige Indgriben i Tingenes
Sammenhæng at gjøre, men udelukkende dreier sig om naturlig Selvdannelse,
om Anlæg og Formaal, Realbegreb og Existens. De tre storartede Hypotheser:
den kosmologiske, om Planetsystemets Oprindelse, den geologiske, om
Jordklodens Oprindelse, og den biologiske, om de levende Organismers
Oprindelse, afgive Exempler paa, hvad vi i videnskabelig Forstand maae
kalde naturlige Selvdannelser.

Den physiske Astronomi forudsætter, at Materien med dens Kræfter er »evig«
som Rummet, og lægger saa an paa: ud fra det første Værende at begribe den
Vorden og trinvise Dannelse, som det Nuværende i Tidernes Løb maa have
gjennemgaaet. Det første Værende er at opfatte som en Tilstand, i hvilken
alt Tilværende fra Begyndelsen kan antages at have befundet sig. Ved en
omfattende Reduction føres vi da tilbage til en i Rummet ligelig udbredt,
endnu indifferent Materie, en overordenlig tynd, men uhyre Dampmasse, hvis
Tvermaal af Nogle anslaaes til 12,000 Billioner Miil. Dannelsen begynder nu
med, at Delene tiltrække hverandre, hvorved der fremkommer Fortætning,
Rotation og Varme. Dampmassen bliver kugleformet, ophøiet i Eqvator;
Rotationen tiltager, og Centrifugalkraften voxer: altsammen naturligt. I
det ophøiede Bælte af Kuglemassen vil der saa efterhaanden danne sig Ringe,
som løsne sig, afsondres og igjen sammenrulles; de sammenrullede Masser
bevæge sig om deres Axe og om Hovedmassen: her have vi et første Udkast til
Systemet i al Naturlighed. Hvad der skete med den store Masse, kan gjentage
sig med de mindre; nogle af disse kunne igjen afsætte een eller flere
Ringe, disse Ringe kunne afsondres og blive til Maaner o.s.v. --- Det er
ikke Hypothesens Fortrin og Mangler i physisk Henseende, vi her skulle
prøve: vor Opgave er at betragte den i Forhold til det Teleologiske.

Det er indlysende, at der ifølge en saadan Forudsætning ikke kan paavises
nogetsomhelst enkelt Phænomen, der særlig skulde vidne om Plan og Hensigt.
At Planetbanerne kun have en ringe Heldning, at Omløbstiderne ere
incommensurable --- $\dfrac{T}{t} = \dfrac{A^{\frac32}}{a^{\frac32}}$ ---,
at Solen, hvad der med Hensyn til »Perturbationerne« jo er særdeles vigtigt,
fik en saa overveiende stor Masse, at den vel omtrent er 700 Gange større
end de Masser tilsammen, der bevæge sig omkring den, er vistnok et talende
Vidnesbyrd om lovbestemt Orden, men ingenlunde noget objectivt Beviis for
en særegen Skaberviisdom. Kun en eneste Omstændighed kunde synes at tyde
paa en saadan Viisdom. Af de tre Keglesnit er kun det ene, nemlig Ellipsen,
egnet til at give et sluttet System af Bevægelser; de to øvrige vilde
derimod gjøre Systemet meningsløst, idet Planeterne, naar de bevæge sig
enten i en Parabel- eller i en Hyperbelgreen, efterhaanden fjerne sig fra
Solen, uden nogensinde at vende tilbage. Vil man nu ikke skrive det paa
Tilfældets Regning, at der er Fornuft i Solsystemet, saa synes det
unægtelig, som om netop den særegne Begyndelseshastighed, af hvilken de
elliptiske Baner ere betingede, maatte forudsætte en særegen Viisdommens
Foranstaltning og altsaa indeholde Svaret paa Laplaces Fordring:
\textit{Philosophe, montre moi la main qui a jeté les Planètes sur la
tangente de leur orbite.} Men Svaret løser sig op i et Ræsonnement af
Ukyndighed. Den rationelle Mechanik beviser paa en exact tilfredsstillende
Maade saa vel, hvorfor Parablen, med en Sandsynlighed $= \dfrac{1}{\infty}$,
er blevet udelukket, som ogsaa, hvad Følgen vilde være blevet, hvis der
havde været Begyndelseshastigheder, som medførte Hyperblen: Planeter med
Hyperbelbaner vilde i Urtiden have forladt Systemet, og vi altsaa beholde
dem, der fik en ringere, for Ellipsen egnet, Begyndelseshastighed tilbage.

Imidlertid er her dog en Dunkelhed, der maa opklares, en Tvetydighed, vi
maae see at faae hævet. Hvorledes forholder det sig egenlig med Hypothesen
om de evige Atomer? Ere de evige, saa ere de jo absolut selvstændige. Men
imellem absolut selvstændige Atomer kan der jo ikke gives noget gjensidigt
Afhængighedsforhold. Have Atomerne Egenskaber og Kræfter, hvorved de kunne
virke paa hverandre, saa ere de ikke uforanderlige, selvfølgelig heller
ikke evige; ere de evige, i den Forstand, at de ere, fordi de ere, saa tør
vi jo hverken tillægge dem Kræfter eller Egenskaber. Om vi end ikke kunne
trænge ind til »de inderste Aarsager« eller udforske Tingenes første
Begyndelse, maae vi dog lægge an paa at befrie vore Forestillinger om det
Oprindelige fra vitterlige Selvmodsigelser. At der er noget Misligt ved
Hypothesen om de evige Atomer i den evige Urdamp, har Physiken let ved at
overbevise sig om. Det uendelige Rum kan ikke have været fyldt med Atomer,
thi det uendelige Rum kan aldrig blive fyldt. Men sæt endog, at det
umaalelige Verdensrum fra Evighed af har været fyldt med Urdamp: hvorledes
skulde denne Urdamp saa være kommet i Bevægelse? --- »Ved Tiltrækning!«
Her er jo blandt de evige Atomer ingen Tiltrækning mulig. --- »Til
Centrum!« Her er jo intet Centrum. --- Jo længere man fortsætter disse
Overveielser, desto klarere indsees det, at hvad her søges, er mere end
Masser, Atomer og Kræfter, at det er Naturen selv, Tilværelsens evige
Grund. Men hvad er saa egenlig det, vi kalde Natur? En oprindelig Eenhed af
Begreb og Existens. Et Begreb, om end et Alt opfattende Realbegreb, kan
ikke være Naturen, thi Realiteten mangler; et Indbegreb af alle Slags Ting
og Realiteter kan Naturen heller ikke være, thi Grunden mangler. Naturen
maa altsaa være den oprindelige Eenhed af Grund og Realitet i Grundens
Form. Men i den oprindelige Grund, der har alle Betingelser i sig selv,
efterdi den er Grunden til Alt, har det hele Indbegreb af vordende
Existenser ikke blot sine Muligheder, men, i væsenlig Forstand, sit Anlæg.

Planetsystemets Udvikling er altsaa en Anlægsudvikling, dets Dannelse en
naturlig Selvdannelse, en Dannelse af et Alt præformerende Naturanlæg. Den
saaledes bestemte Forudsætning er afgjørende for Opfattelsen af den
geologiske Hypothese. Sammenfatte vi de Formationer, Systemer og Grupper
--- Graavakkegruppen, Kulgruppen, Zechsteensgruppen, Trias-, Jura-, Kridt-
og Molassegruppen --- som Geologien for Tiden har indordnet i store
Mellemperioder, under et almindeligt Overblik, da kan man vanskelig værge
sig imod den Grundforestilling, at her under den hele physiske Udvikling
gaaer en drivende Strøm af Midler, Motiver og Maal; men seer man skarpt paa
de materielle Aarsager, saa er her jo ingen. Hvorfor blev den glødende Jord
indhyllet som i en Steenkappe med mange Folder? Fordi Lagene i dens
Overflade afkjøledes ved Udstraaling af Varme, og Eruptivmasser
gjennembrøde de slentede Lag. Hvorfor har den faaet et Hav? Fordi der i det
første, omgivende Gasrum var saa store Qvantiteter af Ilt og Brint, der
indgik i chemisk Forening. Hvorfor er der i Jordskorpen en saa overordenlig
Udbredelse af Silicater, Carbonater, Alkalier o.s.v.? Fordi de tilsvarende
Grundstoffer alle forud have været tilstede i den roterende Dampkugle. Det
er den blot physiske Begrundelse, som idelig gjentager sig. Men de physiske
Realgrunde ere kun halve Grunde, hvad øieblikkelig sees, naar der spørges
om de første væsenlige Forudsætninger. Hvorfor har Jordmassen oprindelig
indeholdt de mange forskjelligartede Substanser? Her staaer Begrundelsen
stille.

Den væsenlige Sammenhæng mellem den kosmologiske og den geologiske
Hypothese viser sig her i det klareste Lys. Er Dannelsen af Planetsystemet
en naturlig Selvdannelse, fordi den er fremgaaet af et i Realbegrebet
bestemt Naturanlæg, da maa det Samme gjælde om Jordlegemets Dannelse.
Jordlegemets Anlæg maa da siges at være indeholdt i Solsystemets Anlæg. Al
Naturudvikling er, væsenlig seet, en Anlægsudvikling, og enhver
Anlægsudvikling en Udvikling af Formaal, en teleologisk Udvikling. For at
bestemme, hvad og hvor meget der ligger i Anlæget, maatte vi altsaa forud
kjende de relative Formaal; for at bestemme de relative Formaal, maatte vi
forud kjende det sidste, afsluttende Formaal, Endemaalet; for at bestemme
Endemaalet, maatte vi kunne overskue Gangen i det Hele, men det er umuligt.
Formaalserkjendelsen bliver kun en menneskelig Tydning, en Gjætning, om man
vil, som, hvad Indholdets Realitet angaaer, ikke kan gjøre Fordring paa
objectiv-videnskabelig Gyldighed. Det Teleologiske beriger os ikke med
Kundskaber, men det betinger vor Indsigt og er, for Grundkategoriernes
Skyld, af afgjørende Betydning i Videnskaben.

Paa Kundskabens og den reelle Erkjendelses Vegne maae vi med Forskeren
holde os til Kjendsgjerninger og lærvillig følge hans Spor, idet han,
arbeidende med Enkeltheder, trænger dybere og dybere ind i Tingenes
virkelige Sammenhæng. Men naar Begrundelsens allervæsenligste Kategorier:
Anlæg og Motiv, Formaal og Middel, enten oversees eller endog ringeagtes,
vil det være Forskeren umuligt at gjøre videnskabelig Adskillelse mellem
Bisag og Hovedsag, mellem Væsenligt og Uvæsenligt; ved de blot physiske
Realgrunde bliver Alt, for Sammenhængens Skyld, lige væsenligt og just
derfor lige uvæsenligt. I samme Nu som Forskeren erklærer Nogetsomhelst,
det være nu Underlag af Steenmasser, Dannelse af Hav, Opskydning af Bjerge,
Udvidelse af Fastland, Bundfældning af Kalk, Rensning af Atmosphære, hvad
det nu er, for vigtigt og betydningsfuldt, underskyder han et Formaal og
indlægger et Motiv.

Heraf forstaaes da endelig den indre organiske Forbindelse mellem den
geologiske og den biologiske Hypothese, som i sit hele Omfang opfattet og
udtalt af Darwin, sætter Forskningen i Bevægelse, for at komme til sin
fulde Ret. Ligesom Jordklodens Anlæg er indeholdt i Solsystemets Anlæg,
saaledes er det organiske Livs Anlæg væsenlig indeholdt i Jordklodens. Den
organiske Dannelse er Naturdannelse, fordi den, teleologisk forstaaet, er
væsenlig Selvdannelse, en Selvdannelse, der har holdt Skridt med
Jordklodens Udvikling.


\chapter{Teleologisk-empiriske Videnskaber}

\section{Teleologi og Biologi}
\label{sec:17}
% pp. 172--183
% Key section for the Darwinism debate: teleological vs. purely
% mechanical interpretation of Darwin. Connects to ``Et Synspunkt
% for Darwinismen'' (1873b).

\noindent Under Bestræbelsen for at paavise en lovbestemt continuerlig
Overgang fra den uorganiske Natur til den organiske, er den physisk exacte
Forskning nærved at udslette Grændsen mellem det Livløse og det Levende.
Ganske naturligt. Livets væsenligste Særkjende er at være teleologisk; men
for det Teleologiske har den physisk-exacte Forskning ingen Brug; en
teleologisk Forklaring kan ikke tilfredsstille den Erkjendelse, der kun
spørger om iagttagelige Betingelser og materielt virkende Aarsager. En
naturlig Følge heraf er, at Forskeren selv, istedenfor at klare sig det
Videnskabelige, der ligger i Teleologiens Begreb, helst holder sig til
mythiske Forestillinger, mod hvilke der i Videnskabens Navn maa
polemiseres. Et overnaturligt Væsen, der har indblæst og endnu maaskee af
og til indblæser Materien Liv, hjælper ikke Videnskaben frem; et Livsstof
for sig eller en enkelt og særegen Livskraft er ingensteds at opdage: det
ligger da nær at fremhæve Ligheden imellem de uorganiske og organiske
Dannelser saa kraftig og alsidig, at Overgangen bliver umærkelig, og
Naturlovene de samme.

Forskjellen, hedder det, mellem det Organiske og det Uorganiske, hvad
Aggregationstilstanden angaaer, er kun den, at de tre Tilstandsformer, den
faste, den flydende, den luft- eller gasformige, der for det Uorganiske
ere saa betegnende, i de organiske Dannelser combineres til en
fast-flydende Form, der meddeler dem en egen »Imbibitionsevne«. Paa Grund
af denne Imbibitionsevne kunne Organismerne adskilles fra Krystallerne, med
hvilke de for Resten, hvad navnlig de laveste Skeletformer udvise, have en
paafaldende Lighed. --- Ved Hjælp af nogle hinkende Analogier og især ved
Henvisning til Kulstoffets mærkværdige Evne til at forbinde sig med de
andre Elementer i de forskjelligste Forhold kan Forskeren tilsidst bilde
sig ind, at det nu endelig er blevet beviist, hvorledes det Levende med
»absolut Nødvendighed« frembringes af det Livløse. En saadan Indbildning
kan være ganske uskyldig, naar blot Forskeren ikke bliver dogmatisk; men
skulde han faae et Anfald af Dogmatisme og reentud paastaae, at Livets
Oprindelse hermed er tilstrækkelig forklaret, og at navnlig ethvert Forsøg
i teleologisk Retning maa betragtes som Tegn paa forældede
Hjernesvingninger: da maa der i Videnskabens Navn protesteres.

Undersøgelserne over det Teleologiske blive sjældent gjennemførte med
fuldstændig Consequens, idet de alt for ofte belæmres med dobbelsidige
Misforstaaelser. Teleologiens Forsvarere gaae gjerne ud fra den
Forudsætning, at det er Forskeren som Forsker, der trænger til et Cursus i
Læren om Formaal og Midler, at en grundig, speculativ-dogmatisk
Øiemedslære i Forbindelse med nogle Skabelsesmirakler, destillerede i
sindrig Forklaring af en paalidelig Tradition, bedst vilde hjælpe ham frem
paa Forskningens Vei og vise ham Broen over den dybe Kløft fra det Livløse
til det Levende. Men Forskeren forkaster den tilbudte Anviisning, og det af
den gode Grund, at den slet ikke passer til hans videnskabelige Opgave.
Misforstaaelsen vender sig, idet Forskeren, som vitterlig ikke har Brug for
det Teleologiske, føler sig forpligtet til at jage Teleologien ud af
Videnskaben ved physisk exacte Beviser. Og hvad beviser han saa? Intet, som
i mindste Maade rører Teleologiens Begreb. Han gaaer ud fra den
Forudsætning, at hvad der beviislig fremgaaer som Resultat af physiske
Aarsager og skeer med betinget Nødvendighed, umulig kan være teleologisk;
men, hvor stor denne Misforstaaelse er, kan sees af Realbegrebernes
Analyse. Han gjør gjældende, at Indbildningen om det Teleologiske hæmmer
Forskningen, idet man, navnlig i Biologien og Physiologien, finder det
bekvemmere at betjene sig af Talemaader om »Livskraften«, »Vitalkraften«,
»Livsprincipet«, »Ideen« o.\ dsl., istedenfor at arbeide med Vævenes
Anatomi og Kulforbindelsernes Chemi. Men at Teleologien paa ingen Maade vil
fritage Forskeren fra Arbeide, følger ligefrem af Sagens Natur.

At Teleologi og Biologi ere sammenhørende Begreber, giver sig uvilkaarlig
tilkjende i det betegnende Slagord: Kamp for Livet; det er en teleologisk
Parole. Trods de Analogier, der ellers kan paavises mellem uorganiske og
organiske Legemer, ere Forskjellighederne dog saa iøinefaldende, at de
umulig kunne undgaae den sunde Sands; de give sig umiddelbart tilkjende ved
betegnende Ord og Vendinger. At uorganiske Legemer, især de faste, gjøre
Modstand mod hvad der vil trænge ind i dem og forandre Delenes Sammenhæng,
er unægteligt, og man kan for saa vidt betragte enhver Ting som et Individ,
der i Forhold til sine Kræfter hævder sin Selvstændighed; men Ingen vil dog
aflede det uorganiske Individs Modstand mod Opløsning af en indre Trang til
Selvopholdelse. Ingen vil paastaae, at Isen af indre Trang til at holde sig
fast værger sig mod Varmen, eller at Vandet føler en stille Længsel efter
at krystallisere. De uorganiske Individer, der punktlig opfylde Loven for
Action og Reaction, ere aldeles ligegyldige mod de Tilstande, hvori de
befinde sig, og de Forandringer, de undergaae; det er Klippen lige saa
ligegyldigt, om den bevarer sin Reisning eller splintres af Lynet, som det
er Jordskorpen ligegyldigt, om den revner i Solen eller overskylles af
Havet: denne fuldstændige Ligegyldighed er netop Særkjendet for det
Livløse. Men at de organiske Individers Selvvedligeholdelse maa betragtes
fra et ganske andet Synspunkt, indskærper Videnskaben selv, idet den sætter
Kampen for Livet som uafviseligt Vilkaar for alt Levende. Selvopholdelsen
bliver jo herved opfattet som udtrykkeligt Formaal, d.v.s.\ som den Opgave,
Naturen har sat hvert levende Individ.

Denne Opgave paatrykker alle organiske Functioner et særeget Præg, saa at
de trods physiske Ligheder væsenlig adskille sig fra tilsvarende
uorganiske. En Krystal kan voxe saa vel som en Plante og et Dyr; men at det
med den organiske Væxt har et ganske andet Sammenhæng end med den
uorganiske, sees paa Ernæringen. Planten voxer paa Grund af Ernæring, men
Krystallen ved Opstablen af Dele, dens Væxt er uden Ernæring. Ernæringen,
der i sig forener Stofskifte med Assimilation og Vedligeholdelse af en
typisk Form, betegner et teleologisk Kredsløb; Ernæringen foregaaer ikke
blot, saa at, men udtrykkelig, for at det organiske Individ kan opholde
Livet.

Spørge vi om de materielle Aarsager og Betingelser, hvoraf den ernærende
Virksomhed afhænger, saa har Forskeren Ret, naar han Punkt for Punkt
henviser til Delenes, Molekulernes physiske Beskaffenhed, Leiring og
Omsætning; men han har ikke Ret, naar han ved Analogier hentede fra det
Uorganiske, mener at kunne begrunde den sig fortsættende Virksomhed paa en
fortsat Ophævelse af mechanisk og chemisk Ligevægt og dermed bortforklare
det Teleologiske. Exemplerne med »Diosmosen« bevise netop det Modsatte, af
hvad de skulde. De bevise nemlig, naar man holder sig strengt til det
Uorganiske, at Udligningen standser efter kortere eller længere Tids
Forløb, saa at Bevægelsen standser; men det, som skulde bevises, er jo
netop, hvorledes det gaaer til, at Virksomheden stadig vedligeholdes. Hvad
der anføres om Imbibitionsevnen, større Spændighed i Molekulernes Forhold,
Sollysets Andeel i Kulsyrens Adskillelse o.s.v., hører vistnok med til de
physiske Betingelser; men hvad Ernæringsphænomenet selv angaaer, ere de
herpaa støttede Forklaringsgrunde saa svage, at Forskeren, hvis han havde
Udsigt til at opnaae andre og bedre, uden Tvivl selv vilde være den Første,
der blottede deres Svaghed. Her savnes Noget, men hvad er det, og hvorfra
kommer det?

Hvad er det? Livet selv! For, om muligt, at gjøre det fuldstændig
indlysende, hvad dette vil sige, maae vi fremfor Alt see at sikkre os en
klar og utvetydig Sprogbrug. Der er Virksomhed saa vel i det Uorganiske som
i det Organiske; dersom der nu er nogen principiel Forskjel mellem Livløst
og Levende, da maa der være en tilsvarende Forskjel imellem uorganisk
Virksomhed og Livsvirksomhed, og dersom denne Forskjel skal kunne blive os
ganske tydelig, maae vi søge at faae den betegnet saa skarpt, som muligt.
Vi begynde desaarsag med at gjøre en Adskillelse mellem Stofvirksomhed og
Selvvirksomhed, idet vi ved Stofvirksomhed forstaae enhver Kraftyttring,
der udgaaer fra Stoffet som Stof, medens vi ved Selvvirksomhed tænke os
Functioner, der ikke have deres Udgangspunkt i et særeget Stof, men i en
til Stofbetingelser svarende ideel Causalitet.

Hvilke ere nu de for al Stofvirksomhed afgjørende Kriterier? Forestiller
man sig to Massedele $m_1$ og $m_2$ virkende paa hinanden, og lader den
Lov, hvorefter Virksomheden foregaaer, være betegnet ved $f$, Afstanden ved
$u$, Virksomheden ved $V$, saa kan Ligningen $V = m_1 m_2\, f(u)$ opstilles
som Symbol og Grundformel for al mulig Stofvirksomhed. Ved Anvendelse af de
saakaldte Potentialfunctioner og »Potentialets« Forhold til »levende
Kraft«, kan det exact bevises, at $f(u)$ kan være hvilken som helst, $u$
variere mellem Grændserne $0$ og $\infty$, Massedelenes Antal være saa
stort, og deres Combinationer saa forskjellige, man vil, idet man hver Gang
mærker sig Forskjellen mellem bevægelige og ubevægelige Massedele, selv de
Dele, Potentialet angaaer, og endda vil Stofvirksomhedens Symbol ---
Virksomheden $f(u)$ være frastødende eller tiltrækkende, mechanisk,
magnetisk, elektrisk, chemisk, kort sagt, af hvad Art man vil --- holde sig
uforandret. For at komme til Erkjendelse af, hvad der egenlig skal forstaaes
ved Stofvirksomhed, behøve vi altsaa kun at underkaste Grundformlen en
logisk Analyse.

Af Formlen $V = m_1 m_2\, f(u)$, naar den med Hensyn til Stillinger,
Projectioner o.s.v.\ udvikles, fremlyser da: $\alpha$) at Materien under
alle Betingelser virker efter Inertiens Lov: $m_1$ virker paa $m_2$, ikke
ifølge nogen Selvbestemmelse eller indre Tilskyndelse, men alene fordi
$m_2$ virker paa $m_1$; det er selvløs Virksomhed; $\beta$) hvor vidt der i
et System af virksomme Masser eller Molekuler er Hvile eller Bevægelse,
afgjøres ikke ved nogen selvstyrende Eenhed, enten i det Hele eller i
nogensomhelst enkelt Deel, men alene ved de samvirkende Kræfters Resultant;
$\gamma$) Systemets Overgang fra Bevægelse til Hvile, eller omvendt, har
ikke sin Grund i Systemet selv, men afhænger af udvortes Aarsager, udvortes
Betingelser; $\delta$) de virksomme Led ere lutter Stofsubjecter, og det i
Virksomheden Eiendommelige henføres til en eller anden for det paagjældende
Stof charakteristisk Egenskab, i mechanisk Physik til Tyngden, i Chemi til
Affiniteten. Er dette nu klart, saa behøve vi ikke at fortsætte Analysen
videre, for at overbevise os om, at Stofvirksomheden, hvormange Betingelser
og Combinationer man end vil underskyde, umulig kan omdannes til
Livsvirksomhed. Her mangler aabenbart Noget, og det er dette Noget, vi nu
skulle charakterisere som det Selvvirksomme.

Hvad er Selvvirksomhed, og hvorpaa kjende vi den? Det ligger nær at vise
hen til en umiddelbar Bevidsthed om Villen og Villie; thi uden Hensyn til,
hvad Realitet, der kan tillægges vor Villen og Selvbestemmen, saa er det
uimodsigeligt, at vi heri have en Virksomhedsform, som er den materielle
diametralt modsat. Hvad vi kalde vort Selv er intet for sig tilværende
Stof, og den fra Selvet udgaaende Virksomhed ingen Stofvirksomhed. Det
virksomme Selv kan ikke, i Lighed med en Massedeel, existere som Led i
Modsætning til sin Virksomhed, thi det har ingen umiddelbar Bestaaen; dets
Bestaaen er en ideel Virken, og denne Virken en uafbrudt Reflecteren. Ud
fra sit ideelle Midtpunkt gaaer Virksomheden over paa Materien og fra
Materien tilbage til Udgangspunktet; Selvvirksomheden er en Kredsbevægelse,
hvorunder Modsætningen af Midtpunkt og Peripheri stadig afspeiler Eenheden.

Lad det Selvbestemmende, $s$, være omgivet af materielle Punkter, $a$, $b$,
$c$, saa vil Forholdet ikke være det, at $s$ kun kan virke, for saa vidt
det umiddelbart reagerer mod Paavirkning. Vexelvirkningen foregaaer
tvertimod saaledes, at $s$, afficeret af $a$, kan derved bevæges til at
virke paa $b$, eller, hvis det paavirker $a$, da virke snart tiltrækkende
snart frastødende, eftersom Tilstanden er, hvori det befinder sig.
Selvvirksomheden beroer paa, at $s$ virker efter sit Befindende, og
desaarsag væsenlig har Virksomhedsmotivet i sig selv. Et System af
materielle Dele giver kun en Sammenhængseenhed; først naar det
Selvbestemmende kommer med, viser der sig en styrende Eenhed, en
Selveenhed.

At en styrende Eenhed, en formgivende Selveenhed spores overalt i alle
levende Organismer og organiske Dannelser, i Moneren og Cytoden lige saa
vel som i Menneskets selvbevidste Sjæl, er en saa iøinefaldende
Kjendsgjerning, at ingen Physiolog vil forsøge paa at faae den
bortforklaret, med mindre det skeer i Fortvivlelse over, ikke at kunne
forklare den.

Hvorfra kommer den formende, styrende Eenhed, den ideelle Causalitet, den
virkende Selveenhed? Fra Realbegrebernes System, fra Livets Anlæg. Dette
Svar maa naturligvis forkastes af dem, der paa Forhaand ere hildede i den
Fixidee, at virksomme Realbegreber med væsenlige Anlæg ikke høre med i
Naturens Orden. Men selv om man indrømmer Gyldigheden af et almindeligt
Livsanlæg, er Sagen i sig dog dunkel og gaadefuld. Hvorledes skal man her,
uden at indsmugle uvidenskabelige Forudsætninger eller blive staaende ved
den almene, og af den Grund intetsigende, Henvisning til Forholdet mellem
Realbegreb og Existens, kunne danne sig nogen bestemt Forestilling om det
Selvvirksommes Forhold til Stoffet?

Her er flere Hypotheser at vælge imellem. Man kan f.\ Ex.\ gaae ud fra en
med Lyshypothesen analog Livshypothese, og antage, at ligesom Lyset har sin
reale Mulighed i et fra de veielige Masser forskjelligt Medium, saaledes
har Livet sit Anlæg i et tilsvarende overalt udbredt Medium. Naar Lyset
tændes, skeer det under Betingelse af, at Lysmediet sættes i Virksomhed. At
den ideelle Selvvirksomhed maa være betinget af en tilsvarende
Stofvirksomhed, ville vistnok Alle, der kjende Naturlovene, uden videre
indrømme. Men det Selvvirksomme kan jo meget vel være betinget af Stoffet,
uden derfor at være Stofsubject. At Selvvirksomheden indtræder, beroer da
paa, at det Mechaniske, Chemisk-Physiske og Chemiske i Stofvirksomheden
sætter Livsmediet i Svingning: under Livsmediets Svingninger danne sig
Udgangspunkter for Selvvirksomheden. I en Molekule $A$ være Masseatomerne
$a$, $b$, $c$ o.s.v., Ætheratomerne $\alpha$, $\beta$, $\gamma$ \ldots;
Masseatomerne ville tiltrække hverandre, men Ætheratomerne frastøde dem og
fjerne dem fra hverandre. Hvis der nu ingen Selvvirksomhed er i Molekulen,
saa vil Spændingen kunne udlignes physisk, d.v.s.\ Molekulen er livløs. Har
derimod Stofvirksomheden sat Livsmediet i Bevægelse, saa vil der have
dannet sig en Mængde Selvhedspunkter $s$, $s$, $s$ mellem Atomerne, disse
Selvhedspunkter kunne hvert for sig være saa svage, at deres Virksomhed
nærmer sig til en Differentialvirksomhed: de ville dog bevirke, at
Molekulen bliver levende. Bestaaer nu en levende Celle af levende
Molekuler, saa fremkommer den styrende Eenhed derved, at $s$, $s$, $s$ ikke
blot ere i Vexelvirkning med de materielle Dele, men under denne
Vexelvirkning tillige i et saadant Forhold til hverandre indbyrdes, at de
danne en Resultant $S$, som deri er forskjellig fra enhver Resultant af
materielle Kræfter, at den som Selvhedsresultant behersker sine Composanter
og nedsætter dem til Tjenere for sine særegne Fordringer. Bestaaer
Organismen af flere Celler, saa ville $S$, $S$ gaae ind under en høiere
Resultant, som i teleologisk Selvvirksomhed vil lede Arbeidets Deling
o.s.frmdl.

Vil man foretrække en anden Hypothese og antage, at det er Atomerne selv,
der, ifølge en Duplicitet i Materiens Væsen, snart holde sig paa et
mechanisk og chemisk Standpunkt som umiddelbare, Inertiens Lov undergivne
Stofsubjecter, snart igjen, især i Kulstofforbindelser, hæve sig til at
afgive Selvhedsresultanter, da kan dette kun skee paa den udtrykkelige
Betingelse, at der gjøres videnskabeligt Rede for, hvor Inertien ophører,
og hvor Selvvirksomheden begynder, samt hvorledes det gaaer til, at de
samme Stofdele kunne paa een Gang bevare deres Inerti og tillige fungere
som Momenter i den organiserende Selvvirksomhed.

Vil man endelig opgive alle Hypotheser, beraabe sig paa Videns
Indskrænkning og erklære det biologiske Problem for slet og ret uløseligt,
da er dette unegtelig en Udvei; og Forskningen kan uforstyrret gaae sin
rolige Gang. Men afskaffe Selvvirksomhedens Factor og bortforklare
Problemet ved Hjælp af hule Analogier og physiske Tvetydigheder gaaer paa
ingen Maade an. Grundformlen $m_1 m_2\, f(u)$ vil vedblivende vidne om
Stofvirksomhedens eiendommelig begrændsede Sphære, og derved om qvalitativ
Forskjel mellem Livløst og Levende.

Livsvirksomhed er hverken Stofvirksomhed for sig ei heller Selvvirksomhed
for sig, men en gjennemgaaende Eenhed af begge. Enhver organisk Function er
paa een Gang en Stoffunction og en Selvhedsfunction. Fra dette Synspunkt
gjælder det, at »Hjernen tænker, ligesom Maven fordøier«; dog maa det, ikke
blot i Hjernens og Mavens, men fremfor Alt i Videnskabens Navn paa det
Eftertrykkeligste fremhæves, at Forholdet mellem Stofmomentet og
Selvhedsmomentet i de organiske Dobbelfunctioner er indtil de mindste
Enkeltheder articuleret.

Som forskende Videnskab kan Physiologien ikke direct gjøre Brug af det
Teleologiske; den er og bliver en empirisk Videnskab. Men fornægte eller
oversee Teleologiens Betydning kan den heller ikke, thi Physiologi er i
væsentligt Slægtskab med Biologi. Physiologien kan ikke lægge an paa at
gjøre Biologien umulig; men Biologi og Teleologi ere nødvendig
sammenhørende Begreber.

\section{Botanik og Zoologi}
\label{sec:18}
% pp. 184--195

% [text to be added]

\section{Den biologiske Afstamningslære}
\label{sec:19}
% pp. 195--205

% [text to be added]

\section{Anthropologi og Psychologi}
\label{sec:20}
% pp. 205--215

% [text to be added]


%%% TREDIE DEEL: AANDSVIDENSKAB %%%

\part{Aandsvidenskab}

\chapter{Menneske-Aanden}

\section{Aandens Naturgrund: Nødvendighed og Frihed}
\label{sec:21}
% pp. 216--229

% [text to be added]

\section{Menneske-Aandens Selvudvikling: Historie}
\label{sec:22}
% pp. 230--243

% [text to be added]


\chapter{Menneske-Aandens Selvbefrielse}

\section{Selvbefrielse i Phantasiens Medium: Æsthetik}
\label{sec:23}
% pp. 243--258

% [text to be added]

\section{Selvbefrielse i Villiens Medium: Ethik}
\label{sec:24}
% pp. 258--269

% [text to be added]


\chapter{Den absolute Aand}

\section{Videnskaben og den absolute Aand}
\label{sec:25}
% pp. 270--277

% [text to be added]

\section{Religionsphilosophi}
\label{sec:26}
% pp. 277--281
% The final §: the *uensartede størrelser* thesis restated in the
% context of the completed Videnskabslære. The religious life
% ``vedblive at danne sig af sin egen Aand, efter sine egne
% Forbilleder, og gaae sin Gang uafhængigt af Videnskab og
% videnskabelige Indsigelser'' (Forord).

% [text to be added]

\end{document}
